% LaTeX fragment. Validated with XeLaTeX, fontspec (DejaVu Serif), amsmath, amssymb and hyperref.
% An optional wrapper and the complete successful compilation log are in scripts/.
% All the arguments in this file are this run's, not a transcription.
\section*{Conventions and common proofs}
The labels below are shared lemmas. Each statement section specifies its application, any literal counterexample, and its second check. A common proof is not a claim that Grothendieck supplied that proof. ``Closed'' without a topological subscript means the manuscript's finite closed-interval calculus. A singleton is allowed as a pseudo-troncon. The empty set is allowed as an empty union, never as a nonempty component.

\paragraph{G: the actual ordered-troncon axioms.}\label{lem:G}
The order is a partial order. The set $L$ has at least two elements. Every two elements have a common upper bound and a common lower bound. If $a<b$, some $c$ satisfies $a<c<b$. If $a<x,y$, some $z$ satisfies $a<z\leq x,y$; dually, if $x,y<a$, some $z$ satisfies $x,y\leq z<a$. The cone condition is imposed when the cone is nonempty. ``Filtered'' includes nonemptiness; the separate $\operatorname{card}L\ne1$ then excludes the remaining trivial case. Neither total order nor Dedekind completeness is assumed.
A maximal element is greatest: a common upper bound of it and any $x$ must equal it. The dual holds for a minimal element. Thus there are at most two absolute extremal points. Nonempty strict cones contain at least two elements, by interpolation. They have the global bounds required by G: in $L_{>a}$, lower bounds follow from the local axiom at $a$, and upper bounds from global upper bounds; all remain $>a$. Its local axioms and density are inherited. The dual proves the assertion for $L_{<a}$. For an interval with endpoint $b$, upper bounds below $b$ use the local axiom at $b$; if $b$ belongs to the interval it itself may be used. Repeat dually at the lower endpoint. Consequently every nondegenerate interval of any endpoint type, and every nonempty ray, has G. Reversing the order exchanges every axiom with its dual.
\emph{Independent check.} For any finite set of points strictly below a common endpoint, repeated use of the binary local axiom gives a common bound still strictly below it. Apply this finite-cone observation to the two points inside each interval. This directly checks the only issue that would fail for an arbitrary convex subset of an arbitrary directed poset.

\paragraph{C: finite cuts and their partitions.}\label{lem:cuts}
Given a symmetric set $S_0\subseteq\mathcal P_2(L)$, call a finite set tame if all of its two-element subsets belong to $S_0$. Put
$C(T)=\{x\notin T:T\cup\{x\}\text{ is tame}\}$.
If $T\subseteq S$, then $C(S)\subseteq C(T)$ by restriction of pairwise compatibility. If $S\cup T$ is tame, the converse implication follows: $t\in T\setminus S$ lies in $C(S)$ and not in $C(T)$. This extra compatibility hypothesis is necessary in general.
For an ordered G, $S_0$ is distinct comparability. If $T=\{t_1<\cdots<t_n\}$, every $x\in C(T)$ occupies exactly one of
\[
 L_{<t_1},\quad (t_1,t_2),\ldots,(t_{n-1},t_n),\quad L_{>t_n}.
\]
Indeed $x$ is comparable with every $t_i$ and unequal to them. Its lower indices form an initial segment; transitivity excludes any other placement. Conversely every point of a displayed set belongs to $C(T)$. Internal sets are nonempty by density; a left or right tail is empty exactly when its endpoint is the corresponding ambient extremum. For $T=\varnothing$ use the single class $L$. These are definitions of combinatorial cut components, not assertions about connected components of a topology. Assign each nonempty gap the adjacent cut points as its frontier. Adding cut points refines the gaps. An old cut point on the boundary of a smaller gap is a boundary point of the containing old gap: this proves $\dot A\cap T\subseteq\dot B$ for $T\subseteq S$. There are at most $n+1$ classes, exactly $n+1$ without ambient endpoints.
\emph{Independent check.} Record for each $x$ the integer $\#\{i:t_i<x\}$. Its fibers are precisely the displayed sets. Deleting cut points merges consecutive integer fibers and gives the same refinement and boundary formula.

\paragraph{R: reconstruction from betweenness and cuts.}\label{lem:reconstruction}
Let $C(x,z)$ mean reflexive comparability and let
\[
 B(x,y,z)\iff x<y<z\text{ or }z<y<x,\qquad
 N(x,y,z)\iff B(x,y,z)\text{ or }((x=y\text{ or }y=z)\land C(x,z)).
\]
Then $N$ means $x\leq y\leq z$ or the reverse. In G, for distinct $x,z$, comparability is equivalent to $\exists y\,B(x,y,z)$, by density. Fix $a<b$. Define
\[
 D(x)\iff x=a\text{ or }B(x,a,b),\qquad
 U(y)\iff y=b\text{ or }B(a,b,y),
\]
\[
 Q_{a,b}(u,v)\iff\exists x,y\,[D(x)\land U(y)\land N(x,u,y)\land N(x,v,y)\land N(x,u,v)].                 \tag{R1}
\]
If $u\leq v$, take a common lower bound $x$ of $a,u$ and a common upper bound $y$ of $b,v$. They satisfy the formula. Conversely $D(x),U(y)$ imply $x\leq a<b\leq y$. Hence $N(x,u,y),N(x,v,y)$ put $u,v$ in $[x,y]$. In $N(x,u,v)$ the alternative $v\leq u\leq x$ forces $u=v=x$; otherwise $x\leq u\leq v$. Thus in all cases $u\leq v$. This proves R1. The alternative extra condition $N(u,v,y)$ works by the dual argument. In particular the four-element lemma on scan 14 needs no totality assumption and allows equalities.
For three distinct pairwise comparable points $a,b,c$, let $A$ be the $R_{\{a,b\}}$-class of $c$ and $A'$ the $R_{\{b,c\}}$-class of $a$. If $a<b<c$, these are $L_{>b}$ and $L_{<b}$ and are disjoint. If $b$ is not the middle point, their intersection is the nonempty interval between the two appropriate neighboring points; for example, if $b<a<c$, it is $(b,a)$. Reversal handles the other three orders. Therefore
\[
 B(a,b,c)\iff C(a,b)C(b,c)C(a,c),\quad |\{a,b,c\}|=3,\quad A\cap A'=\varnothing.                 \tag{R2}
\]
Thus the two-point cut data determine the order up to simultaneous reversal, and determine every finite canonical cut partition and frontier by C.
\emph{Independent re-proof.} Suppose two orders have the same $C,N$ and orient $a,b$ alike. For $u<v$ in the first choose $p\leq a,u$ and $q\geq b,v$. The chain $p\leq a<b\leq q$ orients $p,q$ the same way in the second order. The relations $N(p,u,q)$ and $N(p,v,q)$ place $u,v$ between them. If $p<u<v$, $B(p,u,v)$ orients $u<v$; if $p=u$, use $B(p,v,q)$, unless $v=q$, when the already oriented endpoints suffice. The case $v=q$ with $p<u$ uses $B(p,u,q)$. Interchanging the orders gives equality. If the anchor is reversed, apply this to the opposite order. This proof uses only global bidirectedness once $C,N$ are given. The finite enumeration in \texttt{scripts/check\_finite.out} separately checks R1 and the four-point lemma on all naturally labelled bounded posets of sizes 2 through 7. It does not claim those finite examples satisfy density.

\paragraph{A: a fully explicit recognition scheme.}\label{lem:axioms}
In the language with one ternary predicate $B$, abbreviate
$C(u,v)$ by $u=v\vee\exists w B(u,w,v)$ and $N$ as in R. Use R1 as a formula, not as an extra relation symbol. Require that there exist $a\ne b$ such that $C(a,b)$ and the formula $Q=Q_{a,b}$ satisfies the following finite list:
\[
\begin{gathered}
\forall x\,Q(x,x),\quad\forall x,y\,[Q(x,y)\land Q(y,x)\Rightarrow x=y],\\
\forall x,y,z\,[Q(x,y)\land Q(y,z)\Rightarrow Q(x,z)],\\
\forall x,y\,\exists p,q\,[Q(p,x)\land Q(p,y)\land Q(x,q)\land Q(y,q)],\\
\forall x\ne y\,[Q(x,y)\Rightarrow\exists z\,(Q(x,z)\land Q(z,y)\land z\ne x\land z\ne y)],\\
\forall a,x,y\,[(a<_Qx\land a<_Qy)\Rightarrow\exists z\,(a<_Qz\land Q(z,x)\land Q(z,y))],
\end{gathered}
\]
the dual of the last formula, $|L|\geq2$, and
\[
\forall x,y,z\,[B(x,y,z)\ \Longleftrightarrow\ (x<_Qy<_Qz)\vee(z<_Qy<_Qx)].                 \tag{A1}
\]
All symbols $<_Q$ abbreviate $Q$ plus inequality. Every formula quantifies only over elements. Necessity is R1 and G; sufficiency follows because the definable $Q$ itself is the required order. This is a finite first-order axiomatization, not an independent or economical list.
For cut data use a symmetric irreflexive predicate $S(a,b)$ and a four-place relation $R(a,b;x,y)$. Its domain is $D_{a,b}=\{x:S(a,b)\land S(a,x)\land S(b,x)\}$. Require $R$ to be symmetric in $a,b$, an equivalence relation exactly on this domain, and false outside it. The factor $S(a,b)$ is essential: the original cuts are indexed only by tame distinct pairs; our total four-place encoding extends the relation by false at every other pair. Define $B$ by pairwise $S$, distinctness and
$\neg\exists z\,[R(a,b;c,z)\land R(b,c;a,z)]$.
Impose A1 and G through the anchored formula, require $S$ to equal distinct $Q$-comparability, and require, whenever $a<_Qb$, that $R(a,b;x,y)$ hold exactly when $x,y$ lie in the same one of $L_{<_Qa},(a,b)_Q,L_{>_Qb}$. This last condition is an explicit disjunction of three pairs of inequalities. R2 proves necessity; the imposed displayed formulas prove sufficiency. All finite cuts and their boundary data are then the formulas in C. No unknown higher-cut relations are silently assumed.
\emph{Independent check.} Given a model of these sentences, actually set $u\leq v$ iff R1. The first seven clauses verify every order axiom and every G axiom. The last clauses identify, rather than merely approximate, the supplied $B,S,R$. Conversely substitution of any G-order satisfies each clause. This two-way construction is independent of any literature classification.

\paragraph{D: the doubled rational line and the basic obstruction.}\label{lem:double}
Set $D=\mathbb Q\times\{0,1\}$ and put $(q,i)<(r,j)$ iff $q<r$, with equality added for $\leq$. Reflexivity, antisymmetry and transitivity follow from the strict order of heights. Choose heights below the minimum and above the maximum for global bounds. Interpolate rational heights for density. For $a<x,y$, choose a height strictly between $h(a)$ and $\min(h(x),h(y))$; either point at that height is a common lower bound still above $a$. The dual proves the local upper-bound axiom. Thus D satisfies every G axiom and has no endpoints.
Let $a_0=(0,0)$ and $a_1=(0,1)$. They are distinct incomparable twins with identical strict upper and lower rays and identical $C(\{a_i\})$. This refutes injectivity and the converse inclusion assertion of scan 3 even in the intended canonical class.
Put $V=D_{<a_0}=D_{<a_1}$. Both $H_i=D_{\leq a_i}$ are closed tame supersets. A least closed tame superset $H$ would lie in their intersection $V$. No such $H$ can contain all of $V$: every closed interval component of $H\subset V$ has a finite upper endpoint of negative height (an unbounded upper ray would contain height zero), and finitely many such endpoints have a largest height $r<0$; choose $r<s<0$. Hence $V$ has no least closed tame hull. Also $V=H_0\setminus\{a_0\}=H_1\setminus\{a_1\}$, so its purported endpoint set is not intrinsic. An inclusion $V\subset H_1$ can be compatible or incompatible according to which of these presentations is chosen. The complements prescribed for $\{a_0\}$ and $\{a_1\}$ coincide, so they cannot define an involution on actual subsets.
The topology generated by strict intervals and rays on D is pulled back from the usual topology on $\mathbb Q$ under height. Its singleton closures are the two-point fibers. It is not $T_0$; a fiber is irreducible closed with two generic points, so it is not sober either.
In the join of this topology with the topology having principal closed intervals and rays as a closed subbasis, $V\subset H_1$ and $a_0\notin H_1$ show $a_0\notin\overline V$. For $a_0<b_0=(1,0)$, the closed interval $[a_1,b_0]$ contains $(a_0,b_0)$ but not $a_0$. Thus $\overline{(a_0,b_0)}\ne[a_0,b_0]$ as well. The open neighborhood $D\setminus[a_1,b_0]$ of $a_0$ misses the whole right-hand interval $(a_0,b_0)$; it contains no $(a_0,\varepsilon)$ with $a_0<\varepsilon\leq b_0$.
\emph{Independent check.} Every displayed set in this example is determined by its rational height except for the explicitly included endpoint twin. Intersecting the two closed rays deletes both zero-height points; the rational sequence of heights $-1/n$ then contradicts any finite list of negative endpoints. This separately checks both the absence of a least hull and the direction of all closure counterexamples.

\paragraph{E: the exact endpoint condition and its proof.}\label{lem:E}
The following is OUR added condition, absent from the definition on scan 9:
\[
(E)\qquad a=\sup\{x:x<a\},\qquad a=\inf\{x:a<x\}\quad(a\in L).
\]
At a least or greatest element use $\sup\varnothing=\min L$ or $\inf\varnothing=\max L$. Equivalently every upper bound of the strict predecessors of $a$ is at least $a$, and dually. Within G, E is equivalent to the existence of a least closed tame hull for every nonempty strict principal ray.
For necessity, let $H$ be the least closed tame hull of $L_{<a}$. It lies in $L_{\leq a}$. If $a\notin H$, each of its finitely many closed components has a greatest endpoint strictly below $a$. The local axiom gives a common bound $c<a$ for those endpoints, and density gives $c<z<a$, contradicting $z\in H$. Hence $H=L_{\leq a}$. If $b$ bounds all $x<a$, $L_{\leq b}$ is another closed tame superset, so $a\leq b$. Duality proves the other half of E. If a strict cone is empty, global directedness makes $a$ the appropriate extremum, and the convention is automatic.
For sufficiency, if $a<b$, the interval $(a,b)$ is cofinal in $L_{<b}$. Given $x<b$, use local filtering below $b$ on $x,a$ and then density to obtain $x\leq z<b$ and $a<z$. Its dual gives coinitiality above $a$. E therefore makes $a,b$ the actual infimum and supremum of each nonempty interval with these displayed endpoints, whether or not those endpoints are retained. A bounded ray has the corresponding actual endpoint as well. Consequently a nonempty interval has unique intrinsic endpoints, a unique closed hull, and a unique endpoint-removal presentation. Global least hulls for finite tame unions follow from F below. The condition E is inherited by every nondegenerate interval or ray: at an internal point, the strict predecessors within the interval are cofinal among ambient strict predecessors, by the same finite local-filtering argument; the dual handles successors. At an included relative extremum the empty-cone convention applies. Singletons have the trivial relative calculus.
There is also an independent topological equivalence. In G, every canonical neighborhood of a nonminimal $a$ meets $L_{<a}$: a basic neighborhood has finitely many lower and upper bounds around $a$; the finite local filtering lemma and density insert a point immediately below $a$ satisfying all of them. Dually every neighborhood meets the strict upper cone. If all principal closed rays are canonically closed, their containment of a strict ray forces the E inequalities by this closure observation. Conversely, under E, if $x\not\leq a$, some $p<x$ satisfies $p\not\leq a$; otherwise E would imply $x\leq a$. The open neighborhood $L_{>p}$ of $x$ misses $L_{\leq a}$. This proves that $L_{\leq a}$ is closed; the dual handles $L_{\geq a}$. It follows that the closed-subbasis topology is contained in the canonical topology, so their supremum is the canonical topology under E, and that canonical closures of nonempty open intervals and rays are exactly their endpoint closures. Endpoint accumulation is the finite-local-filtering argument just given. E also implies $T_1$, since a singleton is the intersection of its two principal closed rays. $T_1$ does not imply E, by P below.
\emph{Independent check.} The topological characterization gives a second route to every least-hull assertion: the canonical closure of a finite union is the union of the finitely many endpoint closures, hence is closed tame after merging touching closed pieces, and is contained in every closed tame superset. It also shows that E is necessary, not merely a convenient sufficient strengthening.

\paragraph{F: interval calculus, including the non-set-theoretic join.}\label{lem:flags}
Assume G and E unless a statement below is explicitly about closed pieces alone. Two pieces are compatible when their intrinsic finite endpoints form a chain (ambient extrema may be adjoined without changing this). For $T<T'$ meaning $t<t'$ for all $t\in T,t'\in T'$, write $T\ll T'$ when a strict separator $z$ satisfies $T<L_z<T'$; explicitly $t<z<t'$ for all such points. A finite union of successively strongly separated pieces is tame.
First, normal components are unique even without E if $\ll$ is defined by strict separators and the pieces are actual convex subsets. If one convex piece in one representation met two components in another representation, a separator between two intervening components would lie between two of its points and hence in that piece, but outside the union. This is impossible. Each piece is therefore contained in one piece of the other representation; applying the same argument in reverse forces equality. The argument also proves that if $X\subseteq Y$ are tame, each component of $X$ is contained in one component of $Y$.
For a nonempty interval $I\subseteq J$, intrinsic lower and upper endpoints compare in the expected order, since infima and suprema are monotone. Open or closed endpoint flags resolve equality: an included endpoint of $I$ must belong to $J$. Thus containment implies compatibility, after comparing endpoints of all other components through their strong order. For closed pieces this argument needs no E: their finite endpoints actually belong to them, and unbounded sides are detected by the absence of bounds. Compatible interval intersections have the larger lower endpoint and smaller upper endpoint, with inclusion at a tied endpoint precisely when both intervals include it. They are empty if these constraints are inconsistent.
Fix a finite chain $S=\{s_1<\cdots<s_n\}$ and list the nonempty cells
\[
 L_{<s_1},\{s_1\},(s_1,s_2),\{s_2\},\ldots,\{s_n\},L_{>s_n}.
\]
Use a single cell $L$ if $S$ is empty. For a set of selected cells, replace each maximal consecutive selected run by the interval whose outer endpoints and endpoint flags are those of that run. Let $F_S(m)$ be the union of those intervals. A run may contain points incomparable with an intermediate cut; therefore this replacement is essential. Distinct masks give distinct subsets, because intersection with each cell is either that entire cell or empty. The endpoint intersection formula gives
\[
 F_S(m)\cap F_S(n)=F_S(m\cap n),\qquad
 F_S(m)\subseteq F_S(n)\Longleftrightarrow m\subseteq n.
\]
For completeness, intersect two run intervals: the overlap has exactly the maximum lower and minimum upper cell indices, and exactly their endpoint flags. A nonempty overlap therefore corresponds to a nonempty overlapping index run. Taking all run pairs proves the displayed intersection formula. The containment implication also follows by testing any point in each nonempty cell; the reverse follows by containment of interval endpoints.
Refining $S$ by a compatible finite chain subdivides each old cell into consecutive new cells. Selecting all its subdivisions and then normalizing recovers the old interval, including points incomparable with intermediate new cuts. The intersection formula and the equality of outer endpoints prove that this gives Boolean embeddings. Every compatible finite family is represented in one such flag using all its endpoints. Its meet is actual intersection; its join is $F_S(m\cup n)$; its complement is $F_S(\text{all cells}\setminus m)$. To verify that the join is a least upper bound among ALL tame parts, let $Z$ be an arbitrary tame upper bound. The containment lemma makes $Z$ compatible with both operands. Their endpoints together with those of $Z$ are a chain, so refinement to that flag makes the assertion ordinary inclusion of masks. Complement is exactly the set of points outside $X$ comparable with every endpoint of $X$: inspect each component and the intervening gaps. It is the largest compatible disjoint part, is involutive and order reversing, and obeys both De Morgan identities. Finite pairwise compatible families generate finite Boolean subalgebras because there are finitely many endpoint cells. Arbitrary pairwise compatible families have every finite subfamily in such a block, which proves all finitary Boolean identities.
Closure fills each omitted endpoint of each selected nondegenerate run and merges the runs when these closures touch. It is a closed tame superset. Any closed tame upper bound is compatible by containment, and its mask in a common refinement must contain those endpoint cells; therefore this closure is least. It preserves compatible finite joins and preserves compatibility, because it introduces no new endpoint. Interior is the complement of the closure of the complement. The intersection of the two closures is the finite relative frontier: all component endpoints except an ambient extreme included in a non-singleton component. A singleton at an ambient extremum remains a frontier point. Removing the frontier from a nondegenerate run leaves a nonempty open interval, by density. Thus a tame set has empty tame interior exactly when all its components are singletons, exactly when it is finite; it then equals its frontier. These statements concern tame interior, not arbitrary topological connectedness.
For closed compatible unions alone, all endpoints are intrinsic without E, the same intersection and normalized-union operations give global meets and joins among CLOSED tame parts, and uniqueness of closed components holds. The topology need not call those parts closed.
\emph{Independent re-proof.} Instead of masks, intersect every pair of component intervals by the maximum/minimum endpoint rule. For a join, order all endpoint events and sweep from left to right: an uncovered strict gap separates components; at a touching point or an occupied intervening gap, convexity of any upper component forces the entire interval between the outermost occupied endpoints. The resulting intervals are forced in every upper bound and are themselves a tame upper bound. Endpoints omitted from an interval must be added in any closed upper bound by E and cofinality. This reproduces meet, join and closure directly. Boolean and De Morgan computations have also been checked exhaustively on up to eight abstract cells in \texttt{check\_finite.out}; those computations do not substitute for the endpoint proof.

\paragraph{P: canonical topologies and two further countermodels.}\label{lem:topology}
The strict intervals and rays form a basis in G. Given a point in the intersection of finitely many such sets, use the finite local cone lemma to find a common lower bound for their upper constraints still above the point, and a common upper bound for their lower constraints still below it. The resulting strict interval lies in the intersection. At an ambient endpoint use the appropriate ray. These sets cover L: a point without a predecessor or successor is the corresponding ambient endpoint by global directedness. A least element is closed because every other point has a strict predecessor above it or equal to it; its upper ray avoids the least element. The dual proves closedness of a greatest element. If there is no least element, a point in a left ray can be enclosed by a strict interval, and similarly for right rays without a greatest element.
There is an exact order-only separation criterion. Put $P(x)=L_{<x}$ and $U(x)=L_{>x}$. Since strict lower and upper rays generate the canonical topology,
\[
 x\in\operatorname{cl}\{y\}\quad\Longleftrightarrow\quad
 P(x)\subseteq P(y)\ \hbox{and}\ U(x)\subseteq U(y).
\]
Indeed every generating ray containing x contains y exactly when the corresponding displayed inclusion holds; finite intersections and arbitrary unions preserve this implication. Hence the topology is $T_0$ exactly when $x\mapsto(P(x),U(x))$ is injective. It is $T_1$ exactly when, for distinct x,y, the two componentwise inclusions never hold simultaneously. Equivalently, for every ordered pair $x\ne y$ there is $p<x$ with $p\not<y$, or $q>x$ with $y\not<q$. This is a necessary and sufficient condition in the order language, not a totality assumption. The doubled line fails injectivity; the strict plane satisfies the stronger antichain condition while remaining non-total.
On a total G-order, if x<y choose x<c<y. The disjoint rays below and above c separate x,y. The complement of a principal closed ray is the opposite strict ray, and the complement of [a,b] is the union of the exterior strict rays. Thus the canonical topology is Hausdorff and all principal closed pieces are closed. E also follows directly: if b is an upper bound of all predecessors of a but b<a, interpolate b<z<a to contradict it.
For a non-total Hausdorff example take P=$\mathbb R^2$, with p<q exactly when both coordinates of p are strictly smaller, and add equality. Antisymmetry and transitivity are immediate. Choose both coordinates farther down or up for global bounds, strictly between for interpolation, and between a and the coordinatewise minima of two points above a for the local lower-bound axiom; the dual verifies the other cone. Strict bounded intervals are exactly open rectangles. Hence the canonical topology is Euclidean and T\textsubscript{1}, although (0,1) and (1,0) are incomparable. E fails: every strict predecessor of a=(0,0) is below b=(0,1), whereas a is not below b. Thus b is an upper bound of the strict predecessors without bounding a. The weak coordinatewise order would not work: above (0,0), (1,0) and (0,1) have no common strict predecessor.
Let $\tau_c$ have principal closed pieces as a closed subbasis. Every basic $\tau_c$-open is the complement of a finite union of these pieces. For all sufficiently large n, (n,-n) belongs to that complement: it is in no bounded interval, in no lower ray (first coordinate), and in no upper ray (second coordinate). Thus no nonempty such basic open is contained in a bounded rectangle. A bounded rectangle is not $\tau_c$-open, and $\tau_c$ need not contain the canonical topology. For the supremum topology one must use complements of ALL members of the closed subbasis and finite intersections of those and canonical basic opens. The bare [a,b] printed on scan 24 is not that complement. On the real line, including [0,1] as an open generator strictly refines the claimed supremum, which is the usual topology.
For a T\textsubscript{0} but nonsober example take J to be all nondegenerate real compact intervals [l,r], l<r, ordered by strict separation r(I)<l(J), with equality added. Global bounds are remote short intervals; interpolation and both local cone conditions follow by putting a short interval between the adjacent real endpoint bounds. Its canonical basis is
\[ U_{s,t}=\{[l,r]:s<l\leq r<t,\ l<r\},\quad s<t, \]
with rays covered by such bounded sets. Two different intervals have different endpoint constraints, hence different neighborhood filters. Directly,
\[\operatorname{cl}\{I\}=\{J:J\supseteq I\}.\]
Indeed every basic open containing a larger interval contains its subinterval I; if J fails to contain I, choose a real bound between the offending endpoints to find a basic open containing J and missing I. The set F=\{J:0$\in$ J\} is closed: its complement consists of intervals wholly to the left or wholly to the right of 0, an open condition. If U\textsubscript{s,t} meets F, s<0<t. Two such basic opens contain a common small nondegenerate interval about 0. Thus any two open sets meeting F have intersection meeting F, which is precisely irreducibility. A generic point I of F would be contained in every interval of F, hence in their intersection \{0\}; this contradicts nondegeneracy. F has no generic point, so J is not sober.
\emph{Independent check.} For P, all basic-open assertions reduce to two real strict inequalities in each coordinate. For J, the specialization calculation is checked by the two separate possibilities l(I)<l(J) and r(J)<r(I); choosing a separating real threshold excludes I. Intersecting the explicit family [-1/n,1/n] proves the generic-point obstruction without any abstract sobriety theorem.

\paragraph{H: arbitrary closed hulls, separators and components.}\label{lem:hulls}
Even in L=[0,1], A=\{0\}$\cup$\{1/n:n$\geq$1\} has no least finite union of closed intervals containing it. Any such union H has a component containing infinitely many 1/n. Since its left endpoint is at most their infimum 0, that component contains [0,$\epsilon$] for some $\epsilon$>0. Choose an irrational c$\in$(0,$\epsilon$). Choose N with 1/N<c. Then [0,1/N] together with the finitely many singletons \{1\},\ldots,\{1/(N-1)\} contains A and excludes c. Thus H is not least. The same proof works on $\mathbb R$. This proves the manuscript's warning; it is not the counterexample to the stronger all-tame-hulls claim, which requires D.
Suppose actual a=$\sup$ T and b=$\inf$ T' exist. If T<T', every point of T' bounds T, so a$\leq$ b. A strict separator x satisfies a$\leq$ x$\leq$ b, with x$\ne$ a when a$\in$ T and x$\ne$ b when b$\in$ T'. Conversely these inequalities and endpoint exclusions imply t<x<t' for every t$\in$ T,t'$\in$ T': equality t=x would force x=a=t, and likewise at b. Therefore the separator set is [a,b] with exactly those excluded endpoints. It is nonempty iff a<b, or a=b belongs to neither T nor T'. Density handles the first case. This equivalence needs no E because actual supremum and infimum were assumed explicitly on scan 29.
Without E it need not have intrinsic endpoints or the asserted unique raccordement data: take T=\{a\textsubscript{0}\}, T'=\{b\textsubscript{0}\} at heights 0,1 in D. The separator interval has lower bounds a\textsubscript{0},a\textsubscript{1} with no greatest lower bound, and upper bounds b\textsubscript{0},b\textsubscript{1} with no least upper bound. Under E its intrinsic endpoints are a,b, and the endpoint flags show that it raccords to T and T' exactly as claimed. Uniqueness follows because any such interval has these forced endpoints and flags. Exterior separator sets are the corresponding exterior ray; their absence is characterized by the absence of a STRICT lower or upper bound. In [0,1], a left exterior ray may contain 0 and hence has a lower bound: ``sans minorant'' must mean without a strict lower bound. The strict qualifier appears just above the corollary, supporting the slip judgment.
The number of complementary components of a nonempty n-component tame part under E is $n-1+e_-+e_+$, where each e is 1 exactly when the corresponding strict exterior separator set is nonempty. Hence the difference in numbers is at most one. For X=$\varnothing$ the complement is L and the difference is again one; the table involving T\textsubscript{1},T\textsubscript{n} applies only to n>0.
Within a pseudo-troncon T with G and E, a nonempty proper tame part has nonempty relative frontier. Indeed a proper nonempty interval component must have an internal finite endpoint, and distinct components have a strict intervening gap. The endpoint is retained in the frontier unless it is an ambient end in a non-singleton component; if all finite endpoints were of this harmless kind, the sole component would be T. Therefore a clopen tame part is only empty or T. For a tame A with components T\textsubscript{i}, closure fills only their endpoint omissions; strict separators ensure $\overline{T_i}$$\cap$ A=T\textsubscript{i}. The complementary union of the other components is also relatively closed, so each T\textsubscript{i} is clopen in A and connected in this tame sense. Their finite clopen algebra has them as its atoms. This does not assert topological connectedness: a rational open interval has a cut at an irrational real number into two canonical clopen subsets, neither forming such a finite endpoint cut in $\mathbb Q$.
\emph{Independent check.} The separator proof can instead be read from the four endpoint-flag cases (both included, only left included, only right included, neither included). Each contributes precisely the indicated removal from [a,b]. Counting the n-1 internal gaps and the two optional exterior gaps then proves the count separately. The clopen-component assertion follows independently from the finite endpoint-cell construction F: a clopen mask must be constant throughout each component, since a change at an internal endpoint leaves that endpoint in both closures.

\paragraph{M: the abstract axioms, relativization, and a false hope.}\label{lem:abstract}
Here M is an abstract bounded poset with a reflexive symmetric compatibility relation. Comparable pairs are compatible. Compatible pairs have global meets and joins. There is the compatible complement x', the greatest compatible element disjoint from x, with x$\vee$ x'=1, order-reversing involution and compatibility invariant under either complement. Every finite pairwise compatible family is contained in a Boolean algebra for these global operations (the strengthened marginal axiom on scan 33). The distinguished closed elements include 0,1; each x has a least closed upper bound c(x); closure preserves compatible joins and c(x) is compatible with y whenever x is. These are the stipulated hypotheses of scans 34--37, not consequences asserted for every G in this proof. For an arbitrary pairwise compatible family, the finite-family condition also supplies a single generated Boolean algebra: take the union of the subalgebras generated by its finite subfamilies. These subalgebras are directed under inclusion, because all operations are the stipulated global meet, join and complement, and a Boolean block containing the union of two finite subfamilies contains both generated subalgebras. Every finite Boolean identity and every compatibility assertion about two generated elements is witnessed in one such block. Thus no unasserted infinitary block postulate is needed.
Fix A. In M\textsubscript{A}=\{x:x$\leq$ A\}, inherited compatible meet and join still lie below A, so are also the global relative meet and join. The top is A, the bottom 0. Since x$\leq$ A, x and A are compatible and so are x' and A. The relative complement is A$\wedge$ x'. If y$\leq$ A is compatible with x and disjoint from x, then y$\leq$ x', hence y$\leq$ A$\wedge$ x'. Conversely the latter is compatible with x and disjoint from it in the Boolean block containing x,A. That block also proves x$\vee$(A$\wedge$ x')=A and involutivity. A compatible family in M\textsubscript{A}, together with A, lies in a Boolean block; its relative operations are those in the principal Boolean ideal below A, which is itself a Boolean algebra with top A. This verifies the strengthened finite-family clause and every relative complement compatibility statement.
Declare x relatively closed iff x=A$\wedge$ c(x). This is equivalent to x=A$\wedge$ z for some closed z compatible with A: the forward implication uses z=c(x), which is compatible with A by the closure axiom. Conversely x$\leq$ z implies c(x)$\leq$ z, and hence x$\leq$ A$\wedge$ c(x)$\leq$ A$\wedge$ z=x. The least relative closed upper bound is c\textsubscript{A}(x)=A$\wedge$ c(x). It is relatively closed by the equivalent description; and if x$\leq$ A$\wedge$ z with z closed, then c(x)$\leq$ z. For compatible x,y below A, all of A,c(x),c(y) are pairwise compatible, so distributivity in their Boolean block yields
\[c_A(x\vee y)=A\wedge(c(x)\vee c(y))=c_A(x)\vee c_A(y).\]
The same block and closure-compatibility axiom show c\textsubscript{A}(x) compatible with y. Both relative extrema are closed. This verifies all axioms 1--7 in M\textsubscript{A}.
The Boolean ring identities on scan 34 hold on each compatible block: identify its finite subalgebra with subsets of its atoms. Meet is bitwise AND, x+y is XOR, and x$\vee$ y=x+y-xy (in characteristic two minus equals plus). The equation x$\leq$ y iff xy=x is the subset test. These identities do not define multiplication of incompatible elements.
For a counterexample to universal compatibility of clopens, take M=\{0,1,a,a',b,b'\}, with 0 below every element and 1 above every element, and four incomparable middle elements. Compatible pairs are exactly pairs in either block \{0,1,a,a'\} or \{0,1,b,b'\}. Complement swaps a,a', b,b', 0,1. Compatible meet and join are the Boolean ones; they are also global bounds in M. A pairwise compatible family lies in one block (unless it contains only 0,1). Take every element closed, c=identity. This verifies all axioms 1--7 including the strengthened clause, and every element is open too. Nevertheless a and b are incompatible. Thus the hope on scan 37 is false under precisely its stated abstract hypotheses.
If one ADDS that each clopen is compatible with every element, then clopens form a Boolean algebra: closure preserves compatible joins; the complement of a clopen is clopen; De Morgan gives meets. Its Stone space is the set of ultrafilters, with basic sets $\widehat x$=\{U:x$\in$ U\}. Boolean identities make them clopen and separate distinct ultrafilters. The ultrafilter lemma proves compactness: a cover with no finite subcover gives a family of complemented elements with the finite-intersection property, extendible to an ultrafilter missing the cover. This uses the Boolean prime ideal theorem, stated explicitly. These spectral points encode consistent clopen choices; calling them actual components in an unspecified ambient geometry is an interpretation, not an additional representation theorem. If the Boolean algebra is finite and 0$\ne$1, every ultrafilter contains exactly one atom; its clopens are disjoint joins of the finitely many atoms. Below such an atom the relative clopen algebra is trivial, so it is a component in this algebraic sense. Here relative clopens below a globally clopen a are exactly the globally clopen elements below a. Indeed x relatively closed and x$\leq$ a imply c(x)$\leq$ c(a)=a, whence x=a$\wedge$ c(x)=c(x). If its relative complement a$\wedge$ x' is relatively closed, it is globally closed by the same argument. The Boolean identity x'=(a$\wedge$ x')$\vee$ a', and join preservation of closure, then make x' globally closed. The converse follows immediately from the relative closure formula. Consequently a proper nonzero relative clopen below an atom would contradict minimality among the global nonzero clopens. When 0=1 the spectrum is empty.
\emph{Independent re-proof.} The map z$\mapsto$ A$\wedge$ z on each common Boolean block is a Boolean homomorphism to its principal ideal (relative complement follows by one distributive identity). The least-upper-bound definition of c then gives the relative reflector directly. This checks the relative construction without following the page's sequence. The explicit six-element tables and all finite Boolean identities through eight cells are separately checked by \texttt{check\_finite.py}; the table check includes every pairwise-compatible subset, not just pairs.

\paragraph{O: recovering compatibility from inclusion under E.}\label{lem:ordercompat}
This is OUR exact answer for the repaired concrete calculus, not an axiom of an arbitrary M. In the inclusion poset of actual tame subsets, atoms are exactly singleton subsets, since every nonempty tame subset contains a singleton. Thus the set of atoms below an element determines that element.
For two distinct atoms p,q, their join exists in M if and only if their points are comparable. If they are comparable, their two-point union is tame by density and is their least upper bound. If they are incomparable, any tame upper bound Y puts them in the same component I, since different components are strictly ordered. If I has a finite upper endpoint b, both points are strictly below b; local filtering below b gives a common upper bound c<b. Truncating I at c still contains p,q and gives a proper tame subset of Y. The discarded part is nonempty by interpolation towards b. If I has no greatest or finite upper endpoint, choose a common upper bound c of p,q by global filtering; it lies in I because it is above both and the lower endpoint, and I has points beyond c. Again truncate. Ambient greatest elements are covered by the first case. The other components of Y stay untouched and strongly separated. Thus every upper bound has a strictly smaller upper bound, and no join exists. This proof does not require E.
Write C(p,q) for equality or existence of that atomic join. Under E the canonical tame complement has the following description entirely in terms of atoms and inclusion:
\[
 \operatorname{At}(X')=\{p:p\not\leq X,\quad \forall q\in\operatorname{At}(X),\ C(p,q)\}.                 \tag{O1}
\]
To prove this, if p is outside X and compatible with its endpoints, it lies strictly on one side of each component of X, hence is comparable with every point of X. Conversely assume p is outside X and comparable with every point of X. On each convex component T all those comparisons have the same orientation: otherwise t<p<t' with t,t'$\in$ T would force p$\in$ T. If all t<p, p is an upper bound of T and therefore is above its actual supremum, when finite, by E. It also compares with the lower endpoint. The other orientation is dual. A component unbounded in the relevant direction cannot have p as such a bound; an ambient extremum is already comparable to p. Thus p compares with every component endpoint, as required. F ensures that the atom set in O1 is represented by a tame subset, and atoms uniquely determine this subset. This recovers X' from the inclusion poset alone.
Finally compatibility itself is recovered as follows. Elements X,Y are compatible iff there exist A,B,C pairwise orthogonal, in the sense A$\leq$ B', A$\leq$ C', B$\leq$ C', with
\[X=A\vee C,\qquad Y=B\vee C.\]
If X,Y are compatible, use A=X$\wedge$ Y', B=Y$\wedge$ X', C=X$\wedge$ Y in their Boolean block. Conversely orthogonality implies compatibility by inclusion and complement invariance; the three elements thus lie in a common Boolean block, which contains both displayed joins. The whole criterion uses only inclusion, its extrema/bounds, atoms, and the complement just recovered. This answers the broader order-definability question positively under E. It does not assert that the original ill-defined endpoint calculus on every G has a recoverable relation.
The simpler proposed test by existence of X$\wedge$ Y is false, even under E. Construct a split real line with a common left tail (-$\infty$,0), two incomparable copies of (0,1), and a common right tail (1,$\infty$), omitting 0 and 1. Order each tail and each branch by the real coordinate, every left-tail point below both branches and every right-tail point above both. Common tail bounds give directedness. Between comparable points interpolate on their branch or tail; for two points on different branches with a common strict predecessor use a point farther right in the left tail, and dually use the right tail. This verifies local filtering. E holds: below a branch point its cofinal same-branch predecessors exclude the other branch from being an upper bound; the same argument works for successors and for points of either tail. Thus G and E hold. Two singleton points on different branches are incompatible but their greatest lower bound in the inclusion poset is the empty set. No correction of ``exists'' to a merely set-theoretic intersection test fixes this.
\emph{Independent check.} Formula O1 says that complement is the polar for the graph on point atoms whose edges are the pairs admitting a join. The direct endpoint proof of F and the all-points comparison proof above compute the same polar independently. Applying the three disjoint Boolean pieces to the recovered polar verifies the final compatibility criterion without referring to endpoints or an orientation.

\paragraph{V: the genuine ficelle counterexample and the corrected theorem.}\label{lem:ficelle}
The supplementary primary source is the original Montpellier 156-2 PDF, scans 2--10. The actual definitions are these. A ficelle is $(H,\{A,B\})$, with disjoint nonempty closed nonopen ends, such that $(H\setminus A)/B$ and $(H\setminus B)/A$ are connected. Each end is regularly immersed: after contracting it, the resulting point has a closed conical neighborhood on a compact base, with open cone interior. The word compact is visible on 156-2/002.18. An internal lieu C is closed, disjoint from the ends, and $H\setminus C$ is the disjoint open union $U_{A,C}\amalg U_{B,C}$, containing A and B respectively, with $U_{A,C}/A$ and $U_{B,C}/B$ connected. The lieu is regularly immersed. No connectedness of C is required. The two component spaces are $X_{A,C}=U_{A,C}\cup C$ and $X_{B,C}=U_{B,C}\cup C$, both closed. The oriented relation $C<D$ means that C is an internal lieu of the left component ficelle at D (with the usual end conventions).

The component spaces themselves are indeed ficelles; this does not imply promotion of their lieux. For example, $X_{A,C}\setminus C=U_{A,C}$ gives one connectedness condition. The other, $(X_{A,C}\setminus A)/C$, is a continuous quotient of the connected space $(H\setminus A)/B$ obtained by collapsing $X_{B,C}/B$; finite closed-set gluing identifies the quotient. A retains its local neighborhood. At C, after its contraction, the punctured cone is separated by the two sides; each radial interval lies wholly in one side. Thus its compact base splits into two clopen compact parts, giving conical neighborhoods in both component spaces. If C were open in one component, its opposite open side would become a nontrivial clopen subset of the relevant globally connected end quotient, a contradiction. This verifies nonopenness and the remaining component-ficelle conditions without the false promotion claim.

\emph{Correction history.} The earlier draft incorrectly promoted component-ficelle conical cuts to global lieux without proving their global connectedness condition. That assertion is false under these definitions. A level-t retraction also retains the level and does not alone prove connectedness of the punctured outside. The following counterexample replaces the earlier positive universal verdict; the two legitimate additional-hypothesis results are proved below.

\emph{A compact finite-tree ficelle.}
Take a main interval $[0,1]$, with $a=0$, $b=1$, and attach a side interval $[0,1]_{\rm h}$ at its zero endpoint to the main point $v=1/3$. Write p for the free hair endpoint, at hair coordinate 1. Give the resulting finite tree its usual graph topology. It is compact and Hausdorff. Removing either a or b leaves a path-connected tree; contracting the other point changes nothing. The two original ends have closed half-interval cone neighborhoods with one-point compact base, and are closed, nonempty, nonopen and disjoint. Thus every ficelle hypothesis holds.

Every internal lieu has exactly the following form. Its main-path part is a closed interval $I=[l,r]$ with $0<l\leq r<1$. Its hair part is either empty or a terminal interval $[s,1]_{\rm h}$; if v belongs to I, the entire hair must be included. Conversely every such set is a lieu, subject to that condition. To prove necessity, a set separating a from b must meet the unique main path. Two separated pieces of its main intersection would leave an intervening component containing neither a nor b, so that intersection is an interval. A removed nonterminal portion of the hair would leave a distal component containing neither end; hence the removed hair is terminal. If v is removed, any remaining hair is such a third component, so all of it must be removed. These arguments also show that the complement of a set in the list has exactly two connected open components containing a and b. Finally regular immersion is explicit: the set is a finite union of closed graph intervals; after collapsing it, a sufficiently small closed neighborhood is a finite star, hence a cone on a finite compact set, with its cone interior open. This verifies sufficiency and every source hypothesis, including the extra regular-immersion clause.

This also identifies the false source-level promotion step. In the left component at $D_s$, choose $0<\epsilon<\min(s,d-v)$ and let E consist of the main point $d-\epsilon$ and the hair point $s-\epsilon$. In that component, removing E leaves an A-side and two distal arms; contracting the whole end $D_s$ joins the two distal arms. Hence E is a regular lieu of the component ficelle. In the global tree, removing E leaves three components: the A component, the B component, and the distal hair. The last contains neither original end, so E is not a global lieu. The statement on 156-2/008.3--4 that every component lieu is a global lieu is therefore false under the written definitions. Its ink is retained; no connectedness assumption is inserted into it.

Call a lieu bare when its hair part is empty. A bare interval lies wholly to one side of v. Two non-bare lieux both contain p, so cannot be strictly ordered. For distinct internal C,D the order is as follows:
\[
\begin{array}{ll}
C,D\text{ bare}:& C<D\ \Longleftrightarrow\ r_C<l_D,\\
C\text{ non-bare},\ D\text{ bare}:&
C<D\ \Longleftrightarrow\ \max(r_C,v)<l_D,\\
C\text{ bare},\ D\text{ non-bare}:&
C<D\ \Longleftrightarrow\ r_C<\min(l_D,v).
\end{array} \tag{V1}
\]
For example, in the last case C must be to the left of the main portion of D and also to the left of v, so that all of the terminal hair in the end D is on its D-side. These conditions are necessary by the component-ficelle definition; under them the two relevant tree pieces are connected, also after the prescribed end contractions, so they are sufficient. The other two cases are the same direct path check. End a is least and b is greatest.

The entire lieu order in this example satisfies G, not just a weaker finite surrogate. Global bounds are a,b. Density in the bare/bare case follows by inserting a singleton main cut in the gap, choosing it different from v. If the lower lieu is non-bare, insert a bare singleton between $\max(r_C,v)$ and $l_D$; if the upper lieu is non-bare, insert one between $r_C$ and $\min(l_D,v)$. For local lower bounds in a strict upper cone, if its apex is non-bare, both upper points are bare on the right; insert a bare singleton between the apex's effective upper endpoint $\max(r_C,v)$ and the smaller of their left endpoints. If the apex is bare to the left of v, insert a bare singleton between its right endpoint and the minimum of v and the two upper left endpoints. If the apex is bare to the right of v, both upper points are bare on the right and the ordinary interval choice works. At a, take a sufficiently small positive main singleton. The dual construction proves the strict-lower-cone axiom, including b. The choices are possible by the strict inequalities in V1 and real density. Thus all G axioms hold on the full, infinite lieu set.

Fix $d=3/4$ on the main path, and, for $0<s\leq1$, set
\[
 D_s=\{d\}\cup[s,1]_{\rm h}.
\]
By V1, every $D_s$ has exactly the same strict predecessors: a and the bare main intervals whose right endpoint is less than v. Its strict successors are b and the bare main intervals whose left endpoint is greater than d. Consequently distinct $D_s$ have identical neighborhoods for the canonical topology generated by strict rays and intervals. The lieu space is not $T_0$. For $s<t$, $D_t$ is a proper subset of $D_s$, but $D_t\in\operatorname{cl}\{D_s\}$, although $D_t$ does not contain $D_s$. This directly refutes the claimed reverse-inclusion equivalence as well. The connected lieu
\[
 J=[v,d]\cup[0,1]_{\rm h}
\]
has those same strict cones; the defect is not remedied by observing that some of the loci are connected.

The finite-grid program verifies the classification and V1 by independently computing connected components and relative end contractions. It checks the equal strict cones of two sample $D_s$ and J. Its output is supplementary: density and the infinite G assertions were proved above, not inferred from a finite enumeration.

\emph{What remains true without an extra hypothesis.}
For arbitrary genuine ficelle lieux, subset containment still implies specialization:
\[
 C\subseteq D\quad\Longrightarrow\quad D\in\operatorname{cl}\{C\}. \tag{V2}
\]
Here is the necessary order argument in full. If $E<D$, then $E\subset U_{A,D}$. Since $U_{B,D}/B$ is connected and disjoint from E, it lies in the B-side of E. Thus $U_{A,E}\subset X_{A,D}$. Within $X_{A,D}\setminus E$, the local lieu definition separates A from the whole end D. Connectedness of $U_{A,E}/A$ puts this entire global A-side in the local A-side, so it misses D. Hence $D\subset U_{B,E}$.

We also need a restriction lemma. Suppose global internal E,C satisfy $E\subset U_{A,C}$ and $C\subset U_{B,E}$. Connectedness after contracting A gives $U_{A,E}\subset U_{A,C}$; the dual side argument gives $X_{B,C}\subset U_{B,E}$. In $X_{A,C}\setminus E$, the sides are $U_{A,E}$ and $W=U_{B,E}\cap X_{A,C}$, containing A and C. The first is connected after contracting A. The closed cover
\[
 U_{B,E}=W\cup X_{B,C},\qquad W\cap X_{B,C}=C
\]
shows that $W/C$ is the quotient of the connected space $U_{B,E}/B$ obtained by collapsing $X_{B,C}/B$. The quotient identification follows from finite closed-set gluing, so $W/C$ is connected. Regular immersion of E restricts to $X_{A,C}$ because it contains the open neighborhood $U_{A,C}$ of E. Therefore E is a lieu of that component ficelle: $E<C$.

Now if $C\subseteq D$ and $E<D$, connectedness of $U_{A,D}/A$, disjoint from C, gives $E\subset U_{A,C}$; the previous paragraph gives $C\subset D\subset U_{B,E}$. The restriction lemma yields $E<C$. Dually $D<F$ implies $C<F$. Every generating strict ray containing D therefore contains C, proving V2. At an ambient end, the only lieu contained in that end is itself. The exact general specialization relation, without a false subset description, is the strict-cone-inclusion formula in P.

\emph{A substantive sufficient repair, weaker than requiring every lieu connected.}
For an internal target D impose the additional condition
\[
 (W_D)\qquad X_{A,D}/A\ \hbox{and}\ X_{B,D}/B\ \hbox{are connected}. \tag{V3}
\]
These are the \emph{unpunctured} component spaces. The original definition assumes only the corresponding punctured-side and opposite-end-contracted connectedness, and does not imply V3. In the tree counterexample $X_{B,D_s}/B$ is the disjoint union of the right main piece and the terminal hair segment.

Under V3 the conical cuts around D are global lieux, and
\[
 D\in\operatorname{cl}\{C\}\quad\Longleftrightarrow\quad C\subseteq D. \tag{V4}
\]
To prove this, first note an elementary closed-quotient fact. If a closed surjection has connected fibers and connected target, its domain is connected: a nontrivial clopen separation of the domain would put each fiber on one side, and its two closed images would separate the target. Collapsing a closed connected subset is such a closed surjection. Closedness follows because the saturation of a closed set is either itself or its union with the collapsed closed subset.

Work in the left component ficelle at D. After contraction of D, choose a closed conical neighborhood with compact base and open cone interior. It may be shrunk away from the other end A: a finite product-neighborhood subcover of the compact base at radius zero gives a uniform positive radius. At a small positive radius t, the level $S_t$ is a closed regularly immersed set, with a two-sided product collar.

Its outside A-side after contracting A is connected, including with $S_t$ removed. For every $u>t$, the radial map $r\mapsto\max(r,u)$ on the punctured component, extended by the identity outside the cone, descends after contracting A and has connected image by the ficelle axiom. These images all contain the contracted A point. Their union as u decreases to t is exactly the outside with $S_t$ removed, hence is connected. The pasting is continuous because the cone neighborhood and its exterior are a closed cover and the maps agree at its boundary.

The global other side is the open truncated left cone including D, glued along D to $X_{B,D}$. After contracting B, its closed subspace $X_{B,D}/B$ is connected by V3. Collapsing that whole subspace leaves precisely the open cone with apex D, a connected space. The closed-quotient fact makes this global side connected after contraction of B. The quotient identification again uses the finite closed cover, whose intersection is D. Thus $S_t$ satisfies the \emph{global} lieu definition. In the left component it is a lieu as well: the side toward D becomes an open cone after contracting D, and the other side was just checked. Its two component-ficelle connectedness conditions also follow by collapsing the truncated cone or reversing its cone coordinate. Hence $S_t<D$. The dual construction gives a global cut above D.

Choose these two radii small enough to avoid any prescribed $h\notin D$: on its side, if h lies in the cone, take a radius below its positive coordinate; otherwise a sufficiently small truncated cone already misses it. Every lieu between the two cuts is contained in the resulting closed strip, by the side containments proved before V2. That strip avoids h. If $C\not\subseteq D$, use $h\in C\setminus D$ to obtain a canonical neighborhood of D excluding C. Together with V2 this proves V4.

Connectedness of D is one useful sufficient condition for V3. For either component ficelle, contracting both ends gives a connected reduced ficelle: its complement of one end is connected and the missing end is in its closure because it is not isolated. The quotient from $X_{A,D}/A$ to this reduced ficelle collapses the closed connected image of D, so the closed-quotient fact proves $X_{A,D}/A$ connected; dually on the right. Thus V4 holds at every connected target D. If V3 holds for every internal lieu, the whole canonical lieu topology is $T_0$: equal neighborhood filters give both subset inclusions. This last universal premise is not automatic in a general ficelle. At an ambient end, small global end-cone level cuts separate it from every other closed lieu: compactness shrinks the cone away from that lieu, the punctured outside is connected by the same union of radial retraction images, and the inner side becomes a cone after contraction of the end. Thus no ambient end shares its neighborhood filter with another lieu.

\emph{The genuine segment example used elsewhere.}
For $H=[0,1]$, $A=\{0\}$, $B=\{1\}$, the internal lieux are exactly the intervals $[l,r]$, $0<l\leq r<1$. Necessity follows because a hole would create a third complement component; sufficiency and regular immersion follow from the ordinary interval quotient. They are connected, so V3 holds. Their order is strict separation $r(C)<l(D)$, with equality and the two ends added. Canonical specialization is reverse interval inclusion, as is also computed directly in P. For $a=[3/10,1/2]$, $b=\{4/5\}$ and $x=\{2/5\}$, we have $a\notin[x,b]$, but every lieu z with $a<z<b$ lies in $[x,b]$. Thus the open set $L\setminus[x,b]$ in the supremum topology contains a and no right subinterval $(a,\epsilon)$. This retains the valid narrow counterexample on scan 24; it does not validate the universal ficelle assertion on scan 22.

\emph{Independent verification.} On the tree, compute the two strict cones directly from the three cases in V1; their independence of the terminal-hair parameter proves the obstruction without a neighborhood approximation argument. The separate finite graph computation derives order from connected-component tests, not from V1. For the corrected theorem, the closed-quotient argument checks the previously missing global connectedness, while the union over $u>t$ checks the separately missing punctured-side condition. The specialization formula in P then independently translates the resulting order neighborhoods. No connectedness of the removed locus, or promotion of a component lieu, is silently assumed.


\paragraph{N: what the provisional network axioms do and do not imply.}\label{lem:networks}
Let L=$\mathbb Q$ and S\textsubscript{0}=$\mathcal P_2$(L). Every finite subset is tame and C(T)=L$\setminus$ T. Give every C(T) the indiscrete equivalence relation, with one class A\textsubscript{T}. In one system give it frontier empty; in another give it frontier T. Dec\textsubscript{1} holds (one class). If T$\subset$ S, the single smaller class is contained in the single larger class; the frontier condition is respectively empty$\subset$ empty or S$\cap$ T=T$\subset$ T. There is no maximal finite tame subset, so the finite-maximal-cut axiom holds vacuously. These two systems have exactly the same S\textsubscript{0} and every R\textsubscript{T} but different frontiers. Thus the displayed provisional axioms do not make the frontier datum redundant.
For higher partitions keep empty frontiers, and compare the all-indiscrete system with this system: for |T|$\leq$2 use the indiscrete relation; for |T|$\geq$3 partition L$\setminus$ T into its negative and nonnegative parts. Both parts are nonempty after any finite removal. Adding cuts either retains the same two sign classes or refines the single old class into them. Hence Dec\textsubscript{1} and Dec\textsubscript{2} hold, as does the vacuous maximal-cut axiom. The two systems agree on every two-point cut but not on three-point cuts. This refutes the second forecast under the displayed provisional axioms. It does not refute the programme of finding additional axioms, whose exact ordered-troncon version is settled by A and R.
The discrete structure has equality as every R\textsubscript{T} and empty frontiers. Dec\textsubscript{2} is automatic. Dec\textsubscript{1} applied to T=$\varnothing$ says that L has finitely many singleton classes, hence is finite; conversely a finite L makes every quotient finite. For finite L the unique maximal tame subset is L, so the maximal-cut axiom holds. That axiom is a proposed requirement, not a theorem from S\textsubscript{0}: for example a graph consisting of one edge and an isolated vertex has a maximal tame edge that is not all of L.
\emph{Independent check.} In the two partition examples, every refinement comparison reduces to the three cases |T|,|S|$\leq$2; |T|$\leq$2<|S|; or 3$\leq$|T|$\leq$|S|. The corresponding relations are one-to-one, two-to-one, and sign-preserving. The frontier examples reduce to the single identity S$\cap$ T=T. This directly verifies all provisional clauses without inferring consistency from their names.

\paragraph{Z: zero-dimensional forms.}\label{lem:zero}
A set P with the discrete topology is a zero-dimensional form as defined on 156-1/002; no finiteness or nonemptiness is required there. A model is (L,$\sim$) with an equivalence relation. Define q:L$\to$ L/{$\sim$} by q(x)=[x]. Reflexivity, symmetry and transitivity imply that the classes partition L, so equality of classes is equivalent to x$\sim$ y. Give the quotient set the discrete topology (the text specifies a set quotient, not the quotient of an unspecified topology on L). If a map f:L$\to$ S is constant on equivalence classes, define $\bar f$([x])=f(x). This is well defined and is the unique map with $\bar f$ q=f. Conversely every factorization has that constancy property. Every set S is obtained from the diagonal relation on S; a model need not be its form, as a two-element indiscrete-equivalence model has a one-element quotient. No pro-object, transition map, higher-dimensional gluing rule, or category of arbitrary forms is supplied by these sentences.
\emph{Independent check.} Define the quotient instead as the image of L under x$\mapsto$\{y:y$\sim$ x\} in $\mathcal P$(L). Equality of these images is exactly equivalence, so this construction yields the same set and the same factorization. This checks existence independently of a choice of representatives.

\paragraph{B: the full cell algebra and the actually generated algebra.}\label{lem:generated}
Under G and E, F constructs a finite Boolean algebra B\textsubscript{S} for every finite endpoint chain S, with the nonempty endpoint/gap/tail cells as its atoms (in its Boolean order). If X\textsubscript{1},\ldots,X\textsubscript{k} lie in it, attach to each atom c its membership vector $v(c)=(1_{c\leq X_i})_{i=1}^k$. Every Boolean expression in the X\textsubscript{i} has the same truth value on atoms with equal vectors, by induction on the expression. Conversely the intersection over i of X\textsubscript{i} when v\textsubscript{i}=1 and X\textsubscript{i}' when v\textsubscript{i}=0 has exactly the atoms of that vector. The nonzero such intersections are therefore exactly the atoms of the generated subalgebra. It is all of B\textsubscript{S} if and only if the realized vectors are pairwise distinct.
On L=$\mathbb R$, one generator X=[0,1] has precisely the algebra \{$\varnothing$,X,X',L\}, with two atoms X and X'. But S=\{0,1\} gives the five nonempty cells (-$\infty$,0),\{0\},(0,1),\{1\},(1,$\infty$), and B\textsubscript{S} has 2\textsuperscript{5}=32 elements. Thus the cells are not in general the atoms of the algebra generated by the original family. Even the union of frontiers does not determine that subalgebra: X=\{0\} and Y=(-$\infty$,0] have the same frontier \{0\} and generate different four-element algebras. The exact corrected statement is the vector criterion above, not an identification of generated and ambient cell algebras.
\emph{Independent check.} A Boolean expression is a truth table on the k generator bits; its distinct evaluations are exactly the realized vectors. This truth-table argument independently yields the same atom classes and cardinality 2\textsuperscript{r}, where r is the number of realized vectors. The one-generator and five-cell counts are also computed in \texttt{check\_finite.out}.

\paragraph{H2: the cancelled ordinary-difference assertion remains false under E.}\label{lem:ordinarydifference}
Use the regular split line of O, with common tails (-$\infty$,0),(1,$\infty$), the two copies of (0,1), and no points at 0 or 1. Let a be the height-1/2 point on the first branch. Set X=L, Y=$L_{\leq a}$, both closed tame and Y$\subset$ X. The ordinary difference contains the entire second branch, the first branch above a, and the right tail. For every left-tail point p<0, the closed tame ray $L_{\geq p}$ is a superset of this difference. A least closed tame hull H would lie in the intersection of all these rays, which contains no point of the left tail. Yet no finite union of closed pieces contained there covers the entire second branch: every component meeting that branch either has its lower endpoint on the branch, at a positive height, or has its lower endpoint on the excluded left tail. It cannot have a lower endpoint on the other branch, since then its interval would miss the second branch. A lower-unbounded component would include the left tail. Finitely many permitted positive lower endpoints have a positive minimum, leaving the initial second-branch segment uncovered. This contradiction proves that even G+E does not ensure a hull for the ORDINARY difference.
The valid replacement used by the next manuscript sentence is the compatible difference C\textsubscript{\rm mod}(Y,X), together with E for its closure. Its full proof is F. On a totally ordered G, compatible and ordinary difference coincide; cutting at finitely many retained endpoints gives a finite union of intervals, and its endpoint closure is the least closed tame hull. Thus the cancelled assertion is true on every total G and false even on some endpoint-regular non-total G.
\emph{Independent check.} For the chosen example the compatible difference consists only of the first-branch ray above a together with the common right tail. Its closure is $L_{\geq a}$, which visibly omits the second branch. This verifies both why the author's replacement changes the problem and why using the repaired tame-complement theorem to prove the cancelled assertion would be invalid.


% Per-statement applications and second checks
\section*{H0K8.1}\label{statement:H0K8.1}
\paragraph{Locator.} 156-1/002.3

\paragraph{His words.} Il propose une « construction récurrente » des formes par la dimension; la page commence par la dimension zéro puis passe aux modèles de dimension un.

\paragraph{How he left it.} programme

\paragraph{Precise statement (ours).} A recursive definition of forms in all dimensions is requested. Higher stages, morphisms, equivalence of models and the recursion step are not defined here. The supported zero-dimensional versions are the next two rows.

\paragraph{Status.} programme

\paragraph{Complete proof / counterexample application (ours).} Lemma Z identifies exactly the two defined constructions. A recursion would require its induction data and transition rule; these are absent. This row does not assert their nonexistence. The two precise versions below are separately decided.

\paragraph{Correction and reconstruction (ours).} We separate an unfinished programme from its fully specified zero-dimensional base. No pro-object interpretation is inserted.

\paragraph{Independent second check.} Documentary second check: search of the complete 156-1 transcription located the one-dimensional F1–F6 model axioms, whose object is H0KF; none supplies a category or recursive equivalence relation on general forms.

\paragraph{Confidence.} High for mathematics; decisive wording checked in READINGS.tsv.

\section*{H0K8.2}\label{statement:H0K8.2}
\paragraph{Locator.} 156-1/002.6

\paragraph{His words.} Une « forme de dim 0 » est un ensemble discret; ses éléments sont les « lieux » de la forme.

\paragraph{How he left it.} stated

\paragraph{Precise statement (ours).} A zero-dimensional form is a set P equipped with the discrete topology. The sentence imposes neither finiteness nor nonemptiness.

\paragraph{Status.} True as written

\paragraph{Complete proof / counterexample application (ours).} Lemma Z constructs the discrete topology \(\mathcal P(P)\) on any set and explains that this is a definition, with no further structure. The empty set and arbitrary cardinalities satisfy it.

\paragraph{Correction and reconstruction (ours).} Our use of the discrete topology explicates « ensemble discret »; it does not equip the model L with a topology.

\paragraph{Independent second check.} The diagonal equivalence relation on P produces the same object by the independent quotient construction in Z.

\paragraph{Confidence.} High for mathematics; decisive wording checked in READINGS.tsv.

\section*{H0K8.3}\label{statement:H0K8.3}
\paragraph{Locator.} 156-1/002.6m

\paragraph{His words.} Un « modèle » est un ensemble L de lieux muni de la relation d’équivalence « non disjoints »; la forme du modèle est l’ensemble quotient.

\paragraph{How he left it.} stated

\paragraph{Precise statement (ours).} For an equivalence relation \ensuremath{\sim} on L, the associated form is the set L/\ensuremath{\sim}, regarded as discrete. The model and its associated form are distinct constructions.

\paragraph{Status.} True as written

\paragraph{Complete proof / counterexample application (ours).} Lemma Z proves existence, the equality criterion for classes, and the universal factorization property of \(q:L\to L/{\sim}\). A two-point model with all points equivalent has a one-point form, so identifying model and form is not licensed.

\paragraph{Correction and reconstruction (ours).} The universal property is ours. The margin states the quotient itself; it states no pro-object construction.

\paragraph{Independent second check.} The independent power-set image construction in Z avoids representatives and yields the same quotient.

\paragraph{Confidence.} High for mathematics; decisive wording checked in READINGS.tsv.

\section*{H0KB.1}\label{statement:H0KB.1}
\paragraph{Locator.} 156-3/002.3

\paragraph{His words.} Il veut « une autre axiomatique » de la structure de réseaux sur l’ensemble L des lieux, en termes de découpages.

\paragraph{How he left it.} programme

\paragraph{Precise statement (ours).} Specify tame pairs, equivalence relations on finite tame complements and their frontiers, with sufficient axioms to recognize networks. The general class is not yet axiomatized; the displayed provisional systems and the ordered-tronçon recognition problem are the supported precise versions.

\paragraph{Status.} programme

\paragraph{Complete proof / counterexample application (ours).} Lemma N checks the provisional systems; lemmas R and A settle the exact ordered subclass. There is no defined general-network class whose equivalence to the earlier segment axioms can be inferred from this heading alone.

\paragraph{Correction and reconstruction (ours).} The general programme remains descriptive; every explicit mathematical version below is separately decided.

\paragraph{Independent second check.} Compared the data on scans 2–6 with the concluding exercise on 41: the latter explicitly restricts to the ordered situation.

\paragraph{Confidence.} High for mathematics; decisive wording checked in READINGS.tsv.

\section*{H0KB.2}\label{statement:H0KB.2}
\paragraph{Locator.} 156-3/002.6

\paragraph{His words.} S\textsubscript{0}\ensuremath{\subset}P\textsubscript{2}(L) est l’ensemble des « 0-sphères plongées »; une partie finie T est modérée si P\textsubscript{2}(T)\ensuremath{\subset}S\textsubscript{0}. Deux lieux égaux sont aussi en position relative modérée.

\paragraph{How he left it.} stated

\paragraph{Precise statement (ours).} A symmetric irreflexive graph on L specifies distinct compatible pairs; finite tame sets are its finite cliques, including the empty set and singletons. Reflexive compatibility adds equality.

\paragraph{Status.} True as written

\paragraph{Complete proof / counterexample application (ours).} For any \(S_0\subseteq\mathcal P_2(L)\), the clique condition is well defined and hereditary under subsets. The union of two finite cliques is a clique precisely when every cross pair is equal or belongs to \(S_0\). This proves the stated relative-compatibility rule.

\paragraph{Correction and reconstruction (ours).} His explanation through the earlier segment model is a link to H0KF, not a new axiom silently imposed here.

\paragraph{Independent second check.} Independently inspect each two-element subset of the union: it lies in the first set, second set, or is a cross pair.

\paragraph{Confidence.} High for mathematics; decisive wording checked in READINGS.tsv.

\section*{H0KB.3}\label{statement:H0KB.3}
\paragraph{Locator.} 156-3/002.10

\paragraph{His words.} Dans l’axiomatique par les segments, une paire est modérée si un lieu est singulier, ou si les lieux réguliers sont sur des lignes différentes, ou sont les extrémités d’un segment dans la même ligne.

\paragraph{How he left it.} stated

\paragraph{Precise statement (ours).} Given the earlier network data of singular places, regular lines, and two-ended segments, this rule constructs the symmetric graph S\textsubscript{0} on distinct places. It is an interpretation of the datum, not yet an equivalence theorem for the provisional cut axioms.

\paragraph{Status.} True as written

\paragraph{Complete proof / counterexample application (ours).} Each of the three conditions is invariant under exchanging the two places. Restricting to two-element subsets excludes equality, and adding equality gives the reflexive compatibility relation already stipulated. Thus the construction of S\textsubscript{0} is well defined from those earlier data. Nothing in this bracket proves that arbitrary provisional cut systems arise from such segment networks; N supplies countermodels to weaker claimed consequences.

\paragraph{Correction and reconstruction (ours).} We preserve the earlier model’s terms. H0KF owns its F1–F6 axioms; H0KG/H0KH own the orientation and regular-network constructions. The bracket defines a graph from them without settling a converse.

\paragraph{Independent second check.} The three alternatives are checked directly against the bracket on scan 2 and the definition of singular/regular places in the searched 156-1 model passages.

\paragraph{Confidence.} High for mathematics; decisive wording checked in READINGS.tsv.

\section*{H0KB.4}\label{statement:H0KB.4}
\paragraph{Locator.} 156-3/003.4

\paragraph{His words.} C\_L(T) est l’ensemble des t\ensuremath{\notin}T tels que T\ensuremath{\cup}\{t\} soit modéré; il l’appelle « complémentaire modéré ».

\paragraph{How he left it.} stated

\paragraph{Precise statement (ours).} For a finite clique T, C(T)=\{x\ensuremath{\in}L\ensuremath{\setminus}T: every t\ensuremath{\in}T is compatible with x\}. It can omit points outside T.

\paragraph{Status.} True as written

\paragraph{Complete proof / counterexample application (ours).} Every two-point subset of \(T\cup\{x\}\) already lies in T or is \(\{t,x\}\). Thus the displayed condition is exactly the cross-pair condition, with x excluded from T. This proves both descriptions.

\paragraph{Correction and reconstruction (ours).} No set-theoretic complement is substituted for the manuscript’s tame complement.

\paragraph{Independent second check.} For a singleton in D, the other twin is outside T but outside C(T), explicitly checking that the word complement is qualified.

\paragraph{Confidence.} High for mathematics; decisive wording checked in READINGS.tsv.

\section*{H0KB.5}\label{statement:H0KB.5}
\paragraph{Locator.} 156-3/003.6m

\paragraph{His words.} Si T\ensuremath{\prime}\ensuremath{\subset}T, alors C(T)\ensuremath{\subset}C(T\ensuremath{\prime}).

\paragraph{How he left it.} stated

\paragraph{Precise statement (ours).} T\ensuremath{\prime}\ensuremath{\subset}T implies C(T)\ensuremath{\subset}C(T\ensuremath{\prime}) for finite tame sets in any compatibility graph.

\paragraph{Status.} True as written

\paragraph{Complete proof / counterexample application (ours).} If x avoids T and is compatible with every point of T, it avoids T\ensuremath{\prime} and is compatible with every point of T\ensuremath{\prime}. This is lemma C.

\paragraph{Correction and reconstruction (ours).} Our expansion of the displayed definitions and their hypotheses; no additional claim is attributed to the manuscript.

\paragraph{Independent second check.} Check the two requirements separately: nonmembership descends with subset inclusion, and the universal compatibility condition restricts.

\paragraph{Confidence.} High for mathematics; decisive wording checked in READINGS.tsv.

\section*{H0KB.6}\label{statement:H0KB.6}
\paragraph{Locator.} 156-3/003.6m

\paragraph{His words.} Il demande si T est connu par C(T), puis écrit « Oui » pour l’inverse de l’inclusion, et revient au NB : T est connu par C(T).

\paragraph{How he left it.} sketched

\paragraph{Precise statement (ours).} The proposed map T\ensuremath{\mapsto}C(T) is injective and reflects reversed inclusion for arbitrary finite tame sets.

\paragraph{Status.} False

\paragraph{Complete proof / counterexample application (ours).} In D of lemma D, distinct singleton twins have identical tame complements. All G hypotheses and their singleton tameness are verified there, so the failure persists in the intended ordered subclass. The strongest direct general-graph repair is: if S\ensuremath{\cup}T is tame, then C(S)\ensuremath{\subset}C(T) iff T\ensuremath{\subset}S. For t\ensuremath{\in}T\textbackslash{}S, compatibility of the union gives t\ensuremath{\in}C(S) but t\ensuremath{\notin}C(T), a contradiction.

\paragraph{Correction and reconstruction (ours).} Add the missing hypothesis S\ensuremath{\cup}T tame. It is exactly the condition needed by his attempted witness argument; it is not present in the written assertion.

\paragraph{Independent second check.} The twin computation checks equality point by point at height 0 and at every other height. The repaired converse is also verified by its explicit witness t.

\paragraph{Confidence.} High for mathematics; decisive wording checked in READINGS.tsv.

\noindent Refined statement identifier: H0KB.6.correction (the corrected statement and its proof in this section).

\section*{H0KB.7}\label{statement:H0KB.7}
\paragraph{Locator.} 156-3/003.15

\paragraph{His words.} Le NB écrit « s\ensuremath{\ne}t ou \{s;t\}\ensuremath{\in}S\textsubscript{0} » pour la position relative modérée.

\paragraph{How he left it.} stated

\paragraph{Precise statement (ours).} Taken literally, this makes all distinct points compatible and fails reflexivity. The preceding margin on scan 2 instead says s=t or \{s,t\}\ensuremath{\in}S\textsubscript{0}.

\paragraph{Status.} Slip

\paragraph{Complete proof / counterexample application (ours).} With \(S_0=\varnothing\) on a two-point set, the printed disjunction holds for the distinct pair, contrary to its definition, and fails for equal points. Replacing \(\ne\) by \(=\) gives exactly the finite-clique compatibility rule proved in H0KB.2.

\paragraph{Correction and reconstruction (ours).} Our slip judgment: replace \ensuremath{\ne} by =, using the explicit reflexive margin on 156-3/002.

\paragraph{Independent second check.} The cropped sign agrees with the supplied \ensuremath{\ne} reading; the correction is mathematical and supported by scan 2, not a transcription emendation.

\paragraph{Confidence.} High for mathematics; decisive wording checked in READINGS.tsv.

\noindent Refined statement identifier: H0KB.7.correction (the corrected statement and its proof in this section).

\section*{H0KB.8}\label{statement:H0KB.8}
\paragraph{Locator.} 156-3/003.18;156-3/004.1

\paragraph{His words.} Pour chaque T modéré fini, il donne une relation d’équivalence R\_T dans C(T); pour chaque classe A, une partie Ȧ\ensuremath{\subset}T, appelée « frontière ».

\paragraph{How he left it.} stated

\paragraph{Precise statement (ours).} The data consist of arbitrary equivalence relations R\_T on the specified domains and designated subsets \ensuremath{\partial}\_T A\ensuremath{\subset}T. Calling their classes components makes no topological assertion.

\paragraph{Status.} True as written

\paragraph{Complete proof / counterexample application (ours).} Equivalence relations partition each domain; an arbitrary subset of T may be assigned to each class. The indiscrete and discrete systems of N give explicit realizations. Empty domains carry the unique empty relation and have no classes.

\paragraph{Correction and reconstruction (ours).} Our expansion of the displayed definitions and their hypotheses; no additional claim is attributed to the manuscript.

\paragraph{Independent second check.} The examples in N separately realize one-class and singleton-class choices, showing that no unseen topology is part of the definition.

\paragraph{Confidence.} High for mathematics; decisive wording checked in READINGS.tsv.

\section*{H0KB.9}\label{statement:H0KB.9}
\paragraph{Locator.} 156-3/004.5

\paragraph{His words.} « Je prévois » que c résulte de b, et que b résulte des R\_\ensuremath{\epsilon} pour \ensuremath{\epsilon}\ensuremath{\in}S\textsubscript{0}, après choix d’axiomes.

\paragraph{How he left it.} hoped

\paragraph{Precise statement (ours).} Find additional axioms under which frontiers are determined by partitions and all finite partitions by the two-point partitions. The unspecified additional axioms are part of the programme.

\paragraph{Status.} programme

\paragraph{Complete proof / counterexample application (ours).} With unspecified axioms this is not a fixed implication. The next two rows test the displayed provisional axioms; R and A prove the ordered-troncon version with its actual canonical cuts.

\paragraph{Correction and reconstruction (ours).} We do not refute the existence of a suitable eventual axiom system. We decide both supported provisional implications below.

\paragraph{Independent second check.} The wording explicitly precedes « Axiomes provisoires » and does not claim the forthcoming short list is final.

\paragraph{Confidence.} High for mathematics; decisive wording checked in READINGS.tsv.

\section*{H0KB.10}\label{statement:H0KB.10}
\paragraph{Locator.} 156-3/004.8

\paragraph{His words.} Première prévision : la donnée c, les frontières Ȧ, résulte déjà de b, les relations R\_T.

\paragraph{How he left it.} hoped

\paragraph{Precise statement (ours).} Precise provisional version: S\textsubscript{0} and all R\_T determine all \ensuremath{\partial}\_T A among systems satisfying Dec\textsubscript{1}, Dec\textsubscript{2} and the finite-maximal-cut axiom displayed on scans 5–6.

\paragraph{Status.} False

\paragraph{Complete proof / counterexample application (ours).} Lemma N gives two systems on Q with identical S\textsubscript{0} and identical indiscrete R\_T; their frontiers are respectively empty and T. It verifies every provisional clause. The ordered-troncon correction is true: canonical gaps have exactly their adjacent cut points as frontiers, by C, and R recovers all gap data from pair cuts.

\paragraph{Correction and reconstruction (ours).} This is our explicitly delimited version under the provisional axioms, not a reading that his unknown future axioms are refuted.

\paragraph{Independent second check.} Independent frontier check in N is the identity S\ensuremath{\cap}T=T in the nonempty-frontier system, versus the empty inclusion in the other.

\paragraph{Confidence.} High for mathematics; decisive wording checked in READINGS.tsv.

\noindent Refined statement identifier: H0KB.10.correction (the corrected statement and its proof in this section).

\section*{H0KB.11}\label{statement:H0KB.11}
\paragraph{Locator.} 156-3/004.10

\paragraph{His words.} Deuxième prévision : tous les R\_T résultent des seuls R\_\ensuremath{\epsilon} pour les paires \ensuremath{\epsilon}\ensuremath{\in}S\textsubscript{0}.

\paragraph{How he left it.} hoped

\paragraph{Precise statement (ours).} Precise provisional version: the two-point partitions determine every finite-cut partition under the same displayed provisional axioms.

\paragraph{Status.} False

\paragraph{Complete proof / counterexample application (ours).} Lemma N compares the all-indiscrete system with the system split by sign only at cuts of cardinality at least three. Both have empty frontiers and satisfy all provisional clauses; they agree on pair cuts but differ on triple cuts. The exact canonical ordered-troncon version is proved by R and A.

\paragraph{Correction and reconstruction (ours).} Additional recognition axioms are necessary; A supplies a complete finite scheme for the ordered subclass.

\paragraph{Independent second check.} The three cardinality cases in N check every possible refinement, including a two-point cut refined by a three-point one.

\paragraph{Confidence.} High for mathematics; decisive wording checked in READINGS.tsv.

\noindent Refined statement identifier: H0KB.11.correction (the corrected statement and its proof in this section).

\section*{H0KB.12}\label{statement:H0KB.12}
\paragraph{Locator.} 156-3/005.1

\paragraph{His words.} Si U est à la fois modérée finie et modérée cofinie, il prend S,T finis disjoints, S\ensuremath{\cup}T modéré, et écrit que « tout u\ensuremath{\in}L modéré » par rapport à S est dans T; il en déduit S\ensuremath{\cup}T maximal puis discret.

\paragraph{How he left it.} stated

\paragraph{Precise statement (ours).} The argument needs C(S)=T: every point OUTSIDE S compatible with S lies in T. Its literal clause omits outside S, although relative compatibility was declared reflexive. With that repair and the proposed maximal-finite-cut axiom, S\ensuremath{\cup}T=L and the structure is discrete.

\paragraph{Status.} Slip

\paragraph{Complete proof / counterexample application (ours).} For a two-point complete graph with S=\{s\}, T=\{t\}, the intended C(S)=T holds, while the literal clause fails at u=s. The corrected equivalence is direct: S\ensuremath{\cup}T tame gives T\ensuremath{\subseteq}C(S), and the repaired universal clause gives C(S)\ensuremath{\subseteq}T. Any point outside S\ensuremath{\cup}T extending that clique would belong to C(S)=T, a contradiction; thus S\ensuremath{\cup}T is maximal. Applying the explicitly proposed maximal-cut axiom gives the stated discrete conclusion. Conversely, a finite discrete structure has precisely these cofinite complements.

\paragraph{Correction and reconstruction (ours).} Our omitted-convention judgment is supported by the immediately preceding definition of C(S), which excludes S. The supplied phrase is retained, not silently rewritten. The discreteness conclusion uses the proposed axiom, not the graph data alone.

\paragraph{Independent second check.} The counterexample uses the explicit reflexive convention on scan 2. A second check compares each of the two inclusions in C(S)=T separately, then tests a hypothetical point extending S\ensuremath{\cup}T.

\paragraph{Confidence.} High for mathematics; decisive wording checked in READINGS.tsv.

\noindent Refined statement identifier: H0KB.12.correction (the corrected statement and its proof in this section).

\section*{H0KB.13}\label{statement:H0KB.13}
\paragraph{Locator.} 156-3/005.14

\paragraph{His words.} La structure « discrète » a S\textsubscript{0}=P\textsubscript{2}(L), C\_T=L\ensuremath{\setminus}T, R\_T discrète et Ȧ=\ensuremath{\varnothing}.

\paragraph{How he left it.} stated

\paragraph{Precise statement (ours).} These data define a coherent cut system for every set L. It satisfies Dec\textsubscript{1} exactly when L is finite, and always satisfies Dec\textsubscript{2}.

\paragraph{Status.} True as written

\paragraph{Complete proof / counterexample application (ours).} Lemma N verifies the equality partitions and empty frontiers. Applying Dec\textsubscript{1} to T=\ensuremath{\varnothing} proves necessity of finite L; a finite L makes every quotient finite. His data alone do not claim that an infinite discrete model obeys the later finiteness condition.

\paragraph{Correction and reconstruction (ours).} We distinguish the valid definition from the restriction imposed one page later.

\paragraph{Independent second check.} Count the singleton classes directly: |C(\ensuremath{\varnothing})/R\_\ensuremath{\varnothing}|=|L|.

\paragraph{Confidence.} High for mathematics; decisive wording checked in READINGS.tsv.

\section*{H0KB.14}\label{statement:H0KB.14}
\paragraph{Locator.} 156-3/005.23

\paragraph{His words.} « Axiome » : si T est une partie finie modérée maximale, alors T=L et L est discret.

\paragraph{How he left it.} stated

\paragraph{Precise statement (ours).} This is a proposed additional axiom: any maximal finite clique must equal L, with the discrete cut data. It is consistent, but is not a consequence of an arbitrary S\textsubscript{0}.

\paragraph{Status.} True as written

\paragraph{Complete proof / counterexample application (ours).} Finite discrete systems satisfy it; canonical G-orders have no maximal finite clique, since density inserts another point in any comparable gap and global bounds extend singleton or empty chains. Hence it holds vacuously there. The edge-plus-isolated-vertex graph in N shows why it must be imposed.

\paragraph{Correction and reconstruction (ours).} Status concerns a coherent proposed axiom, not a theorem falsely derived from the preceding data.

\paragraph{Independent second check.} Check the two cases independently: maximal finite clique in a finite complete graph is all of L; no finite chain in a nontrivial dense globally bidirected order is maximal.

\paragraph{Confidence.} High for mathematics; decisive wording checked in READINGS.tsv.

\section*{H0KB.15}\label{statement:H0KB.15}
\paragraph{Locator.} 156-3/006.1;156-3/008.4m

\paragraph{His words.} Les hypothèses de finitude sont « peut-être provisoires »; il note « refaire, sans se borner à une ligne ».

\paragraph{How he left it.} programme

\paragraph{Precise statement (ours).} Extend the line example to tronçons with possible endpoints and reassess finiteness. Supported versions are the G definition, endpoint cut counts and interval inheritance below.

\paragraph{Status.} programme

\paragraph{Complete proof / counterexample application (ours).} G and C give the complete supported extension. The words propose no further quantified finiteness theorem. The discrete restriction is already decided in H0KB.13.

\paragraph{Correction and reconstruction (ours).} Our expansion of the displayed definitions and their hypotheses; no additional claim is attributed to the manuscript.

\paragraph{Independent second check.} The later definition on 9 and the explicit endpoint cases on 11 supply the versions, so no unsupported higher-dimensional finiteness claim is invented.

\paragraph{Confidence.} High for mathematics; decisive wording checked in READINGS.tsv.

\section*{H0KB.16}\label{statement:H0KB.16}
\paragraph{Locator.} 156-3/006.4

\paragraph{His words.} Dec\textsubscript{1} : l’ensemble des composantes connexes de tout U=C(T) est fini.

\paragraph{How he left it.} stated

\paragraph{Precise statement (ours).} Dec\textsubscript{1} is finiteness of each quotient C(T)/R\_T. Canonical finite cuts in G satisfy it with at most |T|+1 classes.

\paragraph{Status.} True as written

\paragraph{Complete proof / counterexample application (ours).} Lemma C explicitly lists the finite fibers, including the empty-cut case; there are at most n+1. This proves consistency and the claimed ordered realization.

\paragraph{Correction and reconstruction (ours).} Our expansion of the displayed definitions and their hypotheses; no additional claim is attributed to the manuscript.

\paragraph{Independent second check.} The integer index x\ensuremath{\mapsto}\#\{t\ensuremath{\in}T:t<x\} has at most n+1 values.

\paragraph{Confidence.} High for mathematics; decisive wording checked in READINGS.tsv.

\section*{H0KB.17}\label{statement:H0KB.17}
\paragraph{Locator.} 156-3/006.7

\paragraph{His words.} Pour U=C(S), V=C(T), il écrit « U\ensuremath{\subset}V i.e. T\ensuremath{\subset}S ».

\paragraph{How he left it.} stated

\paragraph{Precise statement (ours).} The equivalence of these two conditions is proposed for arbitrary finite tame cuts. Only T\ensuremath{\subset}S\ensuremath{\Rightarrow}C(S)\ensuremath{\subset}C(T) holds without extra compatibility.

\paragraph{Status.} False

\paragraph{Complete proof / counterexample application (ours).} Use the two singleton twins of D: C(S)=C(T), while neither singleton contains the other. The implication from cut inclusion is H0KB.5. If S\ensuremath{\cup}T is tame, the converse is exactly the corrected witness argument in H0KB.6.

\paragraph{Correction and reconstruction (ours).} Replace « i.e. » by the explicit assumption T\ensuremath{\subset}S, or require S\ensuremath{\cup}T tame when asserting equivalence.

\paragraph{Independent second check.} Direct equality at every height checks the counterexample without appealing to an abstract noninjectivity assertion.

\paragraph{Confidence.} High for mathematics; decisive wording checked in READINGS.tsv.

\noindent Refined statement identifier: H0KB.17.correction (the corrected statement and its proof in this section).

\section*{H0KB.18}\label{statement:H0KB.18}
\paragraph{Locator.} 156-3/006.9

\paragraph{His words.} Chaque classe A de U est contenue dans une unique classe B de V; de plus Ȧ\ensuremath{\cap}T\ensuremath{\subset}Ḃ.

\paragraph{How he left it.} stated

\paragraph{Precise statement (ours).} For T\ensuremath{\subset}S, canonical G-cut partitions refine under restriction, and their adjacent-cut frontiers satisfy \ensuremath{\partial}\_S A\ensuremath{\cap}T\ensuremath{\subset}\ensuremath{\partial}\_T B.

\paragraph{Status.} True as written

\paragraph{Complete proof / counterexample application (ours).} Lemma C proves refinement by deleting some of the ordered cut points. A smaller gap lies between consecutive surviving old cuts, so lies in their unique gap. Any old cut retained as an adjacent endpoint of the smaller gap is one of those two bounding old cuts.

\paragraph{Correction and reconstruction (ours).} The genuine refinement assertion is separated from the false preceding « i.e. ».

\paragraph{Independent second check.} The integer-fiber description independently merges adjacent fibers when cuts are deleted.

\paragraph{Confidence.} High for mathematics; decisive wording checked in READINGS.tsv.

\section*{H0KB.19}\label{statement:H0KB.19}
\paragraph{Locator.} 156-3/007.1

\paragraph{His words.} Une « ligne ordonnée » est filtrante croissante et décroissante, sans plus grand ni plus petit élément, localement filtrante dans les deux sens et divisible.

\paragraph{How he left it.} stated

\paragraph{Precise statement (ours).} This is G without ambient extrema. It is a partial order; totality is not among its axioms.

\paragraph{Status.} True as written

\paragraph{Complete proof / counterexample application (ours).} Lemma G expands every axiom. D and the strict-coordinate plane P satisfy them, proving that the class is nonempty and genuinely includes non-total orders.

\paragraph{Correction and reconstruction (ours).} Our expansion of the displayed definitions and their hypotheses; no additional claim is attributed to the manuscript.

\paragraph{Independent second check.} In D both global and local requirements reduce to rational-height interpolation; in P they reduce to two independent strict coordinate inequalities.

\paragraph{Confidence.} High for mathematics; decisive wording checked in READINGS.tsv.

\section*{H0KB.20}\label{statement:H0KB.20}
\paragraph{Locator.} 156-3/007.8

\paragraph{His words.} S\textsubscript{0} est l’ensemble des paires totalement ordonnées; Mod\_f(L) est l’ensemble des drapeaux finis; un drapeau t\textsubscript{1}<\ensuremath{\cdots}<t\_n donne n+1 composantes avec les points frontière indiqués.

\paragraph{How he left it.} stated

\paragraph{Precise statement (ours).} For a nonempty finite chain on an endpoint-free G, C(T) has the n+1 strict gaps and their adjacent-cut frontiers; the empty chain has the one class L.

\paragraph{Status.} True as written

\paragraph{Complete proof / counterexample application (ours).} Lemma C proves that the strict gaps partition C(T). Density makes internal gaps nonempty; absence of extrema makes both exterior gaps nonempty. The boundary assignments are precisely the neighboring cut points.

\paragraph{Correction and reconstruction (ours).} Our expansion of the displayed definitions and their hypotheses; no additional claim is attributed to the manuscript.

\paragraph{Independent second check.} Count index fibers 0 through n, each realized; this independently gives n+1.

\paragraph{Confidence.} High for mathematics; decisive wording checked in READINGS.tsv.

\section*{H0KB.21}\label{statement:H0KB.21}
\paragraph{Locator.} 156-3/008.1

\paragraph{His words.} La structure ne change pas sous l’ordre opposé. Les tronçons fermés de la ligne sont [s,t], les demi-droites fermées et L; leurs réunions finies disjointes mutuellement modérées sont dites modérées.

\paragraph{How he left it.} stated

\paragraph{Precise statement (ours).} The cut and finite closed-piece structures are invariant under simultaneous order reversal. Closed pieces may be singleton intervals; compatible disjoint finite unions have intrinsic components and finite frontier.

\paragraph{Status.} True as written

\paragraph{Complete proof / counterexample application (ours).} Reversal exchanges left and right gaps, preserving their unordered partitions and endpoint sets. For closed pieces all finite endpoints belong to the pieces; the closed-part case of F proves uniqueness and the boundary-union rule without E.

\paragraph{Correction and reconstruction (ours).} Our expansion of the displayed definitions and their hypotheses; no additional claim is attributed to the manuscript.

\paragraph{Independent second check.} Reversing each endpoint inequality leaves the interval subset unchanged and reverses only its presentation order.

\paragraph{Confidence.} High for mathematics; decisive wording checked in READINGS.tsv.

\section*{H0KB.22}\label{statement:H0KB.22}
\paragraph{Locator.} 156-3/008.4m;156-3/018.1

\paragraph{His words.} Un sous-tronçon fermé d’un tronçon est encore un tronçon pour le biordre induit; il permet ensuite de retirer des extrémités.

\paragraph{How he left it.} stated

\paragraph{Precise statement (ours).} Every nondegenerate interval or nonempty ray of any endpoint type in G satisfies G in the induced order. Singleton pieces are separately called large or pseudo-tronçons.

\paragraph{Status.} True as written

\paragraph{Complete proof / counterexample application (ours).} The interval-inheritance proof in G supplies internal common bounds from the local axioms at omitted endpoints, uses included endpoints where available, and inherits density and the smaller local cones. Nonempty strict intervals are nondegenerate by repeated interpolation.

\paragraph{Correction and reconstruction (ours).} This assertion concerns order structure; it does not make these subsets topologically closed.

\paragraph{Independent second check.} The finite-cone proof in G independently checks that both of two selected points have a common bound still inside the interval.

\paragraph{Confidence.} High for mathematics; decisive wording checked in READINGS.tsv.

\section*{H0KB.23}\label{statement:H0KB.23}
\paragraph{Locator.} 156-3/009.4

\paragraph{His words.} Les axiomes a), a\ensuremath{\prime}), b), c), d) définissent le « tronçon ordonné » : filtrations locales, bornes communes, divisibilité et card L>1.

\paragraph{How he left it.} stated

\paragraph{Precise statement (ours).} G is exactly the partial-order axiom list in the common proofs: nontrivial, globally upper and lower directed, order-dense, with every nonempty strict upper cone downward directed and every strict lower cone upward directed.

\paragraph{Status.} True as written

\paragraph{Complete proof / counterexample application (ours).} Lemma G checks consistency, duality and the meaning of each filtering condition. The explanatory variable slips are recorded separately below; they do not replace the clearly stated filtering axiom.

\paragraph{Correction and reconstruction (ours).} Our spelling-out of filtering uses the period nonempty convention; no endpoint regularity E is added to his definition.

\paragraph{Independent second check.} Each of D, P and J independently realizes all five conditions.

\paragraph{Confidence.} High for mathematics; decisive wording checked in READINGS.tsv.

\section*{H0KB.24}\label{statement:H0KB.24}
\paragraph{Locator.} 156-3/009.6

\paragraph{His words.} Les lignes explicatives de a), a\ensuremath{\prime}), b) portent des a là où les autres termes et quantificateurs demandent x, notamment a>a et a,y\ensuremath{\leq}z<a.

\paragraph{How he left it.} stated

\paragraph{Precise statement (ours).} The literal repeated-variable glosses are inconsistent or vacuous and do not express the two filtering conditions. Replace the stray a labels in the x positions by x, retaining the quantified base point a.

\paragraph{Status.} Slip

\paragraph{Complete proof / counterexample application (ours).} The literal a>a antecedent can never hold, so cannot characterize downward directedness of a nonempty strict upper cone. The dual conclusion a\ensuremath{\leq}z<a is impossible: on Q, x=\ensuremath{-}2,y=\ensuremath{-}1,a=0 satisfy the intended premise but no z satisfies that literal conclusion. The corrected glosses are exactly a<x,y\ensuremath{\Rightarrow}\ensuremath{\exists}z(a<z\ensuremath{\leq}x,y) and x,y<a\ensuremath{\Rightarrow}\ensuremath{\exists}z(x,y\ensuremath{\leq}z<a). They are equivalent to the filtering phrases by definition.

\paragraph{Correction and reconstruction (ours).} These are explanatory pen slips; G is read from the explicitly named filtering properties, not weakened by vacuous antecedents.

\paragraph{Independent second check.} The repeated variables and at least the contradictory dual conclusion are visible in the crop. The formulas are also checked against the unambiguous global-bounds display and the later inheritance arguments.

\paragraph{Confidence.} High for the logical correction; medium for individual a/x penforms, with the decisive contradictory placement retained in READINGS.tsv.

\noindent Refined statement identifier: H0KB.24.correction (the corrected statement and its proof in this section).

\section*{H0KB.25}\label{statement:H0KB.25}
\paragraph{Locator.} 156-3/009.18;156-3/010.1

\paragraph{His words.} Segment : deux extrémités; intervalle : aucune; demi-intervalle : une seule. L’opposé est du même type. \ensuremath{\partial}L est l’ensemble du plus petit et du plus grand élément existants.

\paragraph{How he left it.} stated

\paragraph{Precise statement (ours).} Exactly one of the three endpoint types occurs; a half-interval has two orientations. \ensuremath{\partial}L has 0,1 or 2 elements, and reversal preserves it.

\paragraph{Status.} True as written

\paragraph{Complete proof / counterexample application (ours).} In a globally directed poset a maximal element is greatest, and a minimal element least, as proved in G. There is at most one of each; they are distinct because |L|>1. The four existence patterns group into precisely the three named types.

\paragraph{Correction and reconstruction (ours).} Our expansion of the displayed definitions and their hypotheses; no additional claim is attributed to the manuscript.

\paragraph{Independent second check.} Enumerate the two binary endpoint-existence choices; reversal exchanges the two one-end choices.

\paragraph{Confidence.} High for mathematics; decisive wording checked in READINGS.tsv.

\section*{H0KB.26}\label{statement:H0KB.26}
\paragraph{Locator.} 156-3/010.15;156-3/011.1

\paragraph{His words.} Pour un tronçon, S\textsubscript{0}=Drap\textsubscript{2}(L), Mod\_f(L)=Drap(L), et C(S) contient les t\ensuremath{\notin}S qui rendent S\ensuremath{\cup}\{t\} totalement ordonné.

\paragraph{How he left it.} stated

\paragraph{Precise statement (ours).} Finite tame sets are exactly finite chains, and C(S) consists of points outside S comparable to every point of S.

\paragraph{Status.} True as written

\paragraph{Complete proof / counterexample application (ours).} The clique condition for distinct comparability is precisely totality of the restricted partial order. The cross-pair computation of H0KB.4 then gives the displayed C(S).

\paragraph{Correction and reconstruction (ours).} Our expansion of the displayed definitions and their hypotheses; no additional claim is attributed to the manuscript.

\paragraph{Independent second check.} Every pair in S\ensuremath{\cup}\{t\} is tested separately; no transitive orientation needs to be chosen beyond the given order.

\paragraph{Confidence.} High for mathematics; decisive wording checked in READINGS.tsv.

\section*{H0KB.27}\label{statement:H0KB.27}
\paragraph{Locator.} 156-3/011.6

\paragraph{His words.} La liste des composantes porte des crochets fermés [s\textsubscript{1},s\textsubscript{2}], …, [s\_\{n\ensuremath{-}1\},s\_n] au milieu des deux cônes stricts.

\paragraph{How he left it.} stated

\paragraph{Precise statement (ours).} The literal closed internal intervals contain cut points and cannot be classes in C(S). They must be the strict gaps (s\_i,s\_\{i+1\}).

\paragraph{Status.} Slip

\paragraph{Complete proof / counterexample application (ours).} On Q with S=\{0,1\}, the printed [0,1] contains 0,1, excluded from C(S). Replacing it by (0,1) gives the partition proved in C. The preceding list on scan 7 and the subsequent three-point calculation on 16 both use strict gaps.

\paragraph{Correction and reconstruction (ours).} Our slip judgment is supported by his own domain C(S) and his correct lists on 7 and 16.

\paragraph{Independent second check.} Cropped brackets are visibly closed; they are not silently transcribed as open. The integer-fiber proof verifies the corrected partition.

\paragraph{Confidence.} High for mathematics; decisive wording checked in READINGS.tsv.

\noindent Refined statement identifier: H0KB.27.correction (the corrected statement and its proof in this section).

\section*{H0KB.28}\label{statement:H0KB.28}
\paragraph{Locator.} 156-3/011.13

\paragraph{His words.} Il écrit « les seules non vides éventuelles » pour L\_<s\textsubscript{1} si s\textsubscript{1} est le plus petit élément, et L\_>s\_n si s\_n est le plus grand.

\paragraph{How he left it.} stated

\paragraph{Precise statement (ours).} Only the exterior gaps can be empty; the left is empty exactly at an ambient minimum and the right exactly at an ambient maximum.

\paragraph{Status.} Slip

\paragraph{Complete proof / counterexample application (ours).} For L=[0,1] and s\textsubscript{1}=0 the displayed left tail is empty, directly contradicting the word non. Dropping non yields the claim in C. With n=0 there is one component L; with n=1 there are one or two nonempty tails; with n\ensuremath{\geq}2 there are n\ensuremath{-}1 internal gaps plus the nonempty tails.

\paragraph{Correction and reconstruction (ours).} Keep the visible words in the record; delete only « non » in the corrected mathematical statement.

\paragraph{Independent second check.} The endpoint-existence characterization in G and density independently rule out every other empty gap.

\paragraph{Confidence.} High for mathematics; decisive wording checked in READINGS.tsv.

\noindent Refined statement identifier: H0KB.28.correction (the corrected statement and its proof in this section).

\section*{H0KB.29}\label{statement:H0KB.29}
\paragraph{Locator.} 156-3/011.18

\paragraph{His words.} Il dit récupérer le biordre du tronçon, donc a), b), c), à partir des paires modérées et des R\_T pour T\ensuremath{\in}S\textsubscript{0}.

\paragraph{How he left it.} proved on the page

\paragraph{Precise statement (ours).} For every G of any cardinality, distinct comparability and the canonical two-point cut equivalence relations determine its order up to simultaneous reversal, hence all canonical finite cuts and frontiers.

\paragraph{Status.} True as written

\paragraph{Complete proof / counterexample application (ours).} Apply R2 to recover strict betweenness from the disjointness of the indicated two cut classes. Density supplies the nonempty alternative intersections. Recover reflexive comparability and N, select one oriented pair a<b, and apply R1. C then reconstructs every finite cut and its adjacent frontiers. All hypotheses are exactly G.

\paragraph{Correction and reconstruction (ours).} The explicit anchored formula and all-cardinality proof are ours; the page gives the same strategy and its local lemma.

\paragraph{Independent second check.} The independent two-orders orientation argument in R verifies uniqueness without reusing the reconstruction formula.

\paragraph{Confidence.} High for mathematics; decisive wording checked in READINGS.tsv.

\section*{H0KB.30}\label{statement:H0KB.30}
\paragraph{Locator.} 156-3/012.2

\paragraph{His words.} Le biordre est connu par R(a,b,c), « b est strictement entre a et c »; a et b distincts sont comparables ssi un c est entre eux.

\paragraph{How he left it.} proved on the page

\paragraph{Precise statement (ours).} In G, B(x,y,z)\ensuremath{\Longleftrightarrow}x<y<z or z<y<x determines distinct comparability by \ensuremath{\exists}y B(x,y,z), and determines the order up to reversal.

\paragraph{Status.} True as written

\paragraph{Complete proof / counterexample application (ours).} The forward implication to comparability is transitivity; the reverse is density. With this derived C, R1 reconstructs the oriented order. Global bounds and |L|>1 ensure an anchor pair exists.

\paragraph{Correction and reconstruction (ours).} Our expansion of the displayed definitions and their hypotheses; no additional claim is attributed to the manuscript.

\paragraph{Independent second check.} Repeated interpolation between an existing comparable pair also proves the marginal observation that finite chains of arbitrary length exist. The independent proof in R verifies uniqueness.

\paragraph{Confidence.} High for mathematics; decisive wording checked in READINGS.tsv.

\section*{H0KB.31}\label{statement:H0KB.31}
\paragraph{Locator.} 156-3/012.6m

\paragraph{His words.} Le NB remplace l’entre-deux non strict par l’entre-deux strict ou « \{a,b,c\} est un drapeau de cardinal \ensuremath{\leq}2 ».

\paragraph{How he left it.} stated

\paragraph{Precise statement (ours).} N(a,b,c) means a\ensuremath{\leq}b\ensuremath{\leq}c or c\ensuremath{\leq}b\ensuremath{\leq}a. The shorthand using only a two-element chain is too broad when a=c\ensuremath{\ne}b.

\paragraph{Status.} Slip

\paragraph{Complete proof / counterexample application (ours).} On Q, (a,b,c)=(0,1,0) satisfies the printed cardinality/chain condition but not N. The exact repair is B(a,b,c) or ((a=b or b=c) and C(a,c)). All equality patterns are covered: three equal, equal first two, equal last two; equal outer arguments require all three equal.

\paragraph{Correction and reconstruction (ours).} An evident shorthand omission, judged against his explicit inequalities a\ensuremath{\leq}b\ensuremath{\leq}c or the reverse; the crop confirms the supplied \ensuremath{\leq}2 text.

\paragraph{Independent second check.} Enumerate these four repetition patterns separately. This is independent of the proof of order reconstruction.

\paragraph{Confidence.} High for mathematics; decisive wording checked in READINGS.tsv.

\noindent Refined statement identifier: H0KB.31.correction (the corrected statement and its proof in this section).

\section*{H0KB.32}\label{statement:H0KB.32}
\paragraph{Locator.} 156-3/013.1

\paragraph{His words.} Choisir une origine a dans \{a,b\}; les cônes L\_\ensuremath{\leq}a, L\_\ensuremath{\geq}b et l’intervalle ]a,b[ sont déterminés par R. Pour u\ensuremath{\leq}v, il faut et suffit des bornes x\ensuremath{\leq}a, y\ensuremath{\geq}b avec x\ensuremath{\leq}u\ensuremath{\leq}v\ensuremath{\leq}y.

\paragraph{How he left it.} proved on the page

\paragraph{Precise statement (ours).} For fixed a<b, the D,U formulas and R1 of the common proof define the oriented order from B. Either of the two additional nonstrict betweenness conditions in his display suffices.

\paragraph{Status.} True as written

\paragraph{Complete proof / counterexample application (ours).} Global directedness supplies x\ensuremath{\leq}a,u and y\ensuremath{\geq}b,v when u\ensuremath{\leq}v. Conversely the two endpoint-between conditions place u,v in [x,y], and N(x,u,v) forces u\ensuremath{\leq}v; its dual N(u,v,y) also suffices. This is proved in R without assuming u,v already comparable.

\paragraph{Correction and reconstruction (ours).} Our expansion of the displayed definitions and their hypotheses; no additional claim is attributed to the manuscript.

\paragraph{Independent second check.} Exhaustive bounded-poset enumeration in check\_finite.py verifies the relaxed formula on 408 finite posets and 6057 anchors; the independent mathematical proof in R handles all cardinalities.

\paragraph{Confidence.} High for mathematics; decisive wording checked in READINGS.tsv.

\section*{H0KB.33}\label{statement:H0KB.33}
\paragraph{Locator.} 156-3/014.2

\paragraph{His words.} Le lemme nomme x,y,u,v « trois éléments »; a) u,v sont entre x,y; b) u entre x,v, ou c) v entre u,y; d) x<y; alors x\ensuremath{\leq}u\ensuremath{\leq}v\ensuremath{\leq}y.

\paragraph{How he left it.} proved on the page

\paragraph{Precise statement (ours).} For four named, not necessarily distinct elements, the displayed betweenness hypotheses imply the displayed chain; either b) or c) suffices.

\paragraph{Status.} Slip

\paragraph{Complete proof / counterexample application (ours).} The literal word trois mismatches four names and cannot be a general restriction to three distinct points: 0<1<2<3 satisfies the intended lemma. Replace trois by quatre (allowing coincidences). R proves the implication, including the equality cases: the reverse orientation in N(x,u,v) can occur only when u=v=x.

\paragraph{Correction and reconstruction (ours).} Only the count word is corrected; the lemma’s mathematical inequality is true as displayed.

\paragraph{Independent second check.} The program checks all 38169 satisfied four-point antecedents in its bounded-poset sample. The direct inequality proof is independent and exhaustive for arbitrary posets.

\paragraph{Confidence.} High for mathematics; decisive wording checked in READINGS.tsv.

\noindent Refined statement identifier: H0KB.33.correction (the corrected statement and its proof in this section).

\section*{H0KB.34}\label{statement:H0KB.34}
\paragraph{Locator.} 156-3/014.18;156-3/015.1

\paragraph{His words.} Il note que la divisibilité a seulement servi à déduire les paires de l’entre-deux; si l’on donne les deux relations, la filtration croissante et décroissante suffit.

\paragraph{How he left it.} proved on the page

\paragraph{Precise statement (ours).} A nontrivial globally upper and lower directed poset is reconstructible up to reversal from comparability and strict betweenness, with no density or local cone filtering assumption.

\paragraph{Status.} True as written

\paragraph{Complete proof / counterexample application (ours).} Use the given C to define N and R1. Its proof uses only common upper/lower bounds and an oriented comparable pair. A comparable distinct pair exists from common bounds of two distinct points.

\paragraph{Correction and reconstruction (ours).} Our expansion of the displayed definitions and their hypotheses; no additional claim is attributed to the manuscript.

\paragraph{Independent second check.} The finite bounded-poset enumeration targets precisely this weakened statement; its samples are not misleadingly called dense tronçons.

\paragraph{Confidence.} High for mathematics; decisive wording checked in READINGS.tsv.

\section*{H0KB.35}\label{statement:H0KB.35}
\paragraph{Locator.} 156-3/015.7;156-3/016.1

\paragraph{His words.} Pour a<b<c, les classes X\_\{ab\},X\_\{bc\},X\_\{ca\} sont L\_>b,L\_<b,]a,c[; leurs intersections sont \ensuremath{\varnothing},]b,c[,]a,b[. « On gagne ! »

\paragraph{How he left it.} proved on the page

\paragraph{Precise statement (ours).} For a distinct three-point chain, b is the middle point exactly when the class of c in C(\{a,b\}) and the class of a in C(\{b,c\}) are disjoint. The two other intersections are the nonempty adjacent gaps.

\paragraph{Status.} True as written

\paragraph{Complete proof / counterexample application (ours).} Intersect the three displayed interval/ray sets directly. The two nonempty intersections use density. For each of the other choices of middle point relabel the triple, proving R2 for all six oriented orders.

\paragraph{Correction and reconstruction (ours).} Our expansion of the displayed definitions and their hypotheses; no additional claim is attributed to the manuscript.

\paragraph{Independent second check.} The independent check enumerates the three possible middle labels rather than inferring them from a diagram. Reversal leaves all class intersections unchanged.

\paragraph{Confidence.} High for mathematics; decisive wording checked in READINGS.tsv.

\section*{H0KB.36}\label{statement:H0KB.36}
\paragraph{Locator.} 156-3/017.1

\paragraph{His words.} Il cherche la plus grande sous-catégorie pleine des ordres bidirigés non triviaux, divisibles, stable par les cônes stricts non vides et leurs duaux; elle existe, est autoduale, et ses objets sont les tronçons.

\paragraph{How he left it.} stated

\paragraph{Precise statement (ours).} Among full classes of posets with global bounds, nontriviality, density and closure under nonempty strict cones, the largest class is exactly G.

\paragraph{Status.} True as written

\paragraph{Complete proof / counterexample application (ours).} Every object in such a class has downward directed strict upper cones because those cones themselves are globally bidirected objects; the dual supplies the other local axiom. Thus the class is contained in G. Conversely G inherits itself on strict cones by lemma G, so it satisfies all requirements and is the largest class. Reversal is immediate.

\paragraph{Correction and reconstruction (ours).} Our expansion of the displayed definitions and their hypotheses; no additional claim is attributed to the manuscript.

\paragraph{Independent second check.} This is a two-containment proof; it avoids an unsupported size argument about unions of categories. Fullness merely takes all monotone maps between the objects.

\paragraph{Confidence.} High for mathematics; decisive wording checked in READINGS.tsv.

\section*{H0KB.37}\label{statement:H0KB.37}
\paragraph{Locator.} 156-3/017.17;156-3/041.10

\paragraph{His words.} « Il faudrait expliciter » les conditions sur R(a,b,c), et sur S\textsubscript{0} avec les R\_\ensuremath{\epsilon}; à la fin il renonce à expliciter tout l’exercice.

\paragraph{How he left it.} programme

\paragraph{Precise statement (ours).} Find intrinsic necessary and sufficient axioms for an unoriented tronçon’s betweenness or pair-cut data. The next two rows are the precise existence versions; no independence or minimality requirement is stated.

\paragraph{Status.} programme

\paragraph{Complete proof / counterexample application (ours).} Lemma A supplies two explicit finite first-order recognition schemes. The programme row records what the text leaves unwritten; the two supported versions are fully decided below.

\paragraph{Correction and reconstruction (ours).} Our expansion of the displayed definitions and their hypotheses; no additional claim is attributed to the manuscript.

\paragraph{Independent second check.} The concluding wording on 41 changes neither the data nor the demand for recognition; it questions the value of writing the full exercise, not the truth of the reconstruction theorem.

\paragraph{Confidence.} High for mathematics; decisive wording checked in READINGS.tsv.

\section*{H0KB.38}\label{statement:H0KB.38}
\paragraph{Locator.} 156-3/017.17

\paragraph{His words.} Version précise de la demande sur R(a,b,c) : donner les conditions qui définissent un biordre de tronçon.

\paragraph{How he left it.} asked

\paragraph{Precise statement (ours).} There exists an explicit finite first-order sentence in B whose models are exactly the strict betweenness structures of G-orders, in every cardinality.

\paragraph{Status.} True as written

\paragraph{Complete proof / counterexample application (ours).} Use the full scheme A: define C,N,D,U,Q from B, existentially choose an anchor, impose every G axiom on the definable Q, and impose B exactly equal to strict Q-betweenness. Necessity is R1; sufficiency takes that Q as the constructed order. Every quantifier is over points, not over relations.

\paragraph{Correction and reconstruction (ours).} This is our precise existence reading of his requested axiomatization and our solution, not an axiom list attributed to him.

\paragraph{Independent second check.} The independent model-to-order and order-to-model substitutions in A check both directions and every axiom.

\paragraph{Confidence.} High for mathematics; decisive wording checked in READINGS.tsv.

\section*{H0KB.39}\label{statement:H0KB.39}
\paragraph{Locator.} 156-3/017.19;156-3/041.4

\paragraph{His words.} Version précise de la demande sur S\textsubscript{0} et les R\_\ensuremath{\epsilon} dans C(\ensuremath{\epsilon}).

\paragraph{How he left it.} asked

\paragraph{Precise statement (ours).} There exists an explicit finite first-order recognition scheme for a symmetric pair predicate S and a four-place pair-cut relation R(a,b;x,y), recognizing exactly canonical G pair cuts.

\paragraph{Status.} True as written

\paragraph{Complete proof / counterexample application (ours).} A defines B by absence of a common point in the two relevant classes, applies the anchored G recognition clauses, and forces S and R to equal exactly comparability and the three gap equivalence relation of the constructed order. This proves necessity and sufficiency and yields all finite cuts by C.

\paragraph{Correction and reconstruction (ours).} The recognition scheme is ours. It explicitly identifies the supplied predicates; it does not merely produce some unrelated order.

\paragraph{Independent second check.} Direct substitution checks the domain restrictions, symmetry in the cut pair, equivalence on that domain, and the three-gap disjunction.

\paragraph{Confidence.} High for mathematics; decisive wording checked in READINGS.tsv.

\section*{H0KB.40}\label{statement:H0KB.40}
\paragraph{Locator.} 156-3/018.1

\paragraph{His words.} Il donne \ensuremath{\partial}I=\{a,b\}, \{a\}, \{b\}, \ensuremath{\varnothing} pour les sous-tronçons fermés; la marge précise « frontière relative », non bord absolu.

\paragraph{How he left it.} stated

\paragraph{Precise statement (ours).} Relative frontier omits an ambient extremum included in a non-singleton component, and retains an isolated singleton extremum. The first displayed formulas omit this exception.

\paragraph{Status.} Slip

\paragraph{Complete proof / counterexample application (ours).} For L=[0,1] and I=L=[0,1], the first formula gives \{0,1\}, while the separate L case gives empty. The corrected rule removes ambient extrema in non-singleton components, exactly as written on scan 30; for a singleton it retains its point. The interval cases verify it directly.

\paragraph{Correction and reconstruction (ours).} Our slip judgment uses the explicit relative-frontier margin and his later precise exception on 156-3/030. This is a combinatorial frontier, not yet a topological theorem.

\paragraph{Independent second check.} Compare I=[0,1/2], I=L, and I=\{0\}. These three cases independently force both the omission and the singleton exception.

\paragraph{Confidence.} High for mathematics; decisive wording checked in READINGS.tsv.

\noindent Refined statement identifier: H0KB.40.correction (the corrected statement and its proof in this section).

\section*{H0KB.41}\label{statement:H0KB.41}
\paragraph{Locator.} 156-3/018.7

\paragraph{His words.} Deux sous-tronçons fermés larges sont modérés si leurs frontières le sont; leur intersection est vide ou un tel sous-tronçon, et sa frontière est dans la réunion des frontières.

\paragraph{How he left it.} stated

\paragraph{Precise statement (ours).} In G, compatible closed intervals/rays/singletons have empty or closed-piece intersection, with relative frontier contained in the union of the input frontiers.

\paragraph{Status.} True as written

\paragraph{Complete proof / counterexample application (ours).} Their finite endpoints form a chain; take the maximum lower bound and minimum upper bound. If these coincide the intersection is that singleton; if they cross it is empty; otherwise it is the indicated interval/ray. An internal new endpoint is an old one. At an ambient extremum a singleton intersection can occur only if one input already has that singleton frontier.

\paragraph{Correction and reconstruction (ours).} Our expansion of the displayed definitions and their hypotheses; no additional claim is attributed to the manuscript.

\paragraph{Independent second check.} The direct endpoint intersection proof is independent of the finite-mask proof in F. No E is used because all finite endpoints are retained.

\paragraph{Confidence.} High for mathematics; decisive wording checked in READINGS.tsv.

\section*{H0KB.42}\label{statement:H0KB.42}
\paragraph{Locator.} 156-3/018.17

\paragraph{His words.} Après définition de la topologie par les intervalles, il affirme que, pour un pseudo-tronçon fermé, sa frontière combinatoire est la frontière topologique.

\paragraph{How he left it.} stated

\paragraph{Precise statement (ours).} For all G, the displayed relative endpoints of a closed piece are its frontier in the canonical topology.

\paragraph{Status.} False

\paragraph{Complete proof / counterexample application (ours).} In D, the singleton \{a\_0\} has combinatorial frontier \{a\_0\}, but its topological closure and frontier are \{a\_0,a\_1\}; its interior is empty. Thus the two frontiers differ. Under E, principal closed pieces are canonically closed and their only boundary points are their relative endpoints, by the neighborhood and closure proofs E and F.

\paragraph{Correction and reconstruction (ours).} Add E for a theorem about all closed pieces. The author himself later separates the two notions on scans 25 and 30.

\paragraph{Independent second check.} Compute the pulled-back rational topology explicitly. The extra twin is in every neighborhood of a\_0, while no nonempty open set lies in the singleton.

\paragraph{Confidence.} High for mathematics; decisive wording checked in READINGS.tsv.

\noindent Refined statement identifier: H0KB.42.correction (the corrected statement and its proof in this section).

\section*{H0KB.43}\label{statement:H0KB.43}
\paragraph{Locator.} 156-3/019.1

\paragraph{His words.} Th : une réunion finie de pseudo-tronçons fermés disjoints et mutuellement modérés possède une famille I unique; ses T\_i sont les composantes.

\paragraph{How he left it.} stated

\paragraph{Precise statement (ours).} A closed tame subset of G has a unique decomposition into nonempty pairwise compatible disjoint closed pieces.

\paragraph{Status.} True as written

\paragraph{Complete proof / counterexample application (ours).} Compatible disjoint closed pieces are strictly separated: their facing endpoints are distinct and density supplies a missing point between them. The convex-separator proof in F then proves uniqueness by comparing two decompositions component by component. Empty union is uniquely empty.

\paragraph{Correction and reconstruction (ours).} Our expansion of the displayed definitions and their hypotheses; no additional claim is attributed to the manuscript.

\paragraph{Independent second check.} The endpoint sweep gives the same components as maximal runs not separated by a strict missing gap. This second proof uses actual retained endpoints and needs no E.

\paragraph{Confidence.} High for mathematics; decisive wording checked in READINGS.tsv.

\section*{H0KB.44}\label{statement:H0KB.44}
\paragraph{Locator.} 156-3/019.15

\paragraph{His words.} Ẋ est la réunion des Ṫ\_i; Ẋ=\ensuremath{\varnothing} ssi X=L ou X=\ensuremath{\varnothing}.

\paragraph{How he left it.} stated

\paragraph{Precise statement (ours).} With the relative-endpoint convention, a closed tame union has frontier the finite disjoint union of component frontiers. It is empty exactly for the empty or full part.

\paragraph{Status.} True as written

\paragraph{Complete proof / counterexample application (ours).} Uniqueness makes the union intrinsic. Different disjoint closed components have disjoint endpoint sets. Any nonempty proper component either has an internal endpoint or is an isolated singleton; that point contributes to the frontier. If none does, the sole component is L. This proves both directions.

\paragraph{Correction and reconstruction (ours).} Our expansion of the displayed definitions and their hypotheses; no additional claim is attributed to the manuscript.

\paragraph{Independent second check.} Check rays, bounded intervals, singleton ambient endpoints and the full part separately. The correction of the earlier convention is already isolated in H0KB.40.

\paragraph{Confidence.} High for mathematics; decisive wording checked in READINGS.tsv.

\section*{H0KB.45}\label{statement:H0KB.45}
\paragraph{Locator.} 156-3/019.20

\paragraph{His words.} Deux parties fermées modérées sont en position relative modérée si leurs frontières le sont; X\ensuremath{\subset}Y implique cette position.

\paragraph{How he left it.} stated

\paragraph{Precise statement (ours).} For closed tame subsets of G, compatibility is intrinsic and inclusion implies compatibility.

\paragraph{Status.} True as written

\paragraph{Complete proof / counterexample application (ours).} Each X component lies in one Y component by the convex-separator argument. Since endpoints of closed pieces belong to the pieces, its finite endpoints lie between the containing component’s endpoints and compare with all other component endpoints through their strict order. Ambient extrema compare with every point.

\paragraph{Correction and reconstruction (ours).} Our expansion of the displayed definitions and their hypotheses; no additional claim is attributed to the manuscript.

\paragraph{Independent second check.} Direct endpoint comparison gives the same conclusion as a common finite flag; this is the closed-piece case of F and uses no E.

\paragraph{Confidence.} High for mathematics; decisive wording checked in READINGS.tsv.

\section*{H0KB.46}\label{statement:H0KB.46}
\paragraph{Locator.} 156-3/020.1

\paragraph{His words.} Pour X,Y fermés modérés et mutuellement modérés, X\ensuremath{\cap}Y est fermé modéré, X\ensuremath{\vee}Y existe, et les frontières des deux sont contenues dans Ẋ\ensuremath{\cup}Ẏ. Il étend ceci à une famille finie.

\paragraph{How he left it.} sketched

\paragraph{Precise statement (ours).} Compatible closed tame parts have global meet and join in the poset of closed tame parts; their frontiers use only the original endpoints. The operations preserve compatibility with any common compatible part.

\paragraph{Status.} True as written

\paragraph{Complete proof / counterexample application (ours).} Intersect component pairs using the previous interval rule. For the join, sweep the common endpoint chain and merge all touching or overlapping closed pieces; any closed tame upper bound must contain each merged interval by convexity, hence this is least. New endpoints are selected from old endpoints, proving both frontier and compatibility assertions. Iterate for a finite family.

\paragraph{Correction and reconstruction (ours).} At this point « parties modérées » denotes the closed-parts construction just defined; no later open-parts hull axiom is imported.

\paragraph{Independent second check.} The independent finite-flag proof in F gives the same operations and establishes leastness among all closed tame upper bounds.

\paragraph{Confidence.} High for mathematics; decisive wording checked in READINGS.tsv.

\section*{H0KB.47}\label{statement:H0KB.47}
\paragraph{Locator.} 156-3/020.12m

\paragraph{His words.} « Attention » : une partie quelconque A n’a pas toujours de plus petit fermé modéré la contenant, déjà dans [0,1], avec une suite 1/n.

\paragraph{How he left it.} stated

\paragraph{Precise statement (ours).} There need not be a least finite union of closed intervals containing an arbitrary subset, even on [0,1].

\paragraph{Status.} True as written

\paragraph{Complete proof / counterexample application (ours).} Lemma H proves this with A=\{0\}\ensuremath{\cup}\{1/n:n\ensuremath{\geq}1\}. Every finite closed-interval cover contains [0,\ensuremath{\epsilon}]; an irrational point inside it can be omitted by another finite closed tame cover of A. Thus no cover is least.

\paragraph{Correction and reconstruction (ours).} The exact sequence with 0 included is our unambiguous realization of his example; omitting 0 still forces it in every closed cover.

\paragraph{Independent second check.} The competing cover is explicitly [0,1/N] plus finitely many remaining sequence points; it verifies the direction of the leastness obstruction.

\paragraph{Confidence.} High for mathematics; decisive wording checked in READINGS.tsv.

\section*{H0KB.48}\label{statement:H0KB.48}
\paragraph{Locator.} 156-3/020.18;156-3/021.1

\paragraph{His words.} [a,b]\ensuremath{\cup}\_mod[b,c]=[a,c], mais la réunion ordinaire peut différer. L est totalement ordonné ssi L\_\ensuremath{\leq}b\ensuremath{\cup}L\_\ensuremath{\geq}b=L pour tout b.

\paragraph{How he left it.} stated

\paragraph{Precise statement (ours).} For a\ensuremath{\leq}b\ensuremath{\leq}c the closed tame join fills [a,c]; ordinary union equals it precisely when every point of [a,c] is comparable with b. The analogous all-rays property for every b characterizes total order.

\paragraph{Status.} True as written

\paragraph{Complete proof / counterexample application (ours).} Any tame upper bound containing both touching intervals has one component containing them and, by convexity, [a,c]. The ordinary union is exactly the points of [a,c] comparable with b. A twin of b in D lies in [a,c] for outer heights a<b<c but in neither half. The global ray union is exactly the comparability set of b; requiring it to be L for every b is totality.

\paragraph{Correction and reconstruction (ours).} Our expansion of the displayed definitions and their hypotheses; no additional claim is attributed to the manuscript.

\paragraph{Independent second check.} Evaluate membership on an arbitrary point rather than on endpoint masks. This separately checks the join/union distinction and both directions of the totality criterion.

\paragraph{Confidence.} High for mathematics; decisive wording checked in READINGS.tsv.

\section*{H0KB.49}\label{statement:H0KB.49}
\paragraph{Locator.} 156-3/021.8

\paragraph{His words.} Il barre la première affirmation : il existe un plus petit fermé modéré contenant X\textbackslash{}Y, pour Y\ensuremath{\subset}X fermés modérés.

\paragraph{How he left it.} cancelled

\paragraph{Precise statement (ours).} The cancelled universal least-hull assertion for the ordinary difference of closed tame parts is false, even for some G satisfying E.

\paragraph{Status.} False

\paragraph{Complete proof / counterexample application (ours).} Lemma H2 gives the regular split-line example X=L, Y=\ensuremath{\leq}a on the first branch. Its ordinary difference contains the entire second branch. Every left-tail point p supplies a closed-ray cover \ensuremath{\geq}p; their intersection excludes the left tail, and no finite closed tame union inside it can cover the second branch all the way down to its unattained lower cut. Thus no least hull exists. All G and E axioms were verified in O.

\paragraph{Correction and reconstruction (ours).} The statement is true on total G-orders, where ordinary and compatible difference agree; for general endpoint-regular G the correct replacement is C\_mod(Y,X), as in the next row.

\paragraph{Independent second check.} The compatible difference in the same example has the hull \ensuremath{\geq}a and omits the second branch. This independently checks that the author’s replacement really changes the set, and cannot prove the cancelled claim.

\paragraph{Confidence.} High for mathematics; decisive wording checked in READINGS.tsv.

\noindent Refined statement identifier: H0KB.49.correction (the corrected statement and its proof in this section).

\section*{H0KB.50}\label{statement:H0KB.50}
\paragraph{Locator.} 156-3/021.10

\paragraph{His words.} Il remplace X\ensuremath{\setminus}Y par C\_mod(Y,X), les points de X hors Y en position modérée avec Y; il affirme l’existence de son plus petit fermé modéré Z\ensuremath{\subset}X.

\paragraph{How he left it.} stated

\paragraph{Precise statement (ours).} For closed tame Y\ensuremath{\subset}X in arbitrary G, C\_mod(Y,X) is asserted to have a least closed tame hull Z, compatible with Y and containing every compatible closed part of X disjoint from Y.

\paragraph{Status.} False

\paragraph{Complete proof / counterexample application (ours).} In D take X=L and Y=L\_\{\ensuremath{\leq}a\_0\}. Its relative tame complement is the strict positive ray. By the dual of the complete D argument it has no least closed tame hull. Under E, the common-flag calculus makes C\_mod(Y,X)=X\ensuremath{\wedge}Y\ensuremath{\prime} tame and defines Z as its closure. Since X is closed, Z\ensuremath{\subset}X; closure preserves compatibility; every competing disjoint compatible part is already contained in C\_mod(Y,X), hence in Z.

\paragraph{Correction and reconstruction (ours).} E is the exact all-rays additional hypothesis. This is not the cancelled ordinary-difference claim, whose separate row is audited independently.

\paragraph{Independent second check.} The independent topological characterization E gives Z as the canonical closure under E, and immediately proves leastness and Z\ensuremath{\subset}X.

\paragraph{Confidence.} High for mathematics; decisive wording checked in READINGS.tsv.

\noindent Refined statement identifier: H0KB.50.correction (the corrected statement and its proof in this section).

\section*{H0KB.51}\label{statement:H0KB.51}
\paragraph{Locator.} 156-3/022.1

\paragraph{His words.} F, frontière relative de Y dans X, est fini, modéré, contenu dans Z, sans point isolé de Z; C\_mod(Y,X)=C\_mod(F,Z).

\paragraph{How he left it.} stated

\paragraph{Precise statement (ours).} The preceding relative-difference construction has the finite-frontier description stated, wherever its required hull calculus is valid; the universal construction on all G is not valid.

\paragraph{Status.} False

\paragraph{Complete proof / counterexample application (ours).} The same X,Y in D as H0KB.50 have no Z, so the claimed universal construction cannot be made. Under E work in a common endpoint chain: the relative complement consists of the gaps in X outside Y; its closure fills their adjacent excluded endpoints. These are exactly F, each adjacent to a nonempty gap and hence nonisolated by density. Removing F compatibly from Z recovers the original gaps.

\paragraph{Correction and reconstruction (ours).} Complete corrected version: G+E, Y\ensuremath{\subset}X closed tame, Z=closure of X\ensuremath{\wedge}Y\ensuremath{\prime}, and F=relative frontier of Y in X.

\paragraph{Independent second check.} Inspect the two facing endpoint flags for each gap. No isolated component can consist of a newly added frontier point, since every such point bounds a nonempty interval.

\paragraph{Confidence.} High for mathematics; decisive wording checked in READINGS.tsv.

\noindent Refined statement identifier: H0KB.51.correction (the corrected statement and its proof in this section).

\section*{H0KB.52}\label{statement:H0KB.52}
\paragraph{Locator.} 156-3/022.8

\paragraph{His words.} (9 juin) La base d’ouverts est ]a,b[, L\_<a, L\_>b; certaines demi-droites peuvent être omises quand l’extrémité correspondante n’existe pas.

\paragraph{How he left it.} stated

\paragraph{Precise statement (ours).} These strict intervals and rays form a topology basis in every G. Left rays are redundant if there is no minimum, and right rays if there is no maximum.

\paragraph{Status.} True as written

\paragraph{Complete proof / counterexample application (ours).} Lemma P proves coverage and the intersection refinement by finite local filtering. If a point of a left ray has a predecessor, a bounded interval around it remains in that ray; absence of a minimum supplies such predecessors. The right case is dual.

\paragraph{Correction and reconstruction (ours).} Our expansion of the displayed definitions and their hypotheses; no additional claim is attributed to the manuscript.

\paragraph{Independent second check.} The finite-cone lemma of G independently supplies the simultaneous lower and upper constraints needed for an intersection of several basic sets.

\paragraph{Confidence.} High for mathematics; decisive wording checked in READINGS.tsv.

\section*{H0KB.53}\label{statement:H0KB.53}
\paragraph{Locator.} 156-3/022.12

\paragraph{His words.} Une extrémité de L est fermée; en général un singleton n’est pas nécessairement fermé.

\paragraph{How he left it.} stated

\paragraph{Precise statement (ours).} Any ambient least or greatest point is canonically closed. General points in G need not be closed.

\paragraph{Status.} True as written

\paragraph{Complete proof / counterexample application (ours).} P constructs, around every other point, a strict-ray neighborhood missing the endpoint. In D a singleton’s closure is its two-point height fiber, so it is not closed.

\paragraph{Correction and reconstruction (ours).} Our expansion of the displayed definitions and their hypotheses; no additional claim is attributed to the manuscript.

\paragraph{Independent second check.} The pulled-back rational topology in D directly verifies the negative clause, independently of the endpoint argument.

\paragraph{Confidence.} High for mathematics; decisive wording checked in READINGS.tsv.

\section*{H0KB.54}\label{statement:H0KB.54}
\paragraph{Locator.} 156-3/022.16;156-3/023.1

\paragraph{His words.} Pour les lieux d’une ficelle topologique, b\ensuremath{\in}cl\{a\} équivaut à b\ensuremath{\supset}a comme parties de H; l’égalité des adhérences équivaut à a=b, donc T\textsubscript{0}.

\paragraph{How he left it.} stated

\paragraph{Precise statement (ours).} For all genuine ficelles under the written 156-2 definitions, the canonical lieu specialization is claimed to be reverse inclusion of underlying closed subsets, and the lieu topology T\textsubscript{0}.

\paragraph{Status.} False

\paragraph{Complete proof / counterexample application (ours).} V gives a compact finite-tree ficelle: a main interval with a side hair attached at v=1/3. The loci D\_s=\{3/4\}\ensuremath{\cup}[s,1]\_hair are distinct regular closed lieux with exactly the same strict predecessors and successors. The full lieu order satisfies every G axiom, as V proves. Thus its topology is not T\textsubscript{0}, and the claimed reverse-inclusion equivalence fails. V also proves the strongest retained implication, C\ensuremath{\subset}D\ensuremath{\Rightarrow}D\ensuremath{\in}cl\{C\}, the exact strict-cone specialization criterion, and the additional connected-component condition W\_D that restores the converse. Connected target loci satisfy W\_D.

\paragraph{Correction and reconstruction (ours).} Our initial positive universal proof is withdrawn. A component-ficelle level cut need not be a global lieu. The corrected theorem assumes X\_\{A,D\}/A and X\_\{B,D\}/B connected, for the unpunctured component spaces, or the stronger sufficient condition that D is connected. These are additional hypotheses, not supplied ink. The real-segment case remains valid.

\paragraph{Independent second check.} The complete tree classification gives the order formula V1 and computes the identical strict cones independently of neighborhood pictures. check\_tree\_ficelle.py separately computes relative connected components and end contractions on a finite grid; all 56,644 pair comparisons match V1. The continuous density/local-filtering proof is separate. The repaired theorem checks the missing global connectedness by a closed quotient with connected fibers.

\paragraph{Confidence.} High: original definitions and the unqualified promotion sentence were visually checked; the finite-tree model verifies every source hypothesis and all G axioms; an independent graph calculation confirms the order formula.

\noindent Refined statement identifier: H0KB.54.correction (the corrected statement and its proof in this section).

\section*{H0KB.55}\label{statement:H0KB.55}
\paragraph{Locator.} 156-3/023.3

\paragraph{His words.} Il ne sait pas si le caractère T\textsubscript{0} vaut « pour tout tronçon ».

\paragraph{How he left it.} asked

\paragraph{Precise statement (ours).} Every G with its canonical topology is T\textsubscript{0}.

\paragraph{Status.} False

\paragraph{Complete proof / counterexample application (ours).} D satisfies every G axiom but its equal-height twins have identical neighborhoods. V strengthens this to actual lieux of a compact finite-tree ficelle. Under E the canonical topology is T\textsubscript{1}. More exactly, P proves T\textsubscript{0} iff x\ensuremath{\mapsto}(L\_<x,L\_>x) is injective; this is a necessary and sufficient order-language correction.

\paragraph{Correction and reconstruction (ours).} E is a sufficient positive repair for all G. The exact strict-cone injectivity criterion is proved in P. A general ficelle is not a valid automatic T\textsubscript{0} repair.

\paragraph{Independent second check.} The height projection realizes the entire canonical topology as the pullback of Q, so no possible canonical open distinguishes the twins.

\paragraph{Confidence.} High for mathematics; decisive wording checked in READINGS.tsv.

\noindent Refined statement identifier: H0KB.55.correction (the corrected statement and its proof in this section).

\section*{H0KB.56}\label{statement:H0KB.56}
\paragraph{Locator.} 156-3/023.4

\paragraph{His words.} Il demande si un tronçon est sobre : toute partie fermée irréductible aurait un point générique.

\paragraph{How he left it.} asked

\paragraph{Precise statement (ours).} Every G is sober; even the stronger restricted question whether every T\textsubscript{0} G is sober has a negative answer.

\paragraph{Status.} False

\paragraph{Complete proof / counterexample application (ours).} The nondegenerate-interval model J in P is G and T\textsubscript{0}. Its irreducible closed family F=\{I:0\ensuremath{\in}I\} has no generic point because the intersection of its members is the excluded degenerate interval \{0\}. All hypotheses and irreducibility are proved in P. A valid correction is the Hausdorff case: two distinct points of an irreducible closed set have disjoint neighborhoods meeting it, impossible; thus such a set is a singleton and has its unique generic point.

\paragraph{Correction and reconstruction (ours).} EGA I 0\_I 2.1.1–2.1.3 and SGA4 IV 4.2.1 supply the period meanings of irreducible, generic point, T\textsubscript{0} and sober.

\paragraph{Independent second check.} The explicit shrinking family [\ensuremath{-}1/n,1/n] independently proves the obstruction. The Hausdorff repair is a direct argument, not a purported implication T\textsubscript{1}\ensuremath{\Rightarrow}sober.

\paragraph{Confidence.} High for mathematics; decisive wording checked in READINGS.tsv.

\noindent Refined statement identifier: H0KB.56.correction (the corrected statement and its proof in this section).

\section*{H0KB.57}\label{statement:H0KB.57}
\paragraph{Locator.} 156-3/023.8

\paragraph{His words.} Si L est totalement ordonné, les points sont fermés, L est séparé, et [a,b], L\_\ensuremath{\geq}a, L\_\ensuremath{\leq}a sont fermés.

\paragraph{How he left it.} stated

\paragraph{Precise statement (ours).} The canonical topology of a totally ordered G is Hausdorff; every principal closed interval and ray is topologically closed.

\paragraph{Status.} True as written

\paragraph{Complete proof / counterexample application (ours).} For x<y insert c with x<c<y. The strict rays below and above c are disjoint neighborhoods. Complements of closed rays are opposite strict rays, and the complement of [a,b] is the union of the two exterior strict rays. This is P.

\paragraph{Correction and reconstruction (ours).} Our expansion of the displayed definitions and their hypotheses; no additional claim is attributed to the manuscript.

\paragraph{Independent second check.} The same conclusion follows from E for closedness and direct order separation for Hausdorffness; total dense orders satisfy E as proved in P.

\paragraph{Confidence.} High for mathematics; decisive wording checked in READINGS.tsv.

\section*{H0KB.58}\label{statement:H0KB.58}
\paragraph{Locator.} 156-3/023.8m

\paragraph{His words.} À propos de « les pts sont fermés », il ajoute « réciproque ? ».

\paragraph{How he left it.} asked

\paragraph{Precise statement (ours).} If every point in the canonical topology of G is closed, must the order be total?

\paragraph{Status.} False

\paragraph{Complete proof / counterexample application (ours).} The reflexive strict-coordinate order on R\textsuperscript{2} in P satisfies every G axiom and has the Euclidean canonical topology, yet (0,1) and (1,0) are incomparable. P proves the exact replacement: for every x\ensuremath{\ne}y, it is not the case that both L\_<x\ensuremath{\subset}L\_<y and L\_>x\ensuremath{\subset}L\_>y. Equivalently a strict predecessor or successor separates the ordered pair. This condition is necessary and sufficient for T\textsubscript{1} and does not assert totality.

\paragraph{Correction and reconstruction (ours).} Use the exact componentwise strict-cone antichain criterion in P. The weak coordinate order from the prior sweep fails G; the strict-coordinate model checks every local cone condition.

\paragraph{Independent second check.} The two-coordinate proof separately verifies both local cone axioms, correcting the earlier sweep’s invalid weak-product example.

\paragraph{Confidence.} High for mathematics; decisive wording checked in READINGS.tsv.

\noindent Refined statement identifier: H0KB.58.correction (the corrected statement and its proof in this section).

\section*{H0KB.59}\label{statement:H0KB.59}
\paragraph{Locator.} 156-3/023.12;156-3/024.1

\paragraph{His words.} Il définit une autre topologie dont [a,b] et les demi-droites fermées sont une sous-base de fermés, puis note que ]a,b[ ne veut plus être ouvert.

\paragraph{How he left it.} stated

\paragraph{Precise statement (ours).} The principal closed pieces generate a topology \ensuremath{\tau}\_c by arbitrary intersections of finite unions. It need not contain the canonical topology \ensuremath{\tau}.

\paragraph{Status.} True as written

\paragraph{Complete proof / counterexample application (ours).} Any family of subsets is a closed subbasis once finite unions and arbitrary intersections, including empty operations, are allowed. In P every basic \ensuremath{\tau}\_c-open contains (n,\ensuremath{-}n) for all sufficiently large n, so none can lie in a bounded open rectangle. Such a rectangle is \ensuremath{\tau}-open, hence \ensuremath{\tau}⊄\ensuremath{\tau}\_c.

\paragraph{Correction and reconstruction (ours).} Our expansion of the displayed definitions and their hypotheses; no additional claim is attributed to the manuscript.

\paragraph{Independent second check.} Check the escape point against each kind of closed generator: bounded interval, lower ray, upper ray. This verifies the whole finite-union argument.

\paragraph{Confidence.} High for mathematics; decisive wording checked in READINGS.tsv.

\section*{H0KB.60}\label{statement:H0KB.60}
\paragraph{Locator.} 156-3/024.3

\paragraph{His words.} Pour la topologie borne sup, il liste des ouverts stricts, les complémentaires des demi-droites fermées, puis [a,b] sans complément; il dit « base ».

\paragraph{How he left it.} stated

\paragraph{Precise statement (ours).} The supremum topology is generated by the union of the two open subbases; the final generator must be L\ensuremath{\setminus}[a,b], and finite intersections give a basis.

\paragraph{Status.} Slip

\paragraph{Complete proof / counterexample application (ours).} On R both original topologies are the usual order topology, but declaring [0,1] open strictly refines it. Thus the literal list is not their supremum. Taking complements of every closed-subbasis member, and finite intersections with canonical basic opens, is exactly the defining least topology containing both.

\paragraph{Correction and reconstruction (ours).} Our correction uses the preceding phrase « topologie borne sup » and the two adjacent complemented ray generators.

\paragraph{Independent second check.} The crop confirms the missing complement. The real-line test distinguishes the literal topology from the intended join independently of any non-total example.

\paragraph{Confidence.} High for mathematics; decisive wording checked in READINGS.tsv.

\noindent Refined statement identifier: H0KB.60.correction (the corrected statement and its proof in this section).

\section*{H0KB.61}\label{statement:H0KB.61}
\paragraph{Locator.} 156-3/024.8

\paragraph{His words.} « Aura-t-on » cl(]a,b[)=[a,b] dans la topologie borne sup ?

\paragraph{How he left it.} asked

\paragraph{Precise statement (ours).} For every G and a<b, the closure of the strict interval in \ensuremath{\tau}\ensuremath{\vee}\ensuremath{\tau}\_c equals the closed endpoint interval.

\paragraph{Status.} False

\paragraph{Complete proof / counterexample application (ours).} In D choose a\_0 at height 0, its twin a\_1, and b\_0 at height 1. The \ensuremath{\tau}\_c-closed interval [a\_1,b\_0] contains (a\_0,b\_0) and omits a\_0, so a\_0 is not in its closure in the join topology. Under E, \ensuremath{\tau}\_c\ensuremath{\subset}\ensuremath{\tau} and the join is \ensuremath{\tau}; E and local filtering prove the desired closure equality.

\paragraph{Correction and reconstruction (ours).} E is our proved sufficient repair; the source’s intended join topology is used, after the separate generator slip is corrected.

\paragraph{Independent second check.} The genuine segment-ficelle counterexample in V gives the same failure with actual closed interval lieux, independently of D.

\paragraph{Confidence.} High for mathematics; decisive wording checked in READINGS.tsv.

\noindent Refined statement identifier: H0KB.61.correction (the corrected statement and its proof in this section).

\section*{H0KB.62}\label{statement:H0KB.62}
\paragraph{Locator.} 156-3/024.9

\paragraph{His words.} Il demande cl(L\_<a)=L\_\ensuremath{\leq}a si a n’est pas minimum, et la formule duale si a n’est pas maximum.

\paragraph{How he left it.} asked

\paragraph{Precise statement (ours).} The two nonempty strict-ray closure identities are asserted in \ensuremath{\tau}\ensuremath{\vee}\ensuremath{\tau}\_c for arbitrary G.

\paragraph{Status.} False

\paragraph{Complete proof / counterexample application (ours).} In D the negative ray below a\_0 lies in the \ensuremath{\tau}\_c-closed ray \ensuremath{\leq}a\_1, which excludes a\_0. Thus its closure misses a\_0; the dual construction refutes the other identity. Nonextremality is satisfied. Under E both equalities hold by the local-neighborhood accumulation argument, and the join topology is \ensuremath{\tau}.

\paragraph{Correction and reconstruction (ours).} The endpoint qualifications are preserved exactly; they are insufficient without an additional condition such as E.

\paragraph{Independent second check.} The two strict rays are nonempty at every rational height, so the counterexample does not exploit an omitted endpoint qualification.

\paragraph{Confidence.} High for mathematics; decisive wording checked in READINGS.tsv.

\noindent Refined statement identifier: H0KB.62.correction (the corrected statement and its proof in this section).

\section*{H0KB.63}\label{statement:H0KB.63}
\paragraph{Locator.} 156-3/024.15

\paragraph{His words.} Pour montrer a\ensuremath{\in}cl(]a,b[), il suffit que tout ouvert élémentaire contenant a contienne ]a,\ensuremath{\epsilon}[ avec a<\ensuremath{\epsilon}<b.

\paragraph{How he left it.} stated

\paragraph{Precise statement (ours).} For the intended subbasis of \ensuremath{\tau}\ensuremath{\vee}\ensuremath{\tau}\_c, that property is sufficient for a to lie in the closure; the property itself need not hold.

\paragraph{Status.} True as written

\paragraph{Complete proof / counterexample application (ours).} A neighborhood contains a finite intersection of elementary opens. Suppose each contains (a,\ensuremath{\epsilon}\_i) with a<\ensuremath{\epsilon}\_i<b. Local downward filtering above a gives a common lower bound d>a of the \ensuremath{\epsilon}\_i, and density gives a<e<d. Then the nonempty interval (a,e) lies in every elementary open and in (a,b). Hence every neighborhood meets (a,b).

\paragraph{Correction and reconstruction (ours).} The sufficient implication is separated from the false universal premise in the next row.

\paragraph{Independent second check.} This proof explicitly handles finite intersections, which a subbasis-only argument would otherwise leave unjustified.

\paragraph{Confidence.} High for mathematics; decisive wording checked in READINGS.tsv.

\section*{H0KB.64}\label{statement:H0KB.64}
\paragraph{Locator.} 156-3/024.19

\paragraph{His words.} « Mais c’est faux déjà » pour les lieux d’une ficelle; il propose un ouvert élémentaire complémentaire d’un [x,y] contenant tout ]a,b[ mais pas a.

\paragraph{How he left it.} stated

\paragraph{Precise statement (ours).} There exists a genuine ficelle lieu poset and an elementary join-topology neighborhood of a that misses the whole interval (a,b).

\paragraph{Status.} True as written

\paragraph{Complete proof / counterexample application (ours).} In V, H=[0,1], a=[3/10,1/2], b=\{4/5\}, x=\{2/5\}. Then x<b and a\ensuremath{\notin}[x,b], but every z with a<z<b lies in [x,b]. Its complement is the required neighborhood. V verifies that all these are lieux under the actual 156-2 definition and that all G conditions hold.

\paragraph{Correction and reconstruction (ours).} This confirms his narrower example; it is not merely a non-ficelle counterexample to the general question.

\paragraph{Independent second check.} Independently check the endpoint inequalities r(a)=1/2<l(z) and r(x)=2/5<l(z); equality cases for x,a are excluded because they overlap.

\paragraph{Confidence.} High for mathematics; decisive wording checked in READINGS.tsv.

\section*{H0KB.65}\label{statement:H0KB.65}
\paragraph{Locator.} 156-3/025.1

\paragraph{His words.} Il revient à la topologie initiale et appelle néanmoins [a,b], les demi-droites et L sous-tronçons « fermés », « par abus de langage »; les points sont admis comme larges.

\paragraph{How he left it.} stated

\paragraph{Precise statement (ours).} From here “closed” in the tame calculus means a finite closed-endpoint presentation, not necessarily closed in \ensuremath{\tau}. Large/pseudo pieces additionally allow singletons.

\paragraph{Status.} True as written

\paragraph{Complete proof / counterexample application (ours).} The distinction is coherent: the closed-parts calculus F uses order endpoints only, while D shows these need not be \ensuremath{\tau}-closed. The retained definition makes no false topological promise.

\paragraph{Correction and reconstruction (ours).} Our expansion of the displayed definitions and their hypotheses; no additional claim is attributed to the manuscript.

\paragraph{Independent second check.} Compare the explicit abuse-of-language sentence with the cancelled attempted topological-boundary definition below it and the warning on scan 30.

\paragraph{Confidence.} High for mathematics; decisive wording checked in READINGS.tsv.

\section*{H0KB.66}\label{statement:H0KB.66}
\paragraph{Locator.} 156-3/026.1

\paragraph{His words.} Il pose \ensuremath{\partial}T égal au « bord absolu », l’ensemble des plus petit et plus grand éléments de T; \ensuremath{\partial}\{a\}=\{a\}, \ensuremath{\partial}[a,b]=\{a,b\}.

\paragraph{How he left it.} stated

\paragraph{Precise statement (ours).} For a nonempty closed pseudo-piece T, its absolute boundary is the set of its existing least and greatest elements; singleton and bounded-interval formulas follow.

\paragraph{Status.} True as written

\paragraph{Complete proof / counterexample application (ours).} Least and greatest elements are unique when they exist. For [a,b], a is least and b greatest; if a=b the set is the singleton. This is precisely the stated definition, distinct from relative frontier.

\paragraph{Correction and reconstruction (ours).} Our expansion of the displayed definitions and their hypotheses; no additional claim is attributed to the manuscript.

\paragraph{Independent second check.} Directly evaluate L of each of the three endpoint types; its absolute boundary is its ambient extrema, not always empty.

\paragraph{Confidence.} High for mathematics; decisive wording checked in READINGS.tsv.

\section*{H0KB.67}\label{statement:H0KB.67}
\paragraph{Locator.} 156-3/026.9

\paragraph{His words.} Il écrit \ensuremath{\partial}L\_\ensuremath{\geq}a=\ensuremath{\partial}L\_\ensuremath{\leq}a=\{a\}, avec la marge « sans exclure le cas a\ensuremath{\in}\ensuremath{\partial}L ».

\paragraph{How he left it.} stated

\paragraph{Precise statement (ours).} The literal absolute-boundary ray formulas omit a possible opposite ambient extremum.

\paragraph{Status.} Slip

\paragraph{Complete proof / counterexample application (ours).} For L=[0,1] and a=1/2, the absolute boundary of L\_\{\ensuremath{\geq}a\} is \{a,1\}, not \{a\}; that of L\_\{\ensuremath{\leq}a\} is \{0,a\}. The correct formulas adjoin the ambient maximum or minimum when it exists, with repetitions removed. They follow immediately from the extremum definition.

\paragraph{Correction and reconstruction (ours).} Our slip correction follows the immediately preceding absolute-boundary definition; it is not a reversion to relative frontier.

\paragraph{Independent second check.} The scan margin expressly retains endpoint cases; it does not eliminate this counterexample with a strictly internal a.

\paragraph{Confidence.} High for mathematics; decisive wording checked in READINGS.tsv.

\noindent Refined statement identifier: H0KB.67.correction (the corrected statement and its proof in this section).

\section*{H0KB.68}\label{statement:H0KB.68}
\paragraph{Locator.} 156-3/026.11

\paragraph{His words.} Il compare frontière relative et bord absolu, puis encadre \ensuremath{\partial}T=(Ṫ\ensuremath{\cap}T) \ensuremath{\sqcup} (T\ensuremath{\cap}\ensuremath{\partial}L), après avoir barré une exception pour les singletons.

\paragraph{How he left it.} stated

\paragraph{Precise statement (ours).} Ordinary union gives the correct relation between the two boundaries for a closed piece. The displayed disjoint-union assertion fails for a singleton at an ambient extremum.

\paragraph{Status.} Slip

\paragraph{Complete proof / counterexample application (ours).} In L=[0,1], T=\{0\} has absolute and relative boundary \{0\}. Both summands in the boxed formula are \{0\}, so they are not disjoint. Ordinary union remains correct in every case. Alternatively restore the cancelled singleton exception and use disjoint union only for non-singletons.

\paragraph{Correction and reconstruction (ours).} The earlier exception is visibly cancelled and the union mark visibly underlined/disjoint; the literal formula is retained as a slip, not silently repaired.

\paragraph{Independent second check.} For a non-singleton piece, relative frontier omits precisely its ambient extreme endpoints; for a singleton it retains its sole point. These two cases prove the corrected rule.

\paragraph{Confidence.} High for mathematics; decisive wording checked in READINGS.tsv.

\noindent Refined statement identifier: H0KB.68.correction (the corrected statement and its proof in this section).

\section*{H0KB.69}\label{statement:H0KB.69}
\paragraph{Locator.} 156-3/026.20;156-3/026.20m

\paragraph{His words.} Les tronçons sont T\ensuremath{\prime}=T\textbackslash{}\ensuremath{\alpha}, \ensuremath{\alpha}\ensuremath{\subset}\ensuremath{\partial}T; la marge dit « Cette écriture est unique » et que T est le plus petit fermé contenant T\ensuremath{\prime}.

\paragraph{How he left it.} stated

\paragraph{Precise statement (ours).} Every nonempty open/half-open interval has a unique closed-endpoint presentation and a least closed tame hull under G.

\paragraph{Status.} False

\paragraph{Complete proof / counterexample application (ours).} In D, V=L\_\{<a\_0\}=L\_\{<a\_1\} has two presentations (\ensuremath{\leq}a\_i)\textbackslash{}\{a\_i\}; D proves that it has no least closed tame hull. Under E, cofinality and coinitiality of interval interiors force the actual endpoints, giving the unique presentation and least hull. E is also necessary for the universal ray-hull claim.

\paragraph{Correction and reconstruction (ours).} E is the exact additional all-rays hypothesis, stronger and sharper than merely assuming totality. Empty pieces must be excluded from uniqueness.

\paragraph{Independent second check.} The independent canonical-topology characterization in E yields the same hulls under E and proves their leastness among closed tame supersets.

\paragraph{Confidence.} High for mathematics; decisive wording checked in READINGS.tsv.

\noindent Refined statement identifier: H0KB.69.correction (the corrected statement and its proof in this section).

\section*{H0KB.70}\label{statement:H0KB.70}
\paragraph{Locator.} 156-3/026.23

\paragraph{His words.} Dans la liste des tronçons ouverts, l’encre exclut a=\ensuremath{\omega} pour L\_<a et a=\ensuremath{\alpha} pour L\_>a.

\paragraph{How he left it.} stated

\paragraph{Precise statement (ours).} Nonempty strict lower rays require a different from the least element \ensuremath{\alpha}, and nonempty strict upper rays require a different from the greatest element \ensuremath{\omega}. The two written exclusions are reversed.

\paragraph{Status.} Slip

\paragraph{Complete proof / counterexample application (ours).} In L=[0,1], the written a\ensuremath{\ne}\ensuremath{\omega} condition allows a=0 in L\_<a, giving the empty set, and wrongly excludes a=1 although L\_<1 is nonempty. The upper-ray case is dual. By G a point with no strict predecessor is the least element; conversely a least element has no predecessor. This proves the exact corrected exclusions, and the dual proves the upper-ray condition.

\paragraph{Correction and reconstruction (ours).} Our correction reverses the two endpoint symbols. Witness H prints the mathematically corrected exclusions; the crop shows the opposite order, recorded explicitly in READINGS.tsv.

\paragraph{Independent second check.} Directly check the four rays at 0 and 1. Then use the global-lower-bound proof that a minimal point must be least, which checks the general claim independently.

\paragraph{Confidence.} High for the mathematical correction; high for the two enlarged endpoint glyphs, compared with \ensuremath{\alpha} and \ensuremath{\omega} elsewhere on the pages.

\noindent Refined statement identifier: H0KB.70.correction (the corrected statement and its proof in this section).

\section*{H0KB.71}\label{statement:H0KB.71}
\paragraph{Locator.} 156-3/027.1

\paragraph{His words.} Les intervalles stricts et demi-droites strictes sont dits ouverts; il écrit ]a,b] et [b,a[ pour les demi-ouverts, puis note le cas d’une extrémité ambiante.

\paragraph{How he left it.} stated

\paragraph{Precise statement (ours).} The second half-open interval should be [a,b), for a<b. Strict intervals/rays are \ensuremath{\tau}-open; an interval beginning at the ambient minimum is also \ensuremath{\tau}-open when its other endpoint is omitted.

\paragraph{Status.} Slip

\paragraph{Complete proof / counterexample application (ours).} With 0<1 in Q, [1,0) is empty and is not the intended nonempty half-open piece. Replace it by [0,1). The basis definition proves strict openness. If a=min L, [a,b)=L\_\{<b\}; dually at a maximum. Hence his following qualification is correct.

\paragraph{Correction and reconstruction (ours).} The bracket/endpoint reversal is ours to correct; the scan and supplied text record [b,a[. The apparent “exactly the open ones” is read with his immediate exception, not as an unqualified biconditional.

\paragraph{Independent second check.} The example at the ambient minimum distinguishes the classification by endpoint notation from topological openness.

\paragraph{Confidence.} High for mathematics; decisive wording checked in READINGS.tsv.

\noindent Refined statement identifier: H0KB.71.correction (the corrected statement and its proof in this section).

\section*{H0KB.72}\label{statement:H0KB.72}
\paragraph{Locator.} 156-3/027.12

\paragraph{His words.} Un pseudo-tronçon est un point ou un tronçon; pour T\ensuremath{\prime}=T\textbackslash{}\ensuremath{\alpha} il pose \ensuremath{\partial}T\ensuremath{\prime}=\ensuremath{\partial}T, en précisant qu’il ne s’agit pas du bord absolu induit.

\paragraph{How he left it.} stated

\paragraph{Precise statement (ours).} The construction is asserted to define an intrinsic endpoint set for each actual subset T\ensuremath{\prime}, equal to the endpoints of its associated closed hull.

\paragraph{Status.} False

\paragraph{Complete proof / counterexample application (ours).} The two presentations of V in D give respectively endpoint sets \{a\_0\} and \{a\_1\}, so the function on subsets is not well defined. Under E, the unique infimum/supremum endpoints proved in E make it intrinsic; F then supplies the pseudo-piece calculus.

\paragraph{Correction and reconstruction (ours).} Keep the distinction between pseudo-boundary and absolute induced boundary; E corrects the intrinsicness failure without conflating them.

\paragraph{Independent second check.} The ambiguity survives before any algebraic operation is attempted: compare the two endpoint sets directly.

\paragraph{Confidence.} High for mathematics; decisive wording checked in READINGS.tsv.

\noindent Refined statement identifier: H0KB.72.correction (the corrected statement and its proof in this section).

\section*{H0KB.73}\label{statement:H0KB.73}
\paragraph{Locator.} 156-3/027.17

\paragraph{His words.} Il commence « on aura encore » une relation entre \ensuremath{\partial}T\ensuremath{\prime}, sa frontière relative et \ensuremath{\partial}L, puis laisse des points de suspension; la marge distingue le bord absolu.

\paragraph{How he left it.} programme

\paragraph{Precise statement (ours).} Complete the intended absolute-endpoint/relative-frontier formula for pseudo-pieces. Supported precise version under E: relative frontier removes included ambient ends in non-singleton pieces and retains isolated singleton ends.

\paragraph{Status.} programme

\paragraph{Complete proof / counterexample application (ours).} The incomplete display does not supply a determinate universal identity to refute. The precise endpoint/frontier rule is proved in F and is written explicitly on scan 30. The intrinsic-endpoint claim under G is separately refuted in H0KB.72.

\paragraph{Correction and reconstruction (ours).} The supported precise rule is included as the correction in H0KB.84 and proved there.

\paragraph{Independent second check.} The crop preserves the ellipsis and the distinction between dotted frontier and absolute-boundary notation. No set equality is invented from the unfinished line.

\paragraph{Confidence.} High for mathematics; decisive wording checked in READINGS.tsv.

\section*{H0KB.74}\label{statement:H0KB.74}
\paragraph{Locator.} 156-3/027.18m

\paragraph{His words.} Il note que les points frontière sont les bornes inférieure et supérieure dans L quand elles existent, et qu’une borne supérieure existe ssi le pseudo-tronçon est majoré.

\paragraph{How he left it.} stated

\paragraph{Precise statement (ours).} Every bounded pseudo-piece in G has an actual supremum/infimum equal to its displayed finite endpoint; its endpoint set is empty only for a full endpoint-free L.

\paragraph{Status.} False

\paragraph{Complete proof / counterexample application (ours).} In D the negative strict ray is bounded by both twins at height 0 but has no supremum: any supremum would have to be below both twins and bound all negative heights, which is impossible. The dual refutes the infimum clause. Under E, the cofinality argument gives the exact displayed endpoints and the asserted classification of empty boundary.

\paragraph{Correction and reconstruction (ours).} The implication “has a supremum \ensuremath{\Rightarrow} bounded” is always true. E supplies the missing converse for these pieces.

\paragraph{Independent second check.} Use the cofinal sequence of negative heights and inspect all possible heights of a candidate supremum. This independently excludes both a negative candidate and either zero-height twin.

\paragraph{Confidence.} High for mathematics; decisive wording checked in READINGS.tsv.

\noindent Refined statement identifier: H0KB.74.correction (the corrected statement and its proof in this section).

\section*{H0KB.75}\label{statement:H0KB.75}
\paragraph{Locator.} 156-3/027.19

\paragraph{His words.} Deux pseudo-tronçons sont modérés si \ensuremath{\partial}T\ensuremath{\cup}\ensuremath{\partial}T\ensuremath{\prime} est totalement ordonné.

\paragraph{How he left it.} stated

\paragraph{Precise statement (ours).} This rule is asserted to define compatibility of actual pseudo-piece subsets, independently of presentations.

\paragraph{Status.} False

\paragraph{Complete proof / counterexample application (ours).} Take V=L\_\{<a\_0\}=L\_\{<a\_1\} in D and Y=\{a\_0\}. The presentation with endpoint a\_0 makes V compatible with Y; the presentation with endpoint a\_1 does not. Under E endpoints are unique and the rule is well defined, with the complete F calculus.

\paragraph{Correction and reconstruction (ours).} Correct either by retaining presentations as data, which changes the object, or by imposing E to preserve his intended actual-subset object. We prove the latter.

\paragraph{Independent second check.} The two presentations are identical as subsets and differ only in the omitted endpoint; both satisfy the manuscript’s construction.

\paragraph{Confidence.} High for mathematics; decisive wording checked in READINGS.tsv.

\noindent Refined statement identifier: H0KB.75.correction (the corrected statement and its proof in this section).

\section*{H0KB.76}\label{statement:H0KB.76}
\paragraph{Locator.} 156-3/028.1

\paragraph{His words.} Deux pseudo-tronçons en position relative modérée ont une intersection vide ou pseudo-tronçon, avec \ensuremath{\partial}(T\ensuremath{\cap}T\ensuremath{\prime})\ensuremath{\subset}\ensuremath{\partial}T\ensuremath{\cup}\ensuremath{\partial}T\ensuremath{\prime}.

\paragraph{How he left it.} stated

\paragraph{Precise statement (ours).} For compatible endpoint presentations, the intersection has the max/min endpoint formula. As an intrinsic assertion about actual subsets on all G, it inherits the endpoint ambiguity.

\paragraph{Status.} False

\paragraph{Complete proof / counterexample application (ours).} D supplies actual subsets for which even the compatibility premise and claimed endpoint conclusion depend on the chosen presentation, as H0KB.75 shows. Under E, intersect the two intervals by choosing their larger lower endpoint and smaller upper endpoint; tied endpoints are retained iff both input flags retain them. Empty and singleton cases follow, and every surviving endpoint is old.

\paragraph{Correction and reconstruction (ours).} The algebraic intersection formula is valid for fixed compatible presentations; E makes the manuscript’s claimed subset-level statement well defined.

\paragraph{Independent second check.} The endpoint-pair computation independently agrees with the normalized-mask intersection identity in F.

\paragraph{Confidence.} High for mathematics; decisive wording checked in READINGS.tsv.

\noindent Refined statement identifier: H0KB.76.correction (the corrected statement and its proof in this section).

\section*{H0KB.77}\label{statement:H0KB.77}
\paragraph{Locator.} 156-3/028.4

\paragraph{His words.} T<T\ensuremath{\prime} signifie t<t\ensuremath{\prime} pour tout t\ensuremath{\in}T,t\ensuremath{\prime}\ensuremath{\in}T\ensuremath{\prime}; c’est transitif. Il ajoute que cela implique la position relative modérée et la disjonction.

\paragraph{How he left it.} stated

\paragraph{Precise statement (ours).} For nonempty actual pieces the relation is transitive and implies disjointness. Its assertion of intrinsic endpoint compatibility under G is not valid without the preceding well-definedness repair.

\paragraph{Status.} False

\paragraph{Complete proof / counterexample application (ours).} Transitivity follows by choosing an element of the intermediate nonempty piece; disjointness follows from irreflexivity. For endpoint compatibility in D, let T be the ray below height 0 presented with endpoint a\_0 and T\ensuremath{\prime}=\{a\_1\}. Then every point of T is <a\_1, but the endpoint pair a\_0,a\_1 is incomparable. Under E actual sup/inf bounds compare and all endpoints are compatible, proving the full corrected assertion.

\paragraph{Correction and reconstruction (ours).} Retain nonemptiness and add E for the full endpoint assertion; the transitive/disjoint part needs no E.

\paragraph{Independent second check.} The explicit example checks strict order on every point by height, while its two displayed boundary points have the same height and are distinct.

\paragraph{Confidence.} High for mathematics; decisive wording checked in READINGS.tsv.

\noindent Refined statement identifier: H0KB.77.correction (the corrected statement and its proof in this section).

\section*{H0KB.78}\label{statement:H0KB.78}
\paragraph{Locator.} 156-3/028.10

\paragraph{His words.} Après T<T\ensuremath{\prime}, il écrit que T, étant majoré, « a un plus grand élément a »; la phrase duale parle ensuite de \ensuremath{\partial}T\ensuremath{\prime}.

\paragraph{How he left it.} stated

\paragraph{Precise statement (ours).} The literal maximum-of-T assertion fails for open pieces. The nearby reasoning needs the greatest endpoint of its closed hull, namely sup T.

\paragraph{Status.} Slip

\paragraph{Complete proof / counterexample application (ours).} On R take T=(0,1), T\ensuremath{\prime}=(2,3). They are strictly ordered, T is bounded, and has no greatest element. Its supremum is 1, the greatest endpoint of the closed hull. Under E every bounded pseudo-piece has that supremum by E. Without E, even the supremum claim fails as separately recorded in H0KB.74.

\paragraph{Correction and reconstruction (ours).} Our correction replaces “greatest element of T” by “greatest endpoint / sup T”, without pretending this fixes the independent missing-E problem.

\paragraph{Independent second check.} The crop confirms that the first surviving noun is T, while the next clause explicitly uses \ensuremath{\partial}T\ensuremath{\prime}; this supports the local pen-slip judgment.

\paragraph{Confidence.} High for mathematics; decisive wording checked in READINGS.tsv.

\noindent Refined statement identifier: H0KB.78.correction (the corrected statement and its proof in this section).

\section*{H0KB.79}\label{statement:H0KB.79}
\paragraph{Locator.} 156-3/028.17

\paragraph{His words.} T\ensuremath{\ll}T\ensuremath{\prime} renforce T<T\ensuremath{\prime} en excluant le raccordement; il l’exprime par l’existence d’un x avec T\ensuremath{\subset}L\_<x et T\ensuremath{\prime}\ensuremath{\subset}L\_>x.

\paragraph{How he left it.} stated

\paragraph{Precise statement (ours).} Use the intrinsic definition T\ensuremath{\ll}T\ensuremath{\prime} iff there is a strict separator x. It is transitive and forces disjoint, strictly ordered nonempty pieces. Where actual endpoint bounds exist, it is precisely non-raccordement.

\paragraph{Status.} True as written

\paragraph{Complete proof / counterexample application (ours).} If T\ensuremath{\ll}T\ensuremath{\prime}\ensuremath{\ll}T″, any separator of T,T\ensuremath{\prime} lies below every point of T\ensuremath{\prime} and hence below every point of T″, so also separates T,T″. The exact endpoint equivalence is proved in H under the explicit actual-bounds assumption on the next page.

\paragraph{Correction and reconstruction (ours).} A variable written a in the first cone of the last display is read as the locally quantified x in the corrected formula; the typo is recorded in READINGS.tsv.

\paragraph{Independent second check.} The separator definition uses only actual subsets, so it is immune to the endpoint ambiguity. We do not identify it with a presentation-dependent relation without the bounds hypothesis.

\paragraph{Confidence.} High for mathematics; decisive wording checked in READINGS.tsv.

\section*{H0KB.80}\label{statement:H0KB.80}
\paragraph{Locator.} 156-3/029.2

\paragraph{His words.} « Conditions équivalentes (si a=Sup T, b=Inf T\ensuremath{\prime}) » : a) T<T\ensuremath{\prime} sans raccordement; b) il existe un strict séparateur.

\paragraph{How he left it.} stated

\paragraph{Precise statement (ours).} Assume the actual supremum a of T and actual infimum b of T\ensuremath{\prime} exist. Then strict separation is equivalent to a<b, or a=b excluded from both pieces. This holds in G without E.

\paragraph{Status.} True as written

\paragraph{Complete proof / counterexample application (ours).} H proves a\ensuremath{\leq}b and computes all strict separators as [a,b] with a removed if a\ensuremath{\in}T and b removed if b\ensuremath{\in}T\ensuremath{\prime}. Density makes this set nonempty exactly in the stated cases. This proves the two-way implication with the manuscript’s inserted si retained.

\paragraph{Correction and reconstruction (ours).} We explicitly preserve the supplied conditional; an unconditional bound-existence assertion is a different, false row.

\paragraph{Independent second check.} Check all four endpoint-retention combinations independently. The actual-bound hypothesis is visually confirmed and never omitted in a counterexample.

\paragraph{Confidence.} High for mathematics; decisive wording checked in READINGS.tsv.

\section*{H0KB.81}\label{statement:H0KB.81}
\paragraph{Locator.} 156-3/029.7

\paragraph{His words.} Il ajoute un pseudo-tronçon S entre T,T\ensuremath{\prime} se raccordant à chacun, unique, avec \ensuremath{\partial}S=\{a,b\}; S est l’ensemble de tous les stricts séparateurs.

\paragraph{How he left it.} stated

\paragraph{Precise statement (ours).} Under the preceding actual-bounds assumption, the strict-separator subset exists with the stated interval formula. The further intrinsic endpoint/raccordement claims need not follow in arbitrary G.

\paragraph{Status.} False

\paragraph{Complete proof / counterexample application (ours).} In D take T=\{a\_0\}, T\ensuremath{\prime}=\{b\_0\}, heights 0 and 1. Their actual sup/inf are a\_0,b\_0, so the inserted hypothesis holds. The separator set is the strict height interval, whose lower and upper bounds include both endpoint twins and have no infimum/supremum. Its intrinsic \ensuremath{\partial} is not \{a\_0,b\_0\}. Under E, H proves the unique separator piece and both raccordements by its forced endpoint flags.

\paragraph{Correction and reconstruction (ours).} The set formula for S is retained. E makes its endpoint interpretation and uniqueness valid.

\paragraph{Independent second check.} This counterexample keeps the actual sup/inf of T,T\ensuremath{\prime}, so it refutes only the added conclusion about S, not the valid a)\ensuremath{\Longleftrightarrow}b) theorem.

\paragraph{Confidence.} High for mathematics; decisive wording checked in READINGS.tsv.

\noindent Refined statement identifier: H0KB.81.correction (the corrected statement and its proof in this section).

\section*{H0KB.82}\label{statement:H0KB.82}
\paragraph{Locator.} 156-3/029.15

\paragraph{His words.} Si T est « minoré strictement », ses minorants stricts forment le pseudo-tronçon interstitiel gauche, sans minorant, caractérisé par raccordement; dualement à droite.

\paragraph{How he left it.} stated

\paragraph{Precise statement (ours).} Under the intrinsic endpoint calculus, the exterior gap exists exactly when a strict bound exists. Its “without a lower bound” description must mean without a strict lower bound.

\paragraph{Status.} Slip

\paragraph{Complete proof / counterexample application (ours).} On L=[0,1], take T=[1/2,1]. Its strict lower bounds form [0,1/2), which has the lower bound 0 but no strict lower bound. Thus the literal “sans minorant” is false. With strict understood, H gives the exterior-ray formula and characterization under E; the independent G endpoint failure remains in H0KB.74.

\paragraph{Correction and reconstruction (ours).} Restore the strict qualifier consistently. Endpoint intrinsicness still uses the separately stated E repair.

\paragraph{Independent second check.} The manuscript explicitly inserts “strictement” and “minorants stricts” in the preceding lines, providing the local convention needed for this repair.

\paragraph{Confidence.} High for mathematics; decisive wording checked in READINGS.tsv.

\noindent Refined statement identifier: H0KB.82.correction (the corrected statement and its proof in this section).

\section*{H0KB.83}\label{statement:H0KB.83}
\paragraph{Locator.} 156-3/030.1

\paragraph{His words.} Les drapeaux finis \ensuremath{\Phi} de pseudo-tronçons pour \ensuremath{\ll} donnent X\_\ensuremath{\Phi}=\ensuremath{\bigcup}T\_i; l’application \ensuremath{\Phi}\ensuremath{\mapsto}X\_\ensuremath{\Phi} est injective. Les T\_i sont ses composantes, pas pour la topologie canonique.

\paragraph{How he left it.} stated

\paragraph{Precise statement (ours).} Using strict-separator \ensuremath{\ll} and actual nonempty convex pseudo-pieces, finite component decompositions are unique, even without E.

\paragraph{Status.} True as written

\paragraph{Complete proof / counterexample application (ours).} If a convex piece from one decomposition met two pieces in another, a separator between them would lie between two of its points and hence belong to it, but would be absent from the union. Thus each piece lies in a single piece of the other decomposition; reverse the roles to conclude equality. This is F.

\paragraph{Correction and reconstruction (ours).} We retain the intrinsic separator definition of \ensuremath{\ll}. “Components” is explicitly a tame, not topological, term.

\paragraph{Independent second check.} The two-containment argument is independent of endpoint uniqueness; hence the earlier endpoint counterexample does not refute this proposition.

\paragraph{Confidence.} High for mathematics; decisive wording checked in READINGS.tsv.

\section*{H0KB.84}\label{statement:H0KB.84}
\paragraph{Locator.} 156-3/030.16

\paragraph{His words.} Fr(X)= (\ensuremath{\bigcup}\ensuremath{\partial}T\_i)\textbackslash{}\ensuremath{\epsilon}; \ensuremath{\epsilon} retire les extrémités ambiantes appartenant à une composante non ponctuelle, mais non les points isolés.

\paragraph{How he left it.} stated

\paragraph{Precise statement (ours).} The formula defines an intrinsic finite relative frontier of an actual tame subset only if its component endpoint sets are intrinsic. It is correct under G+E with the stated singleton exception.

\paragraph{Status.} False

\paragraph{Complete proof / counterexample application (ours).} The one-component set V in D has two different endpoint presentations, giving different frontier sets; there are no ambient extrema to remove. Under E, F proves unique endpoints, and removing exactly the included ambient ends of non-singleton components yields the displayed relative frontier.

\paragraph{Correction and reconstruction (ours).} No topological frontier is substituted. The page itself warns that this frontier is not for the canonical topology.

\paragraph{Independent second check.} The interval cases L, a proper ray, an internal singleton and an ambient-end singleton independently verify every clause of the corrected formula.

\paragraph{Confidence.} High for mathematics; decisive wording checked in READINGS.tsv.

\noindent Refined statement identifier: H0KB.84.correction (the corrected statement and its proof in this section).

\section*{H0KB.85}\label{statement:H0KB.85}
\paragraph{Locator.} 156-3/031.1

\paragraph{His words.} Un point est modéré avec X s’il l’est avec chaque T\_i, donc avec \ensuremath{\bigcup}\ensuremath{\partial}T\_i, ou avec Fr(X); tout point de X et tout point frontière a cette propriété.

\paragraph{How he left it.} stated

\paragraph{Precise statement (ours).} With intrinsic endpoints, point compatibility is equivalent to comparability with all component endpoints; removing ambient extrema does not change it. Points of X and its frontier satisfy it.

\paragraph{Status.} False

\paragraph{Complete proof / counterexample application (ours).} The presentation-dependent V in D makes even this relation on actual subsets ambiguous. Under E, a point in one component lies between its endpoints and is strictly ordered from all other components; a frontier point belongs to the common endpoint chain. Ambient extrema are universally comparable, so they may be removed. These observations prove every stated implication.

\paragraph{Correction and reconstruction (ours).} E repairs the well-definedness; the elementary endpoint comparison is otherwise correct.

\paragraph{Independent second check.} Check an interior point, an omitted endpoint, and an ambient extremum separately. The same comparisons give the all-points polar formula O1 for outside points.

\paragraph{Confidence.} High for mathematics; decisive wording checked in READINGS.tsv.

\noindent Refined statement identifier: H0KB.85.correction (the corrected statement and its proof in this section).

\section*{H0KB.86}\label{statement:H0KB.86}
\paragraph{Locator.} 156-3/031.8

\paragraph{His words.} X\ensuremath{\prime}=C\_mod(X) est modéré et est la plus grande partie modérée disjointe de X en position modérée avec X.

\paragraph{How he left it.} stated

\paragraph{Precise statement (ours).} The indicated construction is asserted to define a subset-level complement for every G, with the stated maximality.

\paragraph{Status.} False

\paragraph{Complete proof / counterexample application (ours).} On D the complement of \{a\_0\} equals that of \{a\_1\}; applying the claimed endpoint rule to that common set is ambiguous. Thus the full intrinsic complement construction cannot have the asserted properties. Under E, F proves that the gap part is tame and contains every compatible disjoint tame part. O1 independently identifies it as the polar of all points of X for distinct comparability.

\paragraph{Correction and reconstruction (ours).} Use G+E for the actual-subset calculus. The qualified complement is still generally smaller than L\ensuremath{\setminus}X.

\paragraph{Independent second check.} The endpoint-gap proof and the all-points polar proof are independent descriptions of exactly the same subset under E.

\paragraph{Confidence.} High for mathematics; decisive wording checked in READINGS.tsv.

\noindent Refined statement identifier: H0KB.86.correction (the corrected statement and its proof in this section).

\section*{H0KB.87}\label{statement:H0KB.87}
\paragraph{Locator.} 156-3/031.14

\paragraph{His words.} Pour X=T\textsubscript{1}\ensuremath{\cup}\ensuremath{\cdots}\ensuremath{\cup}T\_n, X\ensuremath{\prime} est formé des n\ensuremath{-}1 tronçons interstitiels et des deux interstitiels extérieurs quand ils existent.

\paragraph{How he left it.} stated

\paragraph{Precise statement (ours).} The complement has precisely the consecutive internal gap pieces and the existing exterior pieces, provided the intrinsic endpoint calculus is valid.

\paragraph{Status.} False

\paragraph{Complete proof / counterexample application (ours).} The general G construction fails already for the shared complement of the twin singletons in D; its advertised endpoint data are ambiguous. Under E, the finite flag F lists exactly the unselected runs between selected component runs; H identifies each such run with its separator formula. No further component is possible.

\paragraph{Correction and reconstruction (ours).} G+E is the proved corrected version; n=0 is handled separately by \ensuremath{\varnothing}\ensuremath{\prime}=L.

\paragraph{Independent second check.} The direct separator inequalities of H reproduce the same list without using the Boolean complement mask.

\paragraph{Confidence.} High for mathematics; decisive wording checked in READINGS.tsv.

\noindent Refined statement identifier: H0KB.87.correction (the corrected statement and its proof in this section).

\section*{H0KB.88}\label{statement:H0KB.88}
\paragraph{Locator.} 156-3/032.1

\paragraph{His words.} Il donne n\ensuremath{\prime}\ensuremath{\in}\{n\ensuremath{-}1,n,n+1\}, puis le tableau selon T\textsubscript{1} minoré et T\_n majoré.

\paragraph{How he left it.} stated

\paragraph{Precise statement (ours).} In the corrected calculus, n\ensuremath{\prime}=n\ensuremath{-}1+e\_-+e\_+ for n>0, where e\_± record strict exterior bounds. The unqualified bounds in the printed table are insufficient at ambient extrema.

\paragraph{Status.} Slip

\paragraph{Complete proof / counterexample application (ours).} For X=L=[0,1], n=1 and n\ensuremath{\prime}=0, but the literal last line sees both a lower and upper bound and gives n\ensuremath{\prime}=2. Replace minoré/majoré by strictly minoré/majoré. H proves the corrected count by enumerating n\ensuremath{-}1 internal gaps and two optional exterior gaps. For X=\ensuremath{\varnothing}, n=0,n\ensuremath{\prime}=1.

\paragraph{Correction and reconstruction (ours).} Restore the strict qualifier inherited from scan 29. The broader missing-E construction issue is recorded in the preceding complement rows.

\paragraph{Independent second check.} The cropped first line reads non minoré (not non majoré); the numerical tests at \ensuremath{\varnothing}, L and a proper internal interval check all count possibilities.

\paragraph{Confidence.} High for mathematics; decisive wording checked in READINGS.tsv.

\noindent Refined statement identifier: H0KB.88.correction (the corrected statement and its proof in this section).

\section*{H0KB.89}\label{statement:H0KB.89}
\paragraph{Locator.} 156-3/032.14m

\paragraph{His words.} « Cette relation est-elle déductible de la relation d’ordre » sur M=Mod(L) ?

\paragraph{How he left it.} asked

\paragraph{Precise statement (ours).} On the originally claimed all-G subset calculus the relation need not be well defined. On the exactly repaired G+E calculus, inclusion alone determines compatibility.

\paragraph{Status.} False

\paragraph{Complete proof / counterexample application (ours).} D refutes the claimed premise of an intrinsic relation on all G. The strongest repaired answer is positive, with a complete proof O: atoms are singleton parts; two distinct atoms admit a join iff their points are comparable; formula O1 recovers tame complement as the polar of that graph; the three-piece Boolean decomposition then characterizes compatibility purely by inclusion, bounds and the recovered complement.

\paragraph{Correction and reconstruction (ours).} The positive reconstruction under E is ours. A rejected singleton-adjunction attempt is documented in WORK\_LOG and is not used in the accepted proof.

\paragraph{Independent second check.} The atomic-join obstruction is checked by strictly truncating every proposed upper bound of an incomparable pair. The polar formula is independently verified against the endpoint-gap complement in F.

\paragraph{Confidence.} High for mathematics; decisive wording checked in READINGS.tsv.

\noindent Refined statement identifier: H0KB.89.correction (the corrected statement and its proof in this section).

\section*{H0KB.90}\label{statement:H0KB.90}
\paragraph{Locator.} 156-3/032.14m

\paragraph{His words.} Il propose « p.ex. \{X,Y\} mod ssi X\ensuremath{\wedge}Y existe ??? ».

\paragraph{How he left it.} asked

\paragraph{Precise statement (ours).} Existence of a greatest lower bound in the inclusion poset characterizes compatibility.

\paragraph{Status.} False

\paragraph{Complete proof / counterexample application (ours).} Use the regular split line of O, which satisfies both G and E. Singleton parts on the two incomparable branches are incompatible, but their intersection \ensuremath{\varnothing} is their global greatest lower bound. O verifies all order and endpoint axioms. The correct order-only criterion is the atomic-polar/Boolean-decomposition criterion in H0KB.89.

\paragraph{Correction and reconstruction (ours).} The broader recovery question is not refuted by this failed proposed shortcut.

\paragraph{Independent second check.} No arbitrary intersection theorem is needed: the only subset of two distinct singletons is \ensuremath{\varnothing}, so it is unconditionally their meet.

\paragraph{Confidence.} High for mathematics; decisive wording checked in READINGS.tsv.

\noindent Refined statement identifier: H0KB.90.correction (the corrected statement and its proof in this section).

\section*{H0KB.91}\label{statement:H0KB.91}
\paragraph{Locator.} 156-3/033.1

\paragraph{His words.} Propriété 1 : M possède le plus petit élément \ensuremath{\varnothing} et le plus grand L, noté aussi 1.

\paragraph{How he left it.} stated

\paragraph{Precise statement (ours).} The inclusion poset of actual tame subsets has bottom \ensuremath{\varnothing} and top L.

\paragraph{Status.} True as written

\paragraph{Complete proof / counterexample application (ours).} The empty component family gives \ensuremath{\varnothing}, and the single full piece gives L. Every actual subset lies between them. This assertion is unaffected by endpoint ambiguity.

\paragraph{Correction and reconstruction (ours).} Our expansion of the displayed definitions and their hypotheses; no additional claim is attributed to the manuscript.

\paragraph{Independent second check.} The statement is checked directly in both D and the endpoint-regular examples.

\paragraph{Confidence.} High for mathematics; decisive wording checked in READINGS.tsv.

\section*{H0KB.92}\label{statement:H0KB.92}
\paragraph{Locator.} 156-3/033.3

\paragraph{His words.} Propriété 2 : X\ensuremath{\subset}Y implique \{X,Y\} modéré.

\paragraph{How he left it.} stated

\paragraph{Precise statement (ours).} Inclusion implies the intrinsic tame compatibility relation on all of M.

\paragraph{Status.} False

\paragraph{Complete proof / counterexample application (ours).} In D the ray V below the two twins is contained in the closed ray \ensuremath{\leq}a\_1, but its presentation with endpoint a\_0 is incompatible with that closed ray. The same subset has a compatible alternative presentation, demonstrating the intrinsicness failure. Under E the containment proof in F compares actual component infima/suprema and proves the property.

\paragraph{Correction and reconstruction (ours).} G+E is the corrected all-parts version; G alone suffices for closed parts.

\paragraph{Independent second check.} The closed-only version was proved separately; the counterexample uses precisely an omitted endpoint, so the two results do not conflict.

\paragraph{Confidence.} High for mathematics; decisive wording checked in READINGS.tsv.

\noindent Refined statement identifier: H0KB.92.correction (the corrected statement and its proof in this section).

\section*{H0KB.93}\label{statement:H0KB.93}
\paragraph{Locator.} 156-3/033.4

\paragraph{His words.} Propriété 3 : les éléments compatibles ont infimum et supremum; il ajoute la distributivité pour une famille mutuellement modérée.

\paragraph{How he left it.} stated

\paragraph{Precise statement (ours).} The displayed meet/join and finite distributive laws are asserted for the intrinsic compatible-subset calculus on every G.

\paragraph{Status.} False

\paragraph{Complete proof / counterexample application (ours).} The compatibility premise is not well defined on actual subsets in D, so the stated universal structure is not constructed. Under E, F embeds every finite compatible family in a common normalized-cell Boolean algebra; its meet is actual intersection, its join is normalized union and is globally least. Boolean distributivity follows there.

\paragraph{Correction and reconstruction (ours).} The repaired assertion is finite distributivity. No arbitrary infinite join or completeness claim is inserted; such joins need not exist even on R.

\paragraph{Independent second check.} The independent endpoint sweep verifies leastness beyond a chosen block, preventing an unwarranted inference from local to global joins.

\paragraph{Confidence.} High for mathematics; decisive wording checked in READINGS.tsv.

\noindent Refined statement identifier: H0KB.93.correction (the corrected statement and its proof in this section).

\section*{H0KB.94}\label{statement:H0KB.94}
\paragraph{Locator.} 156-3/033.7

\paragraph{His words.} Propriété 4 : un plus grand compatible disjoint X\ensuremath{\prime} existe; il est l’unique complément compatible, et X\ensuremath{\mapsto}X\ensuremath{\prime} est une anti-involution.

\paragraph{How he left it.} stated

\paragraph{Precise statement (ours).} The all-G actual-subset construction is asserted to have the stated complement, involution and order-reversal properties.

\paragraph{Status.} False

\paragraph{Complete proof / counterexample application (ours).} In D, distinct singleton twins have the same prescribed complement. No involution can take two different inputs to the same output. Under E, F proves maximality and uniqueness inside every common endpoint flag and global maximality by containment; bitwise complement is involutive, and refinement proves the order reversal globally. O independently identifies the complement as a graph polar.

\paragraph{Correction and reconstruction (ours).} G+E supplies the full corrected complement calculus, not merely an arbitrary choice of one presentation.

\paragraph{Independent second check.} The twin equality directly refutes injectivity of an involution. Under E, antitonicity also follows from formula O1 by restriction of the universal comparison condition.

\paragraph{Confidence.} High for mathematics; decisive wording checked in READINGS.tsv.

\noindent Refined statement identifier: H0KB.94.correction (the corrected statement and its proof in this section).

\section*{H0KB.95}\label{statement:H0KB.95}
\paragraph{Locator.} 156-3/033.8m

\paragraph{His words.} Le renforcement marginal exige que les éléments mutuellement modérés engendrent, par \ensuremath{\wedge},\ensuremath{\vee},\ensuremath{\prime}, une algèbre de Boole, finie si la famille est finie.

\paragraph{How he left it.} stated

\paragraph{Precise statement (ours).} The intrinsic concrete all-G structure is asserted to have compatible Boolean blocks. As an independent abstract postulate this is coherent and is explicitly assumed later.

\paragraph{Status.} False

\paragraph{Complete proof / counterexample application (ours).} D refutes the intrinsic complement needed for this universal concrete claim. Under E, F constructs common finite endpoint-cell blocks and their refinement embeddings, proving all Boolean identities and finiteness for finite generators. The abstract model M in lemma M separately shows the strengthened postulate is consistent.

\paragraph{Correction and reconstruction (ours).} The concrete failed assertion is distinguished from the valid abstract hypothesis on scan 34.

\paragraph{Independent second check.} For finite flags, distinct masks are distinguished by a point in a nonempty cell. This checks injectivity of the Boolean realization, not merely its symbolic identities.

\paragraph{Confidence.} High for mathematics; decisive wording checked in READINGS.tsv.

\noindent Refined statement identifier: H0KB.95.correction (the corrected statement and its proof in this section).

\section*{H0KB.96}\label{statement:H0KB.96}
\paragraph{Locator.} 156-3/033.13

\paragraph{His words.} La compatibilité est inchangée si l’on remplace l’un ou les deux termes par leur complément; il donne les deux lois de De Morgan.

\paragraph{How he left it.} stated

\paragraph{Precise statement (ours).} In the intended actual-subset structure, complement preserves compatibility in either argument and obeys the two De Morgan laws for compatible pairs.

\paragraph{Status.} False

\paragraph{Complete proof / counterexample application (ours).} On all G the operation is not intrinsic, by D. Under E, complementation keeps the same endpoint chain and Boolean NOT interchanges AND and OR in F, proving the laws. The same claims hold as part of the abstract assumptions in M.

\paragraph{Correction and reconstruction (ours).} The printed union-like glyph in one De Morgan formula means the tame join in the immediately stated algebraic context, not ordinary set union.

\paragraph{Independent second check.} The exhaustive Boolean truth-table check through eight cells independently verifies both identities; the endpoint proof establishes that those cells actually model the corrected calculus.

\paragraph{Confidence.} High for mathematics; decisive wording checked in READINGS.tsv.

\noindent Refined statement identifier: H0KB.96.correction (the corrected statement and its proof in this section).

\section*{H0KB.97}\label{statement:H0KB.97}
\paragraph{Locator.} 156-3/033.19

\paragraph{His words.} Propriété 5 : \ensuremath{\varnothing} et 1 appartiennent à M\_f, les parties fermées modérées.

\paragraph{How he left it.} stated

\paragraph{Precise statement (ours).} The empty and full parts are closed in the manuscript’s combinatorial sense.

\paragraph{Status.} True as written

\paragraph{Complete proof / counterexample application (ours).} They are respectively the empty union and the closed full piece. This proves the property without invoking any topology or E.

\paragraph{Correction and reconstruction (ours).} Our expansion of the displayed definitions and their hypotheses; no additional claim is attributed to the manuscript.

\paragraph{Independent second check.} Both are also topologically closed in every topology, but that independent fact is not needed for the stated combinatorial convention.

\paragraph{Confidence.} High for mathematics; decisive wording checked in READINGS.tsv.

\section*{H0KB.98}\label{statement:H0KB.98}
\paragraph{Locator.} 156-3/033.20

\paragraph{His words.} Propriété 6 : chaque X\ensuremath{\in}M a un plus petit Y\ensuremath{\in}M\_f avec X\ensuremath{\subset}Y; on pose Y=X̄.

\paragraph{How he left it.} stated

\paragraph{Precise statement (ours).} Every tame actual subset in G has a least closed tame hull.

\paragraph{Status.} False

\paragraph{Complete proof / counterexample application (ours).} D supplies the strict negative ray V, a single pseudo-piece, with no such hull. Lemma E proves the exact repair: within G, the existence of least closed tame hulls for all nonempty strict principal rays is equivalent to E. Under E, F proves all finite tame parts have least hulls. Thus E is necessary and sufficient for this universal hull assertion.

\paragraph{Correction and reconstruction (ours).} This is the central missing hypothesis of the later calculus; it is not a harmless endpoint convention.

\paragraph{Independent second check.} The independent canonical-topology characterization in E computes the same hull as the closure of a finite union of interval pieces.

\paragraph{Confidence.} High for mathematics; decisive wording checked in READINGS.tsv.

\noindent Refined statement identifier: H0KB.98.correction (the corrected statement and its proof in this section).

\section*{H0KB.99}\label{statement:H0KB.99}
\paragraph{Locator.} 156-3/033.20m

\paragraph{His words.} Par passage au complémentaire, il définit X\ensuremath{\circ}=(overline(X\ensuremath{\prime}))\ensuremath{\prime}.

\paragraph{How he left it.} stated

\paragraph{Precise statement (ours).} If the stipulated complement and closure axioms 1–7 hold, this formula defines the greatest open element below X, where open means complement of a closed element.

\paragraph{Status.} True as written

\paragraph{Complete proof / counterexample application (ours).} Order reversal gives (c(X\ensuremath{\prime}))\ensuremath{\prime}\ensuremath{\leq}X. It is open by definition. If U\ensuremath{\leq}X is open, then U\ensuremath{\prime} is closed and X\ensuremath{\prime}\ensuremath{\leq}U\ensuremath{\prime}, so c(X\ensuremath{\prime})\ensuremath{\leq}U\ensuremath{\prime} and U\ensuremath{\leq}(c(X\ensuremath{\prime}))\ensuremath{\prime}. This proves the greatest-open property using only the abstract axioms.

\paragraph{Correction and reconstruction (ours).} This is a valid conditional algebraic definition. The claim that every original G supplies its hypotheses is separately refuted in H0KB.94 and H0KB.98.

\paragraph{Independent second check.} Apply the same order-reversal argument twice to verify idempotence and monotonicity of interior.

\paragraph{Confidence.} High for mathematics; decisive wording checked in READINGS.tsv.

\section*{H0KB.100}\label{statement:H0KB.100}
\paragraph{Locator.} 156-3/033.22

\paragraph{His words.} Pour X,Y compatibles, l’adhérence vérifie overline(X\ensuremath{\vee}Y)=X̄\ensuremath{\vee}Ȳ.

\paragraph{How he left it.} stated

\paragraph{Precise statement (ours).} The all-G closure construction preserves finite compatible joins.

\paragraph{Status.} False

\paragraph{Complete proof / counterexample application (ours).} Closure is not defined for every tame X in D, so the asserted universal operation fails. Under E, F proves closure fills omitted endpoint cells and merges touching runs. Applying it before or after normalized union selects exactly the same occupied cells and forced endpoints. Therefore the displayed equality holds.

\paragraph{Correction and reconstruction (ours).} The exact domain correction is G+E; an infinite-join version is neither stated nor inferred.

\paragraph{Independent second check.} Independently, canonical topological closure under E commutes with finite unions of component intervals; normalizing a compatible join imposes the same outer closed endpoints.

\paragraph{Confidence.} High for mathematics; decisive wording checked in READINGS.tsv.

\noindent Refined statement identifier: H0KB.100.correction (the corrected statement and its proof in this section).

\section*{H0KB.101}\label{statement:H0KB.101}
\paragraph{Locator.} 156-3/033.24

\paragraph{His words.} Propriété 7 : si X,Y sont modérés, X̄,Y le sont, donc aussi X,Ȳ et X̄,Ȳ.

\paragraph{How he left it.} stated

\paragraph{Precise statement (ours).} The universal tame closure preserves compatibility.

\paragraph{Status.} False

\paragraph{Complete proof / counterexample application (ours).} The D ray has no closure in the stated sense, refuting the universal construction. Under E closure introduces no endpoint outside the old finite endpoint chain. Hence all the displayed pairs remain compatible, by F.

\paragraph{Correction and reconstruction (ours).} G+E proves the concrete property; M treats the abstract conditional separately.

\paragraph{Independent second check.} The relative abstract proof M uses this as an explicit hypothesis; it is not deduced from the incomplete G construction.

\paragraph{Confidence.} High for mathematics; decisive wording checked in READINGS.tsv.

\noindent Refined statement identifier: H0KB.101.correction (the corrected statement and its proof in this section).

\section*{H0KB.102}\label{statement:H0KB.102}
\paragraph{Locator.} 156-3/034.1

\paragraph{His words.} Il forme Fr X=X̄\ensuremath{\wedge}overline(X\ensuremath{\prime}), puis (Fr X)\ensuremath{\prime}=X\ensuremath{\circ}\ensuremath{\vee}(X\ensuremath{\prime})\ensuremath{\circ}.

\paragraph{How he left it.} stated

\paragraph{Precise statement (ours).} Under the stipulated abstract axioms 1–7, the frontier meet is defined and its complement has the displayed De Morgan expression.

\paragraph{Status.} True as written

\paragraph{Complete proof / counterexample application (ours).} X is compatible with X\ensuremath{\prime}. Applying closure compatibility twice makes c(X),c(X\ensuremath{\prime}) compatible, so their meet exists. De Morgan gives its complement as c(X)\ensuremath{\prime}\ensuremath{\vee}c(X\ensuremath{\prime})\ensuremath{\prime}, which is exactly (X\ensuremath{\prime})\ensuremath{\circ}\ensuremath{\vee}X\ensuremath{\circ}.

\paragraph{Correction and reconstruction (ours).} The conditional identity is true; finite geometric frontier requires the corrected concrete realization.

\paragraph{Independent second check.} The Boolean identity was independently checked on all finite masks through eight cells. This is an algebraic identity, not a claim that an arbitrary topological frontier is tame.

\paragraph{Confidence.} High for mathematics; decisive wording checked in READINGS.tsv.

\section*{H0KB.103}\label{statement:H0KB.103}
\paragraph{Locator.} 156-3/034.4

\paragraph{His words.} La propriété spéciale de dimension un relie la compatibilité de X,Y à celle de Fr X,Fr Y; la fin répète Fr X\ensuremath{\cup}Fr X.

\paragraph{How he left it.} stated

\paragraph{Precise statement (ours).} For the intrinsic interval calculus, X,Y are compatible iff their finite relative frontiers form a chain together. The last repeated X should be Y. The universal all-G realization is not well defined.

\paragraph{Status.} False

\paragraph{Complete proof / counterexample application (ours).} D gives ambiguous frontiers, so the subset-level assertion is not meaningful on all G. Under E, F says compatibility is exactly comparability of component endpoints; relative frontier deletes only ambient extrema, which are comparable to everything. This proves the equivalence. The ordinary union of the two finite frontier sets is then tame.

\paragraph{Correction and reconstruction (ours).} The substantial repair is E. The repeated Fr X in the final union is a subsidiary pen slip, documented in READINGS.tsv.

\paragraph{Independent second check.} Independently add back all ambient extrema: doing so changes neither the chain property nor compatibility. This checks the effect of the relative-boundary exception.

\paragraph{Confidence.} High for mathematics; decisive wording checked in READINGS.tsv.

\noindent Refined statement identifier: H0KB.103.correction (the corrected statement and its proof in this section).

\section*{H0KB.104}\label{statement:H0KB.104}
\paragraph{Locator.} 156-3/034.11

\paragraph{His words.} Dans le cas envisagé, X\ensuremath{\subset}X•, X=X•, et X fini sont équivalents; X• désigne la frontière, et X\ensuremath{\circ}=\ensuremath{\varnothing} signifie rare.

\paragraph{How he left it.} stated

\paragraph{Precise statement (ours).} For all tame X in the claimed concrete structure, empty tame interior, equality with its frontier and finiteness are equivalent.

\paragraph{Status.} False

\paragraph{Complete proof / counterexample application (ours).} On arbitrary G the required universal complement/closure structure fails by D. Under E, F proves the exact equivalence: every nondegenerate component contains a strict open interval by density, so empty tame interior forces all components to be singletons. Conversely a finite tame set is closed, has empty interior, and equals its finite frontier, including ambient singleton ends.

\paragraph{Correction and reconstruction (ours).} This is a statement about tame interior and finite endpoint sets, not about arbitrary nowhere-dense subsets of the canonical topology.

\paragraph{Independent second check.} The independent endpoint sweep verifies that a change of selected interval state leaves a frontier point, while a singleton has no interval cell in its interior.

\paragraph{Confidence.} High for mathematics; decisive wording checked in READINGS.tsv.

\noindent Refined statement identifier: H0KB.104.correction (the corrected statement and its proof in this section).

\section*{H0KB.105}\label{statement:H0KB.105}
\paragraph{Locator.} 156-3/034.16

\paragraph{His words.} Il prend désormais un ensemble ordonné M satisfaisant 1)–7), avec 4) renforcé, et écrit X+Y=(X\ensuremath{\vee}Y)\ensuremath{\wedge}(X\ensuremath{\wedge}Y)\ensuremath{\prime}, XY=X\ensuremath{\wedge}Y, X\ensuremath{\vee}Y=X+Y\ensuremath{-}XY.

\paragraph{How he left it.} stated

\paragraph{Precise statement (ours).} On every compatible Boolean block, the displayed operations give its Boolean ring; xy=x characterizes x\ensuremath{\leq}y. Incompatible pairs have no ring operation stipulated.

\paragraph{Status.} True as written

\paragraph{Complete proof / counterexample application (ours).} Lemma M identifies finite compatible subalgebras with subsets of their atoms. Addition is symmetric difference, multiplication intersection; the two-bit truth table proves all displayed formulas, with characteristic two so minus equals plus. Every identity involves finitely many elements and is checked in their Boolean block.

\paragraph{Correction and reconstruction (ours).} Our expansion of the displayed definitions and their hypotheses; no additional claim is attributed to the manuscript.

\paragraph{Independent second check.} The independent truth-table program verifies the identities through eight atoms. The formula for order is also the direct subset condition x\ensuremath{\cap}y=x.

\paragraph{Confidence.} High for mathematics; decisive wording checked in READINGS.tsv.

\section*{H0KB.106}\label{statement:H0KB.106}
\paragraph{Locator.} 156-3/035.1

\paragraph{His words.} Pour A\ensuremath{\in}M, il pose M\_A=\{Y:Y\ensuremath{\subset}A\}, avec l’ordre et la modération induits; les fermés sont X=X̄\ensuremath{\wedge}A, équivalemment X=Z\ensuremath{\wedge}A pour Z fermé compatible avec A.

\paragraph{How he left it.} stated

\paragraph{Precise statement (ours).} Under the explicit abstract axioms, this defines a coherent relative structure; the two closedness criteria are equivalent.

\paragraph{Status.} True as written

\paragraph{Complete proof / counterexample application (ours).} M proves equivalence. If X=A\ensuremath{\wedge}Z with Z closed, monotonicity of the least-closed-upper-bound operation gives c(X)\ensuremath{\leq}Z, hence X\ensuremath{\leq}A\ensuremath{\wedge}c(X)\ensuremath{\leq}A\ensuremath{\wedge}Z=X. Conversely use Z=c(X), compatible with A by axiom 7 and X\ensuremath{\leq}A.

\paragraph{Correction and reconstruction (ours).} Our expansion of the displayed definitions and their hypotheses; no additional claim is attributed to the manuscript.

\paragraph{Independent second check.} The second proof views A\ensuremath{\wedge}\ensuremath{-} on a common Boolean block as passage to its principal ideal. The equivalence uses only the leastness of closure, not an unproved meet-preservation property.

\paragraph{Confidence.} High for mathematics; decisive wording checked in READINGS.tsv.

\section*{H0KB.107}\label{statement:H0KB.107}
\paragraph{Locator.} 156-3/035.12

\paragraph{His words.} Dans 1), le plus petit élément relatif est \ensuremath{\varnothing} mais le plus grand est écrit X, bien que l’ambient relatif soit maintenant A.

\paragraph{How he left it.} stated

\paragraph{Precise statement (ours).} The top of M\_A is A. The surviving X is the variable from the earlier cancelled notation.

\paragraph{Status.} Slip

\paragraph{Complete proof / counterexample application (ours).} In the two-element Boolean algebra take A=1 and X=0. Then X is not the largest element below A. For arbitrary A the defining order interval has top A and bottom 0 by direct inclusion.

\paragraph{Correction and reconstruction (ours).} Replace X by A; the stray M\_Y later on the page likewise means M\_A in the explicitly repeated condition Y\ensuremath{\subset}A.

\paragraph{Independent second check.} The crop shows the surviving X while the surrounding indices have been changed to A, confirming a pen slip rather than a new convention.

\paragraph{Confidence.} High for mathematics; decisive wording checked in READINGS.tsv.

\noindent Refined statement identifier: H0KB.107.correction (the corrected statement and its proof in this section).

\section*{H0KB.108}\label{statement:H0KB.108}
\paragraph{Locator.} 156-3/035.13

\paragraph{His words.} Les infimum et supremum relatifs restent ceux de M; le complémentaire relatif de X est X\ensuremath{\prime}\ensuremath{\wedge}A, plus grand Y\ensuremath{\subset}A compatible avec X et disjoint de X.

\paragraph{How he left it.} proved on the page

\paragraph{Precise statement (ours).} For compatible X,Y\ensuremath{\leq}A, their global meet/join lie in M\_A and remain its bounds. The relative complement is A\ensuremath{\wedge}X\ensuremath{\prime}.

\paragraph{Status.} True as written

\paragraph{Complete proof / counterexample application (ours).} M proves the bound claim since A is an upper bound and the global join is least; the meet is below X. Any relative compatible disjoint Y is \ensuremath{\leq}X\ensuremath{\prime} and \ensuremath{\leq}A, hence \ensuremath{\leq}A\ensuremath{\wedge}X\ensuremath{\prime}. The latter itself is disjoint and compatible, in the Boolean block of X,A.

\paragraph{Correction and reconstruction (ours).} Our expansion of the displayed definitions and their hypotheses; no additional claim is attributed to the manuscript.

\paragraph{Independent second check.} The principal-ideal Boolean calculation independently verifies its complement equations and relative involutivity.

\paragraph{Confidence.} High for mathematics; decisive wording checked in READINGS.tsv.

\section*{H0KB.109}\label{statement:H0KB.109}
\paragraph{Locator.} 156-3/036.1

\paragraph{His words.} Les opérations relatives s’obtiennent en remultipliant par A; on trouve un idéal d’une algèbre de Boole, lui-même algèbre de Boole.

\paragraph{How he left it.} proved on the page

\paragraph{Precise statement (ours).} Every compatible family in M\_A generates a Boolean algebra with top A, obtained inside a common global Boolean block by intersecting with A. Relative complements preserve compatibility.

\paragraph{Status.} True as written

\paragraph{Complete proof / counterexample application (ours).} Adjoin A to the compatible family; comparable pairs ensure all its members are compatible with A. The strengthened block axiom gives a common Boolean algebra. Its principal ideal below A has relative complement A\ensuremath{\wedge}X\ensuremath{\prime} and top A, and the inherited meets and joins are global relative bounds.

\paragraph{Correction and reconstruction (ours).} The proof concerns a principal ideal, which has unit A; it does not claim every Boolean ideal has a unit. The isolated horizontal stroke after 1\_M is preserved in READINGS.tsv; the preceding closure operation is visibly cancelled. This row asserts the explicitly named Boolean operations. The closure operation is separately proved by the leastness argument in the following row.

\paragraph{Independent second check.} The map X\ensuremath{\mapsto}A\ensuremath{\wedge}X is a Boolean homomorphism to that principal ideal; checking complement explicitly proves the claimed operation correspondence.

\paragraph{Confidence.} High for mathematics; decisive wording checked in READINGS.tsv.

\section*{H0KB.110}\label{statement:H0KB.110}
\paragraph{Locator.} 156-3/036.18;156-3/037.1

\paragraph{His words.} Les fermés relatifs contenant X sont A\ensuremath{\wedge}Y avec Y fermé contenant X; le plus petit donne overline(X)\textasciicircum{}\{M\_A\}=A\ensuremath{\wedge}overline(X)\textasciicircum{}M.

\paragraph{How he left it.} proved on the page

\paragraph{Precise statement (ours).} The relative closure is c\_A(X)=A\ensuremath{\wedge}c(X), and relative 0,A are closed.

\paragraph{Status.} True as written

\paragraph{Complete proof / counterexample application (ours).} M proves that A\ensuremath{\wedge}c(X) is relatively closed using the equivalent description, contains X, and lies below every A\ensuremath{\wedge}Y with Y closed containing X. The two extrema satisfy the defining fixed-point condition.

\paragraph{Correction and reconstruction (ours).} Our expansion of the displayed definitions and their hypotheses; no additional claim is attributed to the manuscript.

\paragraph{Independent second check.} The same formula follows from the least-closed-upper-bound universal property without reproducing the page’s sequence of calculations.

\paragraph{Confidence.} High for mathematics; decisive wording checked in READINGS.tsv.

\section*{H0KB.111}\label{statement:H0KB.111}
\paragraph{Locator.} 156-3/037.3

\paragraph{His words.} Il calcule la préservation des supremums par l’adhérence relative, puis vérifie la compatibilité de overline(X)\textasciicircum{}A avec Y. « OK ».

\paragraph{How he left it.} proved on the page

\paragraph{Precise statement (ours).} Relative closure preserves compatible finite joins and compatibility; consequently the complete structure on M\_A satisfies all abstract axioms 1–7.

\paragraph{Status.} True as written

\paragraph{Complete proof / counterexample application (ours).} For compatible X,Y\ensuremath{\leq}A, axiom 7 makes A,c(X),c(Y) pairwise compatible. Distributivity in their common Boolean block gives A\ensuremath{\wedge}(c(X)\ensuremath{\vee}c(Y))=(A\ensuremath{\wedge}c(X))\ensuremath{\vee}(A\ensuremath{\wedge}c(Y)). The same block shows A\ensuremath{\wedge}c(X) compatible with Y. Combined with the previous relative rows, every axiom is verified.

\paragraph{Correction and reconstruction (ours).} Our expansion of the displayed definitions and their hypotheses; no additional claim is attributed to the manuscript.

\paragraph{Independent second check.} The principal-ideal homomorphism proof in M independently yields the join identity. The finite six-element model also checks nontrivial relativizations computationally.

\paragraph{Confidence.} High for mathematics; decisive wording checked in READINGS.tsv.

\section*{H0KB.112}\label{statement:H0KB.112}
\paragraph{Locator.} 156-3/037.12

\paragraph{His words.} M est « connexe » si les seuls X ouverts et fermés sont \ensuremath{\varnothing} et 1.

\paragraph{How he left it.} stated

\paragraph{Precise statement (ours).} Connectedness here is triviality of the clopen algebra of a structured partial Boolean poset; it is a definition and does not assert topological connectedness.

\paragraph{Status.} True as written

\paragraph{Complete proof / counterexample application (ours).} Open means complement of a closed element under the stipulated axioms. Thus the condition is well defined. The two extrema are always clopen; requiring no others gives the displayed definition, including the allowed trivial structure 0=1.

\paragraph{Correction and reconstruction (ours).} Our expansion of the displayed definitions and their hypotheses; no additional claim is attributed to the manuscript.

\paragraph{Independent second check.} The later interval example and the explicit “not for the canonical topology” note on scan 30 confirm this scope.

\paragraph{Confidence.} High for mathematics; decisive wording checked in READINGS.tsv.

\section*{H0KB.113}\label{statement:H0KB.113}
\paragraph{Locator.} 156-3/037.16

\paragraph{His words.} « On s’attend » qu’un X ouvert et fermé soit modéré avec tout Y.

\paragraph{How he left it.} hoped

\paragraph{Precise statement (ours).} Under abstract axioms 1–7 with the strengthened Boolean-block postulate, every clopen element is universally compatible.

\paragraph{Status.} False

\paragraph{Complete proof / counterexample application (ours).} The six-element horizontal sum in M satisfies every stated axiom with identity closure. All six elements are clopen, but a and b from different blocks are incompatible. The strongest direct correction is to add universal compatibility of clopens as a separate centrality hypothesis; the following Boolean and spectral conclusions then hold, as M proves.

\paragraph{Correction and reconstruction (ours).} This refutation is of the explicitly abstract hope, not a failure caused by the earlier G endpoint gap.

\paragraph{Independent second check.} check\_finite.py enumerates all axioms and every pairwise compatible subset of this model. The written block-by-block verification independently checks the same hypotheses.

\paragraph{Confidence.} High for mathematics; decisive wording checked in READINGS.tsv.

\noindent Refined statement identifier: H0KB.113.correction (the corrected statement and its proof in this section).

\section*{H0KB.114}\label{statement:H0KB.114}
\paragraph{Locator.} 156-3/037.18

\paragraph{His words.} « S’il en est ainsi », les ouverts-fermés engendrent une algèbre de Boole; les éléments de son spectre jouent le rôle de composantes connexes.

\paragraph{How he left it.} stated

\paragraph{Precise statement (ours).} If clopens are universally compatible, they form a Boolean algebra. Its ultrafilter/Boolean-ring spectrum records consistent clopen choices; a literal geometric component interpretation requires additional structure.

\paragraph{Status.} True as written

\paragraph{Complete proof / counterexample application (ours).} M proves closure under joins, complements and hence meets, and gives the ultrafilter topology with its separation and compactness proof, explicitly using the Boolean prime ideal theorem for arbitrary infinite algebras. The phrase “jouent le rôle” is retained as interpretation, not an unstated geometric representation theorem.

\paragraph{Correction and reconstruction (ours).} The conditional “S’il en est ainsi” is decisive and preserved. We do not mark the conditional false merely because its hoped-for premise fails in general.

\paragraph{Independent second check.} Finite cases require no choice: spectra are exactly the atoms, as the next row proves. This separates the algebraic conclusion from the interpretive phrase.

\paragraph{Confidence.} High for mathematics; decisive wording checked in READINGS.tsv.

\section*{H0KB.115}\label{statement:H0KB.115}
\paragraph{Locator.} 156-3/038.1

\paragraph{His words.} Si l’algèbre de Boole est finie et 0\ensuremath{\ne}1, les composantes sont les ouverts-fermés non nuls minimaux; si 0=1 la structure est vide.

\paragraph{How he left it.} stated

\paragraph{Precise statement (ours).} A finite nontrivial clopen Boolean algebra has atoms as its spectral points and algebraic components. The trivial Boolean algebra has empty spectrum.

\paragraph{Status.} True as written

\paragraph{Complete proof / counterexample application (ours).} Every nonzero element of a finite Boolean algebra contains a minimal nonzero atom. Distinct atoms are disjoint and their join is 1, or its nonzero complement would contain another atom. An ultrafilter contains exactly one atom and consists of its upper elements. The principal clopen algebra below an atom is \{0,a\}.

\paragraph{Correction and reconstruction (ours).} Our expansion of the displayed definitions and their hypotheses; no additional claim is attributed to the manuscript.

\paragraph{Independent second check.} A direct truth-table model as subsets of the finite atom set proves the same classification and the empty-spectrum exception.

\paragraph{Confidence.} High for mathematics; decisive wording checked in READINGS.tsv.

\section*{H0KB.116}\label{statement:H0KB.116}
\paragraph{Locator.} 156-3/038.8

\paragraph{His words.} Pour A=X\_\ensuremath{\Phi}=\ensuremath{\bigcup}T\_i dans Mod(L), les composantes connexes relatives sont les pseudo-tronçons composants T\_i.

\paragraph{How he left it.} stated

\paragraph{Precise statement (ours).} The concrete component identification is asserted for the all-G relative closure calculus. It is correct under G+E.

\paragraph{Status.} False

\paragraph{Complete proof / counterexample application (ours).} In D a component can fail to have a closed tame hull, so the required relative structured M\_A is not supplied by the original construction. Under E, H proves each T\_i clopen in A: its endpoint closure meets A exactly in T\_i. Each T\_i has no nontrivial tame clopen part. Hence the relative clopen algebra has precisely the T\_i as its finitely many atoms.

\paragraph{Correction and reconstruction (ours).} These are tame components. A rational interval can be topologically disconnected while remaining a single tame component.

\paragraph{Independent second check.} The finite-flag closure-mask proof independently shows a relative clopen mask must be constant throughout each T\_i.

\paragraph{Confidence.} High for mathematics; decisive wording checked in READINGS.tsv.

\noindent Refined statement identifier: H0KB.116.correction (the corrected statement and its proof in this section).

\section*{H0KB.117}\label{statement:H0KB.117}
\paragraph{Locator.} 156-3/039.1

\paragraph{His words.} Pour a), il suffit de overline(T\_i)\ensuremath{\wedge}A=T\_i; le NB dit que l’adhérence modérée d’un pseudo-tronçon est le fermé associé, puis s’arrête à a\ensuremath{\in}\ensuremath{\partial}T\ensuremath{\setminus}T…

\paragraph{How he left it.} sketched

\paragraph{Precise statement (ours).} Every pseudo-piece has its associated closed piece as its least tame closure, and separated components remain relatively closed.

\paragraph{Status.} False

\paragraph{Complete proof / counterexample application (ours).} D supplies a single open ray with two purported closed pieces and no least one. Under E, E proves the unique endpoint closure; the strict separator between different components excludes every newly filled endpoint from the other components, so closure(T\_i)\ensuremath{\cap}A=T\_i.

\paragraph{Correction and reconstruction (ours).} The omitted endpoint step is supplied as our proof under E; the dots do not constitute a proof on the page.

\paragraph{Independent second check.} The topological characterization E independently computes the closure. The finite-flag argument verifies relative closedness without assuming a completion at the ellipsis.

\paragraph{Confidence.} High for mathematics; decisive wording checked in READINGS.tsv.

\noindent Refined statement identifier: H0KB.117.correction (the corrected statement and its proof in this section).

\section*{H0KB.118}\label{statement:H0KB.118}
\paragraph{Locator.} 156-3/039.13

\paragraph{His words.} Pour b), il réduit à T=L; une frontière modérée vide force X=\ensuremath{\varnothing} ou X=L. Il conclut « OK ».

\paragraph{How he left it.} proved on the page

\paragraph{Precise statement (ours).} Under the concrete calculus’s hypotheses, a pseudo-tronçon is connected in the tame sense. The valid corrected realization is G+E.

\paragraph{Status.} False

\paragraph{Complete proof / counterexample application (ours).} The claimed all-G structure fails as recorded in H0KB.98, so its universal connectedness deduction lacks its defined closure/interior. Under E the full proof is H: a nonempty proper tame subset has an internal endpoint in its frontier; a singleton is immediate. Empty frontier therefore forces the empty or full part.

\paragraph{Correction and reconstruction (ours).} The argument is valid after the structure is repaired. It is not a claim that arbitrary open interval subsets are connected in \ensuremath{\tau}.

\paragraph{Independent second check.} The independent finite-cell proof in H checks that any nonconstant relative mask creates a frontier point.

\paragraph{Confidence.} High for mathematics; decisive wording checked in READINGS.tsv.

\noindent Refined statement identifier: H0KB.118.correction (the corrected statement and its proof in this section).

\section*{H0KB.119}\label{statement:H0KB.119}
\paragraph{Locator.} 156-3/040.1

\paragraph{His words.} Pour une famille finie compatible, S=\ensuremath{\bigcup}X\_\ensuremath{\alpha}•=\{t\textsubscript{1}<\ensuremath{\cdots}<t\_n\}; il affirme que l’algèbre de Boole engendrée a pour composants les points t\_i et les intervalles de la subdivision, et est connue par S.

\paragraph{How he left it.} stated

\paragraph{Precise statement (ours).} The Boolean algebra generated by the given family is asserted to be the full endpoint-cell algebra determined by the union of their frontiers.

\paragraph{Status.} False

\paragraph{Complete proof / counterexample application (ours).} Lemma B gives the counterexample X=[0,1] on R: one generator yields four elements and two atoms, whereas its five-cell subdivision gives 32 elements. The exact correction identifies generated atoms as joins of cells with equal generator-membership vectors; equality with the full cell algebra holds iff those vectors distinguish every cell. Two generators considered separately with the same singleton frontier also show S alone does not determine the generated algebra.

\paragraph{Correction and reconstruction (ours).} This is a substantive false claim even on a total real interval; E does not fix it.

\paragraph{Independent second check.} The truth-table argument in B and the explicit program count independently verify both the counterexample and the exact corrected classification.

\paragraph{Confidence.} High for mathematics; decisive wording checked in READINGS.tsv.

\noindent Refined statement identifier: H0KB.119.correction (the corrected statement and its proof in this section).

\section*{H0KB.120}\label{statement:H0KB.120}
\paragraph{Locator.} 156-3/040.13

\paragraph{His words.} Dans la liste après \{t\textsubscript{2}\}, l’intervalle est écrit ]t\textsubscript{1},t\textsubscript{3}[.

\paragraph{How he left it.} stated

\paragraph{Precise statement (ours).} The subdivision cell after \{t\textsubscript{2}\} is (t\textsubscript{2},t\textsubscript{3}), not (t\textsubscript{1},t\textsubscript{3}). Empty exterior cells are omitted.

\paragraph{Status.} Slip

\paragraph{Complete proof / counterexample application (ours).} For t\textsubscript{1}=0<t\textsubscript{2}=1<t\textsubscript{3}=2, (t\textsubscript{1},t\textsubscript{3}) contains t\textsubscript{2} and overlaps the preceding cell, so is not a subdivision cell. Replacing the lower index by 2 restores the disjoint ordered list C and F.

\paragraph{Correction and reconstruction (ours).} A local index slip, separate from H0KB.119.

\paragraph{Independent second check.} The crop confirms the surviving lower index 1. Direct membership of t\textsubscript{2} supplies the literal counterexample, independent of the generated-algebra error.

\paragraph{Confidence.} High for mathematics; decisive wording checked in READINGS.tsv.

\noindent Refined statement identifier: H0KB.120.correction (the corrected statement and its proof in this section).

\section*{H0KB.121}\label{statement:H0KB.121}
\paragraph{Locator.} 156-3/041.10

\paragraph{His words.} Il pense que le cas totalement ordonné et divisible a été traité depuis longtemps, « Hilbert ? », au moins par la relation d’entre-deux.

\paragraph{How he left it.} asked

\paragraph{Precise statement (ours).} The historical existence claim is correct: linear order can be reconstructed and axiomatized from betweenness in literature predating 1986. The tentative attribution to Hilbert is not certified here.

\paragraph{Status.} Known

\paragraph{Complete proof / counterexample application (ours).} Huntington and Kline (1917), sections 1 and 4, printed pp.302--303 and 322, give independent postulate systems for betweenness; section 5, p.324, constructs a linear order after choosing two base elements. The source PDF was opened and the construction read. The totally ordered dense case is a special case, with density added as the existence of a middle point for every distinct comparable pair.

\paragraph{Correction and reconstruction (ours).} The word Hilbert remains a question on the page; absence of a verified Hilbert locator is not a refutation of the correctly verified “done long ago” claim.

\paragraph{Independent second check.} Our direct R reconstruction independently proves the dense directed instance. The source’s actual all-points linear-order setting is matched; no claim about arbitrary partial orders is inferred from it.

\paragraph{Confidence.} High for mathematics; decisive wording checked in READINGS.tsv.

